% The Active Set: A Graduate Course in Convex Quadratic Optimization
%
% Thomas Schmelzer and contributors. Built from the companion research papers in
% this repository (matrix_free/, rmt/, non_negative_cg/, quadprog/) and in the
% sibling `homotopy` repository (statsci/, cla/); see the Preface for the map.
%
% Build with `make compile` (pdflatex -> bibtex -> pdflatex x2, plus an extra
% pass for the index). Chapters live in chapters/, front matter in frontmatter/.
\documentclass[11pt,openany]{book}

\usepackage{lmodern}
\usepackage[T1]{fontenc}
\usepackage{amsmath, amssymb, amsthm}
\usepackage{booktabs}
\usepackage{graphicx}
\usepackage{microtype}
\usepackage{algorithm}
\usepackage{algpseudocode}
\usepackage{makeidx}
\usepackage[margin=1.1in]{geometry}
\usepackage{fancyhdr}
\usepackage[colorlinks=true, linkcolor=blue, citecolor=blue, urlcolor=blue]{hyperref}

\makeindex

% ---------------------------------------------------------------------------
% Notation. The book fixes one notation across six papers that did not share
% one; Appendix A tabulates the translation into each field's dialect.
% ---------------------------------------------------------------------------
\newcommand{\R}{\mathbb{R}}
\newcommand{\T}{^{\top}}
\newcommand{\trans}{^{\top}}
\newcommand{\norm}[1]{\lVert #1 \rVert}
\newcommand{\inner}[2]{\langle #1, #2 \rangle}
\newcommand{\Active}{\mathcal{A}}
\newcommand{\Free}{\mathcal{F}}
\newcommand{\Bound}{\mathcal{B}}
\newcommand{\Sig}{\Sigma}
\newcommand{\onevec}{\mathbf{1}}
\newcommand{\Krylov}{\mathcal{K}}
\DeclareMathOperator{\range}{range}
\DeclareMathOperator{\nullsp}{null}
\DeclareMathOperator{\diag}{diag}
\DeclareMathOperator{\tr}{tr}
\DeclareMathOperator{\sign}{sign}
\DeclareMathOperator*{\argmin}{arg\,min}
\DeclareMathOperator*{\argmax}{arg\,max}

% ---------------------------------------------------------------------------
% Theorem environments, numbered by chapter.
% ---------------------------------------------------------------------------
\theoremstyle{plain}
\newtheorem{theorem}{Theorem}[chapter]
\newtheorem{proposition}[theorem]{Proposition}
\newtheorem{lemma}[theorem]{Lemma}
\newtheorem{corollary}[theorem]{Corollary}

\theoremstyle{definition}
\newtheorem{definition}[theorem]{Definition}
\newtheorem{example}[theorem]{Example}
\newtheorem{algorithmbox}[theorem]{Algorithm}

\theoremstyle{remark}
\newtheorem{remark}[theorem]{Remark}

% Exercises are numbered independently, per chapter, so that a problem set can
% be referred to as "Exercise 4.3" without competing with the theorem stream.
\newtheorem{exercise}{Exercise}[chapter]

\newenvironment{exercises}{%
  \section*{Exercises}\addcontentsline{toc}{section}{Exercises}%
}{}

\newenvironment{notes}{%
  \section*{Notes and references}\addcontentsline{toc}{section}{Notes and references}%
}{}

\pagestyle{fancy}
\fancyhf{}
\fancyhead[LE]{\small\leftmark}
\fancyhead[RO]{\small\rightmark}
\fancyfoot[C]{\thepage}
\renewcommand{\headrulewidth}{0.4pt}

\title{\Huge The Active Set\\[0.4em]
       \Large A Graduate Course in Convex Quadratic Optimization}
\author{Thomas Schmelzer}
\date{}

\begin{document}

\frontmatter
\maketitle
\tableofcontents
\chapter*{Preface}
\addcontentsline{toc}{chapter}{Preface}
\markboth{PREFACE}{PREFACE}

This is a course about one problem:
\[
  \min_{x \in \R^n} \; \tfrac12 x\T G x - a\T x
  \qquad\text{subject to}\qquad C\T x \ge b,
\]
a convex quadratic minimised over a polyhedron. That is a narrow subject to
give a book to, and the narrowness is the point. The problem is small enough
that everything about it can be proved, and large enough that the proving
teaches most of what a graduate student needs to know about constrained
optimization: optimality conditions and when they are sufficient, duality,
certificates, active sets and their combinatorics, saddle-point linear algebra,
Krylov subspace methods, conditioning, degeneracy, warm starts, and how to tell
whether a method is fast for a reason or by accident.

It is also a problem the reader has almost certainly met already, under some
other name. Mean--variance portfolio selection is exactly this problem. So is
non-negative least squares. So is every step of a sequential quadratic
programming method, every horizon of a linear model predictive controller, the
dual of a support vector machine, and---after one substitution that
Chapter~\ref{ch:homotopy} makes precise---the LASSO. Four literatures developed
solvers for it in parallel for half a century, largely without reference to one
another. One of this book's aims is to make that duplication visible, and then
useful.

\section*{The organising question}

Nearly all the difficulty in the problem above is combinatorial rather than
analytic. If someone tells you \emph{which} constraints hold with equality at
the solution---the \emph{active set}---the rest is a linear system. Everything
else is a strategy for finding that set. The book is organised around four such
strategies, and they turn out to exhaust the interesting possibilities:

\begin{enumerate}
\item \textbf{Walk to it}, one constraint at a time, maintaining optimality
      conditions on one side and driving violation on the other. This is the
      dual active-set method of Goldfarb and Idnani (Chapter~\ref{ch:walking}).
\item \textbf{Guess it}, wholesale, and repair the guess from the signs that
      come back. This is the primal--dual active set strategy, equivalently
      block principal pivoting on a linear complementarity problem
      (Chapter~\ref{ch:guessing}).
\item \textbf{Trace it}, as a parameter moves, recording the finitely many
      values at which it changes. This is parametric quadratic programming, and
      it is Markowitz's Critical Line Algorithm, least angle regression, the
      support-vector-machine path, and explicit model predictive control, all at
      once (Chapters~\ref{ch:parametric}--\ref{ch:reading}).
\item \textbf{Refuse to form the matrix}, and reach the problem only through the
      three linear actions an active-set method actually performs
      (Chapters~\ref{ch:krylov} and~\ref{ch:operator}).
\end{enumerate}

Two of these strategies come with finite-termination theorems and two do not.
The book's central methodological claim, argued from Chapter~\ref{ch:certificate}
onward, is that this matters less than it appears, because strict convexity makes
the Karush--Kuhn--Tucker conditions \emph{sufficient}. A candidate solution can
therefore be \emph{certified} rather than trusted. A method that may fail, but
that can recognise its own success, needs a certificate and not a convergence
theory. Much of modern computational practice runs on exactly that trade, and it
is worth teaching deliberately rather than leaving students to infer it.

\section*{Who this is for}

The reader is assumed to have had a first course in linear algebra including the
QR and Cholesky factorisations, a course in analysis sufficient for convexity and
compactness arguments, and enough numerical analysis to know that floating-point
arithmetic is not real arithmetic. No prior optimization course is assumed;
Chapter~\ref{ch:certificate} develops the Karush--Kuhn--Tucker conditions from
scratch for this problem class, which is the one setting where they can be proved
in a page and are sufficient as well as necessary. No finance, statistics, or
control background is assumed either: the applications are introduced where they
are used, and Appendix~\ref{app:notation} gives a translation table between the
four dialects.

Programming is not required to read the book, but the material is written by
people who implement, and it shows. Where a formula is a specification rather than
an implementation, the text says so. Where a constant in an algorithm exists to
defeat a floating-point hazard rather than a mathematical one, the text says that
too. Chapters~\ref{ch:numerics} and~\ref{ch:measurement} are given over entirely
to that layer, on the view that a graduate course which stops at the theorem
leaves students to rediscover the hard half in their first job.

\section*{How to use it}

The five parts are cumulative, but not equally so.

\begin{itemize}
\item \textbf{Part I} (Chapters~\ref{ch:oneproblem}--\ref{ch:freeblock}) is the
      foundation and should be read in order. It sets the problem, proves the
      optimality conditions and the sufficiency that everything later leans on,
      develops the combinatorial geometry of active sets, and assembles the
      linear algebra of the free block.
\item \textbf{Part II} (Chapters~\ref{ch:walking}--\ref{ch:krylov}) gives the
      three single-point strategies. The three chapters are largely independent
      of one another and can be taken in any order.
\item \textbf{Part III} (Chapters~\ref{ch:parametric}--\ref{ch:reading}) is the
      parametric theory and the cross-field synthesis. It depends on Part~I but
      only lightly on Part~II.
\item \textbf{Part IV} (Chapters~\ref{ch:operator}--\ref{ch:warmstart}) is about
      structure: reading the quadratic form as an operator, conditioning the
      problem, and exploiting families of related problems. It draws on both
      Part~II and Part~III.
\item \textbf{Part V} (Chapters~\ref{ch:numerics}--\ref{ch:measurement}) is the
      practice: numerics, degeneracy, and honest measurement.
\end{itemize}

A one-semester course can be built from Parts~I and~II with
Chapters~\ref{ch:parametric}, \ref{ch:homotopy}, and~\ref{ch:operator} added, which
is roughly ten weeks of material. A course emphasising statistics or machine
learning should keep Part~III entire and may compress Chapter~\ref{ch:walking}. A
course emphasising numerical linear algebra should keep
Chapters~\ref{ch:freeblock}, \ref{ch:krylov}, \ref{ch:operator},
and~\ref{ch:conditioning}, and may treat Part~III as reading.

Exercises follow every chapter. They are of three kinds, not marked as such but
easy to tell apart: routine verifications that a claim in the text was
understood; derivations that extend a result to a case the text skipped; and
open-ended computational problems, which say ``implement'' and mean it. The last
kind are the ones worth setting.

\section*{Sources}

The book is assembled from a set of companion research papers, and a reader who
wants the primary sources, the full experimental evidence, or the exact
statements as their authors made them should go to them.
Chapters~\ref{ch:walking} and~\ref{ch:guessing} follow \emph{Goldfarb--Idnani
Revisited: Invariants, Certificates, and the Limits of Guessing}.
Chapter~\ref{ch:krylov} follows \emph{Non-Negative Conjugate Gradients}
\cite{schmelzer2026nncg}. Chapters~\ref{ch:operator} and~\ref{ch:conditioning}
draw on \emph{Matrix-Free Methods for Long-Only Portfolio Optimization}
\cite{schmelzer2026matrixfree} and \emph{From Marchenko--Pastur to Woodbury}
\cite{schmelzer2026rmt}. Part~III follows \emph{One Homotopy Across Statistics,
Optimization, and Finance} \cite{schmelzer2026perspective} and the \texttt{cvxcla}
software paper \cite{cvxcla}. Where the book states an empirical number, it is
one of theirs, and the citation says which.

The book's own contribution is the arrangement: one notation, one set of proofs,
and the argument that the four strategies of the organising question are answers
to a single question rather than four separate subjects. Any error in the
translation is the book's and not the papers'.

\bigskip
\noindent\emph{Abu Dhabi}


\mainmatter

\part{Foundations}
\chapter{One Problem, Four Answers}
\label{ch:oneproblem}

\section{The problem}

Fix a symmetric positive definite matrix $G \in \R^{n\times n}$, a vector
$a \in \R^n$, a matrix $C \in \R^{n \times m}$ whose columns $c_1,\dots,c_m$ are
constraint normals, and a right-hand side $b \in \R^m$. The first
$m_{\mathrm{eq}}$ constraints are equalities and the rest inequalities. The
problem is
\begin{equation}
  \min_{x \in \R^n} \; f(x) := \tfrac12 x\T G x - a\T x
  \qquad \text{subject to} \qquad
  \begin{cases}
    c_i\T x = b_i, & i \le m_{\mathrm{eq}},\\
    c_i\T x \ge b_i, & i > m_{\mathrm{eq}},
  \end{cases}
  \label{eq:qp}
\end{equation}
which we will call \emph{the} quadratic program, or QP, throughout.
\index{quadratic program!statement}

Two conventions deserve a word, because they differ from some textbooks and are
inherited from the solver this book's Part~II describes. The linear term is
\emph{subtracted}, so that the unconstrained minimiser is $G^{-1}a$ rather than
$-G^{-1}a$; and the constraints are held \emph{column-wise} in $C$, so the
constraint system reads $C\T x \ge b$ rather than $Cx \le b$. Neither choice
affects anything beyond signs, and both make the formulas in Part~II shorter.

Everything in this book is downstream of two hypotheses in~\eqref{eq:qp}. The
first is that $G \succ 0$, which makes $f$ strictly convex and coercive, so the
problem has exactly one solution whenever it has any. The second is that the
constraints are \emph{linear}, so the feasible set is a polyhedron: an
intersection of finitely many half-spaces and hyperplanes, hence a set with
finitely many faces. Together these give the problem a character that is worth
stating before any algorithm appears.

\begin{quote}
\emph{The analysis is easy and the combinatorics is hard.}
\end{quote}

If an oracle told you which of the $m$ constraints hold with equality at the
solution, you could finish in one linear solve; Chapter~\ref{ch:certificate}
proves this. There are $2^{m-m_{\mathrm{eq}}}$ candidate answers the oracle might
give, and the whole of algorithmic practice is a strategy for searching among
them without enumerating them.

\section{Where it comes from}

The QP~\eqref{eq:qp} is not a mathematical curiosity to be motivated after the
fact. It is solved a great many times a day, in at least four traditions that
have historically not read one another's papers.

\paragraph{Portfolio selection.} An investor spreads weights $w \in \R^n$ across
$n$ assets with covariance $\Sig$ and expected returns $\mu$, subject to being
fully invested, $\onevec\T w = 1$, and---under a long-only mandate---holding no
short positions, $w \ge 0$. Minimising variance against a return tilt $\rho$,
\[
  \min_w \; w\T \Sig w - \rho\,\mu\T w
  \qquad\text{subject to}\qquad \onevec\T w = 1, \quad w \ge 0,
\]
is~\eqref{eq:qp} with $G = 2\Sig$, $a = \rho\mu$, and $C = [\onevec, I]$. This is
Markowitz's problem \cite{markowitz1952}, and it is the running application of
Parts~III and~IV.
\index{portfolio selection}

\paragraph{Non-negative least squares.} Given $M \in \R^{m\times n}$ and
$d \in \R^m$, minimise $\norm{Mx - d}_2^2$ over $x \ge 0$. Expanding the square
gives~\eqref{eq:qp} with $G = 2M\T M$, $a = 2M\T d$, $C = I$, $b = 0$. This is
the problem of Lawson and Hanson \cite{lawson1974}, and it is the setting in
which Chapter~\ref{ch:krylov} works, because $G$ arrives as a factored
Gram matrix and the whole method can be run without ever assembling it.
\index{non-negative least squares}

\paragraph{Model predictive control.} At each sampling instant a controller
minimises a quadratic cost over a finite horizon of inputs subject to linear
state and input constraints. The current state enters as a \emph{parameter} of
the QP. Solving the family once, as a function of that parameter, rather than
solving one instance per instant, is explicit model predictive control
\cite{bemporad2002}, and it is the multi-parameter cousin of the theory in
Part~III.
\index{model predictive control}

\paragraph{Regularised regression.} The LASSO
\cite{tibshirani1996}, $\min_\beta \tfrac12\norm{y - X\beta}_2^2 +
\lambda\norm{\beta}_1$, does not look like~\eqref{eq:qp}: its penalty is not
smooth. But its subgradient condition is a sign pattern on a set of coordinates,
which is an active set by another name, and under the substitution
$G = X\T X$, $a = X\T y$ its solution path is the \emph{same} piecewise-linear
curve as a portfolio frontier. Theorem~\ref{thm:identity} makes this exact. It is
the sharpest illustration in the book of the claim that these are one subject.
\index{LASSO}

The four traditions arrived at the same algorithms independently, four decades
apart in the two extreme cases. Appendix~\ref{app:history} tells that story;
Appendix~\ref{app:notation} translates the vocabulary.

\section{The active set}

\begin{definition}[Active set, working set]
\label{def:activeset}
\index{active set!definition}
For a feasible $x$, the \emph{active set} at $x$ is
$\Active(x) = \{i : c_i\T x = b_i\}$, the constraints holding with equality
there. A \emph{working set} is any index set $\Active \subseteq \{1,\dots,m\}$
that an algorithm is currently \emph{holding} as equalities, whether or not it
is the active set at the current iterate or at the solution.
\end{definition}

The distinction between the two matters and is a frequent source of confusion.
The active set is a property of a point; the working set is a piece of algorithm
state. At a solution the working set of a correct algorithm is a subset of the
active set, and the two coincide except in the degenerate case of
Chapter~\ref{ch:geometry}.

Given a working set $\Active$ of size $k$ with normals $A = C_{\cdot\Active}$, the
\emph{working-set subproblem} drops every constraint outside $\Active$ and
promotes every constraint inside it to an equality:
\begin{equation}
  \min_{x} \; \tfrac12 x\T G x - a\T x
  \qquad \text{subject to} \qquad A\T x = b_\Active .
  \label{eq:sub}
\end{equation}
This is an equality-constrained problem, and equality-constrained problems are
linear algebra. Chapter~\ref{ch:certificate} solves~\eqref{eq:sub} in closed
form, Chapter~\ref{ch:freeblock} shows how to solve it fast and how to update the
solution when $\Active$ changes by one index, and the rest of the book is about
choosing $\Active$.

\section{Four answers}

Here are the four strategies, stated in one paragraph each so that the reader
knows where the book is going. Each gets a chapter, or three.

\paragraph{Walk to it.} Start from a working set you can solve trivially, and
move to the solution one index at a time, maintaining part of the optimality
conditions exactly and driving the violated part to zero. The classical primal
method maintains feasibility and works toward stationarity, which requires a
feasible starting point and therefore a phase-one problem. The \emph{dual}
method of Goldfarb and Idnani \cite{goldfarb1982,goldfarb1983} reverses the
roles: it maintains stationarity and the multiplier signs, and works toward
feasibility. Its starting point is free---with no constraint active, $x = G^{-1}a$
satisfies everything the method asks---so one Cholesky factorisation replaces the
phase-one problem entirely. Chapter~\ref{ch:walking} derives it, and shows that
the whole method follows from two matrix identities.

\paragraph{Guess it.} Solve the subproblem on a guessed working set, read the
signs of the resulting multipliers and slacks, and let those signs tell you which
indices to add and drop---\emph{all of them at once}. This is the primal--dual
active set strategy \cite{hintermuller2002}, equivalently block principal
pivoting on a linear complementarity problem \cite{judice1994,murty1988}. It
converges in a handful of repairs almost independently of $n$, trading arithmetic
(abundant) for iterations (not). Chapter~\ref{ch:guessing} develops it, proves
finite termination in the bound-constrained case $C = I$, and shows exactly why
that theorem does not survive the general constraint class.

\paragraph{Trace it.} Attach a scalar parameter to the linear term,
$a \mapsto \lambda d$, and ask how the solution moves as $\lambda$ varies. The
answer is that it is \emph{piecewise linear}: on any interval where the active set
is constant the optimiser is affine in $\lambda$, and the active set changes at
finitely many breakpoints. So the whole family can be traced exactly, corner to
corner, by a single finite algorithm. Part~III develops this, and shows that
Markowitz's Critical Line Algorithm \cite{markowitz1956}, least angle regression
\cite{efron2004}, the support-vector-machine path \cite{hastie2004svm}, and
explicit model predictive control \cite{bemporad2002} are one algorithm with four
fillings.

\paragraph{Refuse to form the matrix.} Notice how little of $G$ any of the above
actually touches. An active-set method needs a product $v \mapsto Gv$, a solve
against the free block, and a cross product between free and blocked
coordinates---three linear \emph{actions}, not $n^2$ stored numbers. Anything
that supplies those actions can play the role of $G$: a factored Gram matrix
$M\T M$ reached through two products with $M$, a diagonal-plus-low-rank factor
model inverted by the Woodbury identity, a kernel given implicitly by its kernel
function. Chapters~\ref{ch:krylov} and~\ref{ch:operator} develop this reading,
which is where most of the practical speed in this book comes from.

\section{Certificates, and why they are the point}

Two of the four strategies come with finite-termination theorems and two do not.
Guessing can cycle. A parametric trace can be derailed by degeneracy. It would be
natural to conclude that the guaranteed methods are the serious ones and the rest
are heuristics.

That conclusion is wrong, and understanding why is the main methodological lesson
of this book. Because $G \succ 0$, the Karush--Kuhn--Tucker conditions
for~\eqref{eq:qp} are not merely necessary but \emph{sufficient}
(Proposition~\ref{prop:sufficient}). A point equipped with multipliers passing
four inequalities \emph{is} the answer, whatever produced it. So a method that
produces candidates may be as unprincipled as one likes, provided each candidate
is tested; and the test is cheap, being one matrix--vector product and a pass over
the constraints.

This turns the design question around. Instead of asking ``does this method
converge?'' one asks ``can this method recognise its own success, and what does it
do when it fails?'' A method with a certificate and a fallback is safe without
being guaranteed. A guaranteed method whose guarantee rests on an assumption you
cannot check---non-degeneracy, say---may be less safe in practice than an
unguaranteed one that verifies its answer. Chapter~\ref{ch:certificate} proves the
sufficiency; Chapters~\ref{ch:guessing}, \ref{ch:krylov}, and~\ref{ch:warmstart}
each cash it in.

\section{What this book does not cover}

Honesty about scope. The book treats convex quadratic objectives over polyhedra,
and stays there.

\emph{Interior-point methods} are barely mentioned, and only for comparison. They
are the right tool for large sparse problems and they are excellently covered
elsewhere \cite{nocedal2006,boyd2004}. Their relationship to this book's methods is
worth stating once: an interior-point method approaches a face of the polyhedron
asymptotically, whereas an active-set method arrives at one exactly and in finitely
many steps. That is why active-set methods warm-start almost perfectly and
interior-point methods do not, and warm starts are what Part~III lives on.

\emph{Non-convex quadratic programs} ($G$ indefinite) are a different subject:
NP-hard, with local minima, and none of the sufficiency arguments here apply.
\emph{Semidefinite $G$} is a middle case and is genuinely within reach---Boland
\cite{boland1997} extends the dual method to it with finite termination, and the
solver of \cite{arnstrom2022daqp} handles it by proximal regularisation---but the
factorisation of Chapter~\ref{ch:walking} needs $G^{-1}$, so the book keeps
$G \succ 0$ as a matter of scope. Where that restriction bites, the text says so.

\emph{General nonlinear constraints} are out of scope, though the QP is not: a
sequential quadratic programming method solves one instance of~\eqref{eq:qp} per
outer iteration, so everything here is a subroutine of that larger subject.

\begin{notes}
The QP~\eqref{eq:qp} in this form, with its two conventions, is the signature of
the \texttt{cvx-quadprog} package \cite{cvxquadprog}, whose mathematics
Chapters~\ref{ch:walking} and~\ref{ch:guessing} follow. Nocedal and Wright
\cite[Ch.~16]{nocedal2006} is the standard reference for quadratic programming
generally and should be on the desk alongside this book; Boyd and Vandenberghe
\cite{boyd2004} for the convex-analytic background. Golub and Van Loan
\cite{golub2013} is the linear-algebra reference used throughout.

The claim that four fields developed one algorithm independently is argued in
full in \cite{schmelzer2026perspective} and retold in Appendix~\ref{app:history}.
\end{notes}

\begin{exercises}

\begin{exercise}
Write the non-negative least squares problem $\min_{x\ge0}\norm{Mx-d}_2^2$ in the
form~\eqref{eq:qp}, giving $G$, $a$, $C$, $b$, and $m_{\mathrm{eq}}$ explicitly.
For which $M$ is the resulting $G$ positive definite rather than merely
semidefinite?
\end{exercise}

\begin{exercise}
Show that the long-only fully invested portfolio problem has a non-empty feasible
set for every $n \ge 1$, and that the feasible set is compact. Deduce that the
problem has a solution without appealing to coercivity of $f$.
\end{exercise}

\begin{exercise}
The feasible set of~\eqref{eq:qp} with $m_{\mathrm{eq}} = 0$ and $C = I$, $b = 0$
is the non-negative orthant. How many faces does it have? How many candidate
active sets must an enumerating algorithm consider for $n = 50$? Compare with the
number of atoms in the observable universe ($\approx 10^{80}$).
\end{exercise}

\begin{exercise}
Let $x_u = G^{-1}a$ be the unconstrained minimiser. Show that if $x_u$ is
feasible then it solves~\eqref{eq:qp}, and that the active set at the solution is
then contained in $\{i : c_i\T x_u = b_i\}$. Why does this observation make a good
starting heuristic for the guessing strategy of Chapter~\ref{ch:guessing}?
\end{exercise}

\begin{exercise}
Consider~\eqref{eq:qp} with $n = 2$, $G = I$, $a = (0,0)\T$, and the single
constraint $x_1 + x_2 \ge 1$. Solve it by hand, identify the active set at the
solution, and verify that the solution lies on the constraint boundary. Now change
$a$ to $(2,2)\T$ and repeat. What happened to the active set, and why?
\end{exercise}

\begin{exercise}[Implementation]
Implement a brute-force QP solver: enumerate every subset of the inequality
constraints, solve the equality-constrained subproblem~\eqref{eq:sub} on each,
discard the infeasible ones, and return the feasible point of least objective
value. Test it on random instances with $n = 5$ and $m = 8$. Time it as $m$ grows
to $16$. This solver is correct and useless, and the rest of the book explains
why both halves of that sentence are true.
\end{exercise}

\end{exercises}

\chapter{The Convex Quadratic Program and Its Certificate}
\label{ch:certificate}

This chapter proves the four facts everything later rests on: that~\eqref{eq:qp}
has exactly one solution; that the Karush--Kuhn--Tucker conditions characterise
it; that they are \emph{sufficient} and not merely necessary; and that the
working-set subproblem has a closed-form solution. The last section takes the
dual, which explains where the method of Chapter~\ref{ch:walking} gets its name
and why its objective increases.

The mathematics is short. The consequence---that a candidate answer can be
checked rather than trusted---is the organising idea of the book, and
Section~\ref{sec:certificate-principle} states it as a design principle rather
than a theorem.

\section{Existence and uniqueness}

Let
\[
  \Omega = \{x \in \R^n : c_i\T x = b_i \ (i \le m_{\mathrm{eq}}),\ \
                          c_i\T x \ge b_i \ (i > m_{\mathrm{eq}})\}
\]
denote the feasible set, a closed convex polyhedron, and recall
$f(x) = \tfrac12 x\T G x - a\T x$ with $G \succ 0$.

\begin{proposition}[Existence and uniqueness]
\label{prop:existunique}
If $\Omega \ne \emptyset$ then~\eqref{eq:qp} has exactly one solution.
\end{proposition}

\begin{proof}
Since $G \succ 0$, writing $\lambda_{\min} > 0$ for its smallest eigenvalue gives
$f(x) \ge \tfrac12\lambda_{\min}\norm{x}^2 - \norm{a}\norm{x} \to \infty$ as
$\norm{x}\to\infty$, so $f$ is coercive. Fix any $x_0 \in \Omega$; the sublevel set
$\{x \in \Omega : f(x) \le f(x_0)\}$ is then closed and bounded, hence compact and
non-empty, and a continuous function attains its minimum on it. That minimum is a
minimum over all of $\Omega$. Uniqueness: if $x \ne y$ both minimise, then by
strict convexity $f(\tfrac12(x+y)) < \tfrac12 f(x) + \tfrac12 f(y) = \min_\Omega f$,
and $\tfrac12(x+y) \in \Omega$ by convexity of $\Omega$, a contradiction.
\end{proof}

The proposition is unremarkable, but note which hypothesis did which work.
Coercivity came from $G \succ 0$ and gave existence; strict convexity came from
the same place and gave uniqueness. Drop to $G \succeq 0$ and both arguments fail,
though both conclusions may survive---a semidefinite QP on a compact polyhedron
still has a solution, and may have a whole face of them.

\section{Optimality conditions}

The Lagrangian of~\eqref{eq:qp} is
$L(x,\lambda) = \tfrac12 x\T G x - a\T x - \lambda\T(C\T x - b)$, and the
Karush--Kuhn--Tucker (KKT) conditions are
\index{Karush--Kuhn--Tucker conditions}
\begin{align}
  G x - a &= C \lambda, \tag{stationarity}\label{eq:stat}\\
  c_i\T x &= b_i \ (i \le m_{\mathrm{eq}}), \qquad
  c_i\T x \ge b_i \ (i > m_{\mathrm{eq}}), \tag{primal feasibility}\label{eq:pfeas}\\
  \lambda_i &\ge 0 \quad (i > m_{\mathrm{eq}}), \tag{dual feasibility}\label{eq:dfeas}\\
  \lambda_i\,(c_i\T x - b_i) &= 0 \quad (i > m_{\mathrm{eq}}).
      \tag{complementarity}\label{eq:comp}
\end{align}

Multipliers of equality constraints are unrestricted in sign. That single
asymmetry propagates into every formula in Part~II, and forgetting it is the most
common implementation error in this subject.

The four conditions have a reading worth internalising before any algebra.
Stationarity says the objective's gradient at $x$ is a combination of the
constraint normals: there is no direction of descent that the constraints permit.
Dual feasibility says the combination uses the inequality normals with the right
sign, pushing \emph{into} the feasible set rather than out of it. Complementarity
says only constraints that actually touch $x$ get to contribute. Primal
feasibility says $x$ is in $\Omega$ at all.

That the conditions are \emph{necessary} at a solution follows from any standard
constraint qualification, and for linear constraints no qualification is needed:
linearity is itself enough (see \cite[Ch.~12]{nocedal2006}). What matters here is
the converse.

\section{Sufficiency: the certificate}
\label{sec:certificate-principle}

\begin{proposition}[Sufficiency]
\label{prop:sufficient}
\index{sufficiency of KKT}
If $(x,\lambda)$ satisfies \eqref{eq:stat}--\eqref{eq:comp} then $x$ is the unique
solution of~\eqref{eq:qp}.
\end{proposition}

\begin{proof}
Let $\tilde x \in \Omega$. By strict convexity,
$f(\tilde x) \ge f(x) + (G x - a)\T(\tilde x - x)$, with equality only if
$\tilde x = x$. Substituting~\eqref{eq:stat},
\[
  (G x - a)\T(\tilde x - x)
  = \lambda\T\bigl(C\T \tilde x - C\T x\bigr)
  = \sum_i \lambda_i\bigl[(c_i\T\tilde x - b_i) - (c_i\T x - b_i)\bigr]
  = \sum_i \lambda_i (c_i\T \tilde x - b_i),
\]
the last step by~\eqref{eq:comp} for the inequalities and by $c_i\T x = b_i$ for
the equalities. Each surviving term is non-negative: for
$i \le m_{\mathrm{eq}}$ the bracket vanishes, and for $i > m_{\mathrm{eq}}$ both
factors are non-negative by~\eqref{eq:pfeas} and~\eqref{eq:dfeas}. Hence
$f(\tilde x) \ge f(x)$ with equality only at $\tilde x = x$.
\end{proof}

The proof is four lines and it licenses most of the interesting algorithms in this
book. Stated as a design principle:

\begin{quote}
\textbf{The certificate principle.} A method that produces candidate pairs
$(x,\lambda)$ may be as unprincipled as one likes---a guess, a cached answer from
a neighbouring problem, the output of a heuristic---provided every candidate is
tested against \eqref{eq:stat}--\eqref{eq:comp}. Sufficiency is what makes an
unguaranteed method \emph{safe} rather than \emph{usually right}.
\end{quote}

In floating point the four conditions become four tolerances. A candidate is
accepted when
\begin{equation}
  \norm{G x - a - C\lambda}_\infty \le \varepsilon,
  \quad |c_i\T x - b_i| \le \varepsilon \ (i \le m_{\mathrm{eq}}),
  \quad c_i\T x - b_i \ge -\varepsilon, \ \ \lambda_i \ge -\varepsilon,
  \ \ |\lambda_i (c_i\T x - b_i)| \le \varepsilon
  \label{eq:certificate}
\end{equation}
on the inequalities, with $\varepsilon$ scaled by the data. Two points about this
test recur throughout the book and are worth flagging now.

First, \emph{stationarity is checked against $G$ directly}, not inherited from
whatever construction produced $x$. The construction is exactly what an
ill-conditioned working set corrupts, so trusting it defeats the purpose of
checking.

Second, the test is not delicate, because the quantity it thresholds is
\emph{bimodal}. A candidate built from a well-conditioned working set satisfies
\eqref{eq:certificate} to a few multiples of the machine epsilon; one built on a
set that is not has no reason to come anywhere close. The two populations lie
orders of magnitude apart and the threshold sits in the empty space between them.
Chapter~\ref{ch:numerics} returns to this pattern, which is the single most useful
thing to know about tolerances in this subject.

\section{The working-set subproblem}

Fix a working set $\Active$ of size $k$ with normals $A = C_{\cdot\Active}$ and
suppose $A$ has full column rank $k$. Recall the subproblem~\eqref{eq:sub}: minimise
$f$ subject to $A\T x = b_\Active$, with every constraint in $\Active$ held as an
equality regardless of its type in the original program and every constraint
outside $\Active$ ignored.

\begin{proposition}[Solution of the subproblem]
\label{prop:sub}
\index{working set!subproblem}
Problem~\eqref{eq:sub} has the unique solution and multipliers
\begin{equation}
  u = \bigl(A\T G^{-1} A\bigr)^{-1}\bigl(b_\Active - A\T x_u\bigr),
  \qquad
  x = x_u + G^{-1} A\, u,
  \qquad x_u := G^{-1} a .
  \label{eq:subsol}
\end{equation}
\end{proposition}

\begin{proof}
Stationarity for~\eqref{eq:sub} reads $Gx - a = Au$, so
$x = G^{-1}a + G^{-1}Au = x_u + G^{-1}Au$. Substituting into $A\T x = b_\Active$
gives $A\T x_u + (A\T G^{-1}A)u = b_\Active$. The matrix $A\T G^{-1}A$ is positive
definite because $G^{-1} \succ 0$ and $A$ has full column rank, hence invertible,
which yields~\eqref{eq:subsol}; the point so constructed is the unique minimiser by
strict convexity of $f$ on the affine feasible set of~\eqref{eq:sub}.
\end{proof}

\begin{definition}[Dual Hessian]
\label{def:dualhessian}
\index{dual Hessian}
The matrix $A\T G^{-1} A$ is the \emph{dual Hessian} of the working set $\Active$.
\end{definition}

The dual Hessian appears under three names in what follows: as the matrix whose
Cholesky factor the solver of Chapter~\ref{ch:walking} carries; as the coefficient
matrix of the linear system that the guessing method of Chapter~\ref{ch:guessing}
solves; and as the Gram matrix of a least-squares problem in
Lemma~\ref{lem:directions}. Its invertibility is equivalent to $A$ having full
column rank, and \emph{every rank condition in this book comes from that single
fact}.

The pair $(x,u)$ of~\eqref{eq:subsol}, extended by zeros off the working set,
satisfies stationarity~\eqref{eq:stat} and complementarity~\eqref{eq:comp} by
construction: the working-set constraints have zero slack and the rest have zero
multiplier. What it need \emph{not} satisfy is the sign condition~\eqref{eq:dfeas}
on the inequalities in $\Active$, or primal feasibility~\eqref{eq:pfeas} off
$\Active$. Two of the four conditions come free; the other two are what an
algorithm must work for. Which two an algorithm keeps and which it chases is
exactly what distinguishes a primal from a dual method.

\begin{remark}[Cold start]
\label{rem:cold}
For $\Active = \emptyset$ the subproblem is unconstrained and~\eqref{eq:subsol}
degenerates to $x = x_u = G^{-1}a$ with an empty multiplier vector. The sign
condition holds vacuously, so this state satisfies the requirements of a dual
method with no work beyond one Cholesky factorisation. The objective there is
$f(x_u) = -\tfrac12 a\T x_u$. \emph{A dual method's phase-one problem is nothing
more than this.}
\end{remark}

\begin{example}[Two variables, one constraint]
Take $n = 2$, $G = I$, $a = 0$, and the single inequality
$c_1 = (1,1)\T$, $b_1 = 1$. Then $x_u = 0$, which is infeasible. With
$\Active = \{1\}$, $A = c_1$, the dual Hessian is $c_1\T c_1 = 2$, so
$u = (1 - 0)/2 = 1/2$ and $x = 0 + I c_1 (1/2) = (1/2,1/2)\T$. The multiplier is
positive and the point is feasible, so by Proposition~\ref{prop:sufficient} this
is the solution: the closest point of the half-space to the origin, as it should
be. Note that the certificate was checked, not assumed.
\end{example}

\section{The dual problem}

The following explains the name of the method in Chapter~\ref{ch:walking} and,
more usefully, why its objective value can be carried as a running total and why
that total increases.

\begin{proposition}[Dual function]
\label{prop:dual}
\index{duality!Lagrangian dual of the QP}
The Lagrangian dual of~\eqref{eq:qp} is the concave program $\max\,\psi(\lambda)$
over $\lambda_i \ge 0$ $(i > m_{\mathrm{eq}})$, where
\begin{equation}
  \psi(\lambda) = -\tfrac12 (a + C\lambda)\T G^{-1} (a + C\lambda) + b\T\lambda .
  \label{eq:dualfn}
\end{equation}
If $(x,\lambda)$ satisfies stationarity~\eqref{eq:stat} and
complementarity~\eqref{eq:comp}, then $\psi(\lambda) = f(x)$.
\end{proposition}

\begin{proof}
Minimising $L(\cdot,\lambda)$ over $x$ gives $x(\lambda) = G^{-1}(a + C\lambda)$
and, writing $s = a + C\lambda$,
$\psi(\lambda) = \tfrac12 s\T G^{-1}s - s\T G^{-1}s + b\T\lambda$, which
is~\eqref{eq:dualfn}. For the second claim, \eqref{eq:stat} gives $s = Gx$, and
\eqref{eq:comp} together with $c_i\T x = b_i$ on the equalities gives
$b\T\lambda = (Gx - a)\T x$; adding the two yields $f(x)$.
\end{proof}

So along any sequence of states of the form~\eqref{eq:subsol} the tracked
objective value is simultaneously the primal objective at the subproblem
minimiser and the dual objective at the current multipliers. That gives the
sandwich
\[
  \psi(\lambda) \;\le\; f(x^\star) \;\le\; f(\tilde x)
  \qquad\text{for any dual feasible $\lambda$ and any primal feasible $\tilde x$,}
\]
so a method whose iterates are of the form~\eqref{eq:subsol} approaches the
optimal value \emph{from below}. It is the mirror image of a primal method, which
maintains feasibility and descends to $f(x^\star)$ from above. Everything about
the trade between the two follows from this picture: a primal method can be
stopped early and return a feasible suboptimal point; a dual method cannot, but
it needs no phase one.

\begin{notes}
Proposition~\ref{prop:sufficient} and the closed form~\eqref{eq:subsol} are
standard; the presentation here follows \cite{schmelzer2026quadprog}, which states
them in exactly this form because every certificate in its account of the dual method
rests on them. Nocedal and Wright \cite[Ch.~16]{nocedal2006} give
the general treatment of quadratic programming, and \cite[Ch.~12]{nocedal2006} the
constraint qualifications that the linearity of our constraints lets us skip.
Boyd and Vandenberghe \cite{boyd2004} develop Lagrangian duality in the generality
of which Proposition~\ref{prop:dual} is a special case.

The bimodality observation about tolerance~\eqref{eq:certificate} is empirical
rather than a theorem, and the measurements behind it are reported in the
experiments of \cite{schmelzer2026quadprog}; Chapter~\ref{ch:numerics} returns to it.
\end{notes}

\begin{exercises}

\begin{exercise}
Prove that $A\T G^{-1} A \succ 0$ if and only if $A$ has full column rank, given
$G \succ 0$. Where in Proposition~\ref{prop:sub} is each direction used?
\end{exercise}

\begin{exercise}
Verify Proposition~\ref{prop:dual} directly for $n = 1$, $G = 1$, $a = 0$,
$C = 1$, $b = 1$, $m_{\mathrm{eq}} = 0$: write $\psi$ explicitly, maximise it over
$\lambda \ge 0$, and check that the maximiser reproduces the primal solution
through $x = G^{-1}(a + C\lambda)$.
\end{exercise}

\begin{exercise}
Show that the certificate test~\eqref{eq:certificate} costs $O(n^2 + nm)$ for
dense $G$ and $C$, and identify which term dominates for the long-only portfolio
problem of Chapter~\ref{ch:oneproblem}, where $m \approx n$.
\end{exercise}

\begin{exercise}
Give an example with $G \succeq 0$ singular in which the KKT conditions hold at
two distinct points. Which step of the proof of Proposition~\ref{prop:sufficient}
fails?
\end{exercise}

\begin{exercise}
Suppose an algorithm returns $(x,\lambda)$ with $\lambda_i < 0$ for exactly one
inequality $i \in \Active$, all other conditions holding exactly. Show that
dropping $i$ from the working set and re-solving~\eqref{eq:sub} strictly decreases
$f$. (This is the primal-method repair; the dual method of
Chapter~\ref{ch:walking} arranges never to be in this state.)
\end{exercise}

\begin{exercise}
The scaling of the tolerance $\varepsilon$ in~\eqref{eq:certificate} was left
vague. Propose a scaling that makes the test invariant under
$(G,a) \mapsto (\gamma G, \gamma a)$ for $\gamma > 0$, and another that makes it
invariant under $(c_i, b_i) \mapsto (\gamma c_i, \gamma b_i)$ per constraint.
Can both be achieved at once?
\end{exercise}

\begin{exercise}[Implementation]
Take the brute-force solver of Exercise~1.6 and replace its ``discard the
infeasible ones and compare objectives'' step by a certificate check: for each
candidate working set, form $(x,u)$ from~\eqref{eq:subsol} and test
\eqref{eq:certificate}. Confirm that exactly one working set passes on random
non-degenerate instances, and find an instance where more than one does. What is
special about it?
\end{exercise}

\end{exercises}

\chapter{The Geometry of the Active Set}
\label{ch:geometry}

Chapter~\ref{ch:certificate} reduced the QP to a search over working sets. This
chapter studies that search space: what a working set is geometrically, how the
optimality conditions read as a \emph{complementarity} problem, why the
bound-constrained case $C = I$ is structurally privileged, and what degeneracy is
and does. It contains no algorithms. It contains the facts that decide which
algorithms are possible.

\section{Faces, working sets, and the lattice}

The feasible set $\Omega$ is a polyhedron. Its faces are obtained by promoting
some subset of the inequalities to equalities: for
$\Active \subseteq \{m_{\mathrm{eq}}+1,\dots,m\}$,
\[
  F_\Active = \{x \in \Omega : c_i\T x = b_i \text{ for } i \in \Active\}
\]
is a face, possibly empty. These faces, ordered by inclusion, form a lattice, and
the solution of~\eqref{eq:qp} lies in the relative interior of exactly one of
them---namely $F_{\Active(x^\star)}$, where $\Active(x^\star)$ is the active set
of Definition~\ref{def:activeset}.
\index{face lattice}

So the QP is a two-stage problem: choose a face, then minimise a strictly convex
quadratic on the affine hull of that face, which is the subproblem~\eqref{eq:sub}.
The second stage is linear algebra. The first stage has
$2^{m - m_{\mathrm{eq}}}$ candidates, and every method in this book is a way of
visiting a vanishingly small number of them.

Two quantities bound the search. The first is elementary and pessimistic.

\begin{proposition}[Working-set size]
\label{prop:worksize}
If $A = C_{\cdot\Active}$ has full column rank then $|\Active| \le n$. In
particular the solution's working set never exceeds $n$ constraints, however
large $m$ is.
\end{proposition}

\begin{proof}
$A \in \R^{n\times k}$ with rank $k$ forces $k \le n$.
\end{proof}

The second is the observation that makes active-set methods practical: the
solution's active set is usually much smaller than $n$, or---in the
bound-constrained case---much \emph{larger}, in a way that is equally exploitable.
For a long-only portfolio the optimum is a vertex at which most non-negativity
bounds bind: on S\&P~500 data with $n = 494$ assets, roughly $450$ of the $495$
constraints are active at the minimum-variance solution and only $46$ assets are
held. Chapter~\ref{ch:krylov} exploits the small \emph{free} set; the synthetic
test families common in the literature, whose active sets are small, do not
exercise the same code paths, and Chapter~\ref{ch:measurement} argues that this
matters for benchmarking.

\section{Complementarity}

Specialise to bound constraints, $C = I$ and $b = 0$, so the problem is
\begin{equation}
  \min_{x \ge 0} \; f(x) = \tfrac12 x\T A x - b\T x, \qquad A \succ 0,
  \label{eq:bqp}
\end{equation}
where we have switched to the symbols $(A,b)$ that the literature on this case
uses; $A$ here plays the role of $G$ above. Introduce the multiplier
$s \ge 0$ for $x \ge 0$. Stationarity gives $s = Ax - b = \nabla f(x)$, the
\emph{reduced gradient}, and the KKT conditions collapse to
\index{linear complementarity problem}
\begin{equation}
  s = A x - b, \qquad x \ge 0, \qquad s \ge 0, \qquad x_i\,s_i = 0 \ \ \forall i.
  \label{eq:lcp}
\end{equation}

\begin{definition}[Linear complementarity problem]
Given $M \in \R^{n\times n}$ and $q \in \R^n$, the problem
$\mathrm{LCP}(M,q)$ is to find $x$ with $x \ge 0$, $Mx + q \ge 0$, and
$x\T(Mx+q) = 0$.
\end{definition}

So~\eqref{eq:lcp} is $\mathrm{LCP}(A,-b)$, and the bound-constrained QP and this
LCP are the same problem. That identification is not a curiosity: it imports a
theory, and Chapter~\ref{ch:guessing} spends the theory.

Complementarity partitions the indices. Write
\[
  \Free = \{i : x_i > 0\}
  \quad\text{(the \emph{free} set)},
  \qquad
  \Bound = \{i : x_i = 0\}
  \quad\text{(the \emph{bound} set)}.
\]
\index{free set}\index{bound set}
On $\Free$ we have $s_i = 0$, hence $(Ax)_i = b_i$; on $\Bound$ we have
$s_i \ge 0$. A bound variable with $s_i < 0$ would lower $f$ if released to a
small positive value; a free variable with $x_i < 0$ is infeasible and must be
pushed to the bound. Those two violations are precisely the primal and dual steps
of the algorithm in Chapter~\ref{ch:guessing}, and they are visible here, before
any algorithm exists, as the only two ways a candidate partition can be wrong.

Once $\Free$ is fixed the sign constraints drop out and what remains is the SPD
linear system
\begin{equation}
  A_{\Free}\,x_{\Free} = b_{\Free},
  \qquad A_\Free := A_{\Free\Free},
  \label{eq:innerspd}
\end{equation}
which is~\eqref{eq:sub} written for this case. Chapters~\ref{ch:krylov}
and~\ref{ch:operator} are about solving~\eqref{eq:innerspd} well.

\section{$P$-matrices: why bounds are privileged}

Here is the structural fact that separates the bound-constrained case from the
general one, and it is worth stating carefully because a great deal follows from
it.

\begin{definition}[$P$-matrix]
\index{P-matrix@$P$-matrix}
A matrix $M \in \R^{n\times n}$ is a \emph{$P$-matrix} if all its principal minors
are positive.
\end{definition}

\begin{proposition}[SPD implies $P$]
\label{prop:spdp}
If $A \succ 0$ then every principal submatrix $A_{\Free}$ is symmetric positive
definite, hence $\det A_\Free > 0$; so $A$ is a $P$-matrix.
\end{proposition}

\begin{proof}
For nonzero $x_\Free$, extend by zeros to $x \in \R^n$; then
$x_\Free\T A_\Free x_\Free = x\T A x > 0$. Symmetry is inherited. Positive
definiteness gives positive determinant.
\end{proof}

Two consequences, of quite different depth. The shallow one: \emph{every}
candidate free set gives a well-posed inner system~\eqref{eq:innerspd}, so an
algorithm may guess wildly and still have something to solve. The deep one:
$\mathrm{LCP}(M,q)$ with $M$ a $P$-matrix has a \emph{unique} solution for every
$q$, and that solution is reachable by principal pivoting in finitely many steps
\cite{murty1988,judice1994}. This is the hypothesis behind the finite-termination
theorem of Chapter~\ref{ch:guessing}.

Now the general case. With constraints $C\T x \ge b$ and a guessed working set
$\Active$, the system to be solved is
\begin{equation}
  \bigl(C_\Active\T G^{-1} C_\Active\bigr)\lambda_\Active
   = b_\Active - C_\Active\T x_u ,
  \label{eq:generalsystem}
\end{equation}
whose coefficient matrix is the dual Hessian of Definition~\ref{def:dualhessian}.
Compare the two cases:

\begin{center}
\begin{tabular}{lll}
\toprule
 & Bound constraints ($C = I$) & General constraints \\
\midrule
Working-set system & $(G^{-1})_{\Active\Active}$ & $C_\Active\T G^{-1} C_\Active$ \\
Structure & principal submatrix of $G^{-1}$ & Gram matrix of $C_\Active$ in $G^{-1}$ \\
Positive definite when & \emph{always} & $C_\Active$ has full column rank \\
That is a property of & the data & the \emph{guess} \\
\bottomrule
\end{tabular}
\end{center}

\index{dual Hessian!rank condition}
When $C = I$, equation~\eqref{eq:generalsystem} reads
$(G^{-1})_{\Active\Active}\lambda_\Active = \cdots$, a principal submatrix of
$G^{-1} \succ 0$, hence positive definite whatever $\Active$ is. When $C$ is
general, $C_\Active\T G^{-1} C_\Active$ is positive definite exactly when
$C_\Active$ has full column rank---and a guessed set may name $k > n$ constraints,
or $k \le n$ linearly dependent ones. There is no $P$-matrix to appeal to, and
correspondingly no theorem to inherit.

This is the single most important structural difference in the book. It is not a
gap in an argument or a missing anti-cycling rule; it is a property of the two
constraint classes. Chapter~\ref{ch:guessing} measures how often it bites (rarely
on well-posed random data, and in $93\%$ of instances once columns of $C$ are
duplicated), and Chapter~\ref{ch:walking} shows the elegant compensation: the dual
method's \emph{step rules} make linear dependence unreachable, so it needs no rank
test at all.

\section{Degeneracy}

\begin{definition}[Primal and dual degeneracy]
\index{degeneracy}
A feasible point $x$ is \emph{primal degenerate} if more constraints are active at
$x$ than the working set holds---equivalently, if a constraint outside $\Active$
has zero slack. A KKT point is \emph{dual degenerate} (or fails \emph{strict
complementarity}) if some $i \in \Active(x^\star)$ has $\lambda_i = 0$.
\end{definition}

Both are non-generic: each is a codimension-one condition on the data, so a
random instance misses them with probability one. Both nevertheless occur
constantly in practice, because real data is not random. Duplicated assets,
tied expected returns, a constraint entered twice by a model-assembly script, a
group budget that happens to coincide with a sum of individual bounds---each
produces exact ties.

The consequences are specific and worth separating.

\paragraph{Primal degeneracy stalls ascent arguments.} A method that proves
termination by showing the objective strictly increases at each step is defeated
by a step of length zero, and a zero step is exactly what primal degeneracy
permits: a constraint with zero slack can be added to the working set without
moving $x$ at all. The objective does not change, the ascent argument goes silent,
and a cycle is not excluded. This is the hypothesis in
Proposition~\ref{prop:ascent} of Chapter~\ref{ch:walking}.

\paragraph{Dual degeneracy blurs the identification of the active set.} If
$\lambda_i = 0$ for an active $i$, then $i$ may be in or out of the working set at
the solution without violating anything, so the working set at the optimum is not
unique. Methods that promise to ``identify the optimal face finitely'' generally
assume strict complementarity, and Chapter~\ref{ch:krylov} contrasts a family that
does with the method it develops, which does not.

\paragraph{Ties need a rule.} When several events would occur at the same
parameter value, choosing among them arbitrarily can cycle. The remedy is
inherited from linear programming: \emph{Bland's rule}, which breaks ties by
lowest index. It makes the walk deterministic and finite. Murty's least-index
rule for LCPs \cite{murty1988} is the same device, and it is what
Chapter~\ref{ch:guessing}'s termination theorem rests on. Lexicographic
perturbation is the other classical remedy.
\index{Bland's rule}

\begin{remark}[Degeneracy in floating point]
\label{rem:degen-float}
In finite precision the question changes character, because a slack of exactly
zero is not what one observes. What one observes is a slack of $10^{-16}$ or
$-10^{-16}$, and whether a constraint is ``active'' becomes a decision made by a
threshold. This is why Chapter~\ref{ch:numerics} treats degeneracy as a numerical
subject as much as a combinatorial one, and why the two tolerances---the one
deciding what to try next and the one deciding what to return---have very
different jobs and very different delicacy.
\end{remark}

\section{How many steps? The combinatorial question}

An active-set method changes the working set by one index per step, so its step
count is a walk length in the face lattice. Two facts frame what is known.

Empirically, the count is near-linear. A dual method on the QP typically takes
$O(n)$ outer steps; a parametric trace of a frontier typically has $O(n)$
breakpoints; a block-pivoting method converges in a handful of repairs almost
independently of $n$. These are the observations that make the subject practical.

In the worst case, it is not. Mairal and Yu \cite{mairal2012} construct LASSO
paths whose length is \emph{exponential} in the dimension, and the corresponding
bound for multi-parametric quadratic programs---exponentially many critical
regions---has been known in control since \cite{tondel2003}. The simplex method's
history is the obvious analogue: near-linear in practice, exponential on
Klee--Minty examples, and the gap between the two is still not fully explained.

\emph{What separates the benign average from the adversarial worst case is open},
in the shared form that covers all the instances of Part~III. It is the most
interesting unsolved problem this book touches, and it is stated here rather than
in a footnote because a graduate reader should know which parts of a subject are
settled.

\begin{notes}
Cottle, Pang and Stone \cite{cottle1992} is the standard reference for linear
complementarity, Murty \cite{murty1988} for principal pivoting and the least-index
rule, and Júdice and Pires \cite{judice1994} for the block-pivoting construction
that Chapter~\ref{ch:guessing} uses. The $P$-matrix characterisation of LCP
uniqueness is in both.

The $93\%$ figure, and the claim that the rank obstruction never occurs on
$600$ well-posed random instances across four constraint families, are measured in
the Goldfarb--Idnani note \cite{schmelzer2026quadprog}. The observation that a long-only optimum is a vertex holding
few names is from \cite{schmelzer2026matrixfree}.
\end{notes}

\begin{exercises}

\begin{exercise}
Show that the face $F_\Active$ is non-empty for at most $\binom{m}{n}$ of the
$2^m$ index sets when $\Omega$ is bounded and in general position. Why does this
not make enumeration practical?
\end{exercise}

\begin{exercise}
Prove that if $M$ is a $P$-matrix then $\mathrm{LCP}(M,q)$ has at most one
solution. (Hint: if $x \ne y$ both solve it, consider the index $i$ maximising
$|x_i - y_i|$ and the sign of $(M(x-y))_i(x-y)_i$.)
\end{exercise}

\begin{exercise}
Give a $2 \times 2$ example of a $P$-matrix that is not symmetric and not positive
definite. Deduce that Proposition~\ref{prop:spdp} is a strict implication.
\end{exercise}

\begin{exercise}
Let $C$ have two identical columns, $c_1 = c_2$, with $b_1 = b_2$. Show the
feasible set is unchanged by deleting one of them, but that a guessed working set
containing both makes~\eqref{eq:generalsystem} singular. This is the mechanism
behind the $93\%$ figure of Chapter~\ref{ch:guessing}.
\end{exercise}

\begin{exercise}
Construct a two-variable QP whose solution is primal degenerate: three constraints
passing through the optimum in $\R^2$. Verify that two different working sets both
satisfy the KKT conditions there, with different multipliers.
\end{exercise}

\begin{exercise}
For the bound-constrained problem~\eqref{eq:bqp}, show that $x$ solves it if and
only if $x = \max(x - \tau(Ax - b),\,0)$ for any $\tau > 0$, where the maximum is
componentwise. (This fixed-point characterisation is the basis of the projected
gradient method of Chapter~\ref{ch:krylov}.)
\end{exercise}

\begin{exercise}[Implementation]
Generate random long-only minimum-variance instances with $n = 20, 50, 100, 200$
from a factor model, solve each with any QP solver, and plot the number of assets
held (the free set size) against $n$. Confirm that it grows far more slowly than
$n$. Now repeat with $\Sig = I$ and explain the difference.
\end{exercise}

\end{exercises}

\chapter{The Linear Algebra of the Free Block}
\label{ch:freeblock}

Every method in this book solves the same linear system many times, with the
index set changing by one between consecutive solves. This chapter assembles the
linear algebra: the saddle-point form of the working-set subproblem, the two ways
to eliminate its multipliers, the Woodbury identity, and---the part that turns an
$O(n^4)$ algorithm into an $O(n^3)$ one---how to \emph{update} a factorisation
when one index joins or leaves.

Nothing here is new mathematics. It is assembled in one place because the
algorithms of Parts~II and~III are otherwise unreadable, and because the choice
between the alternatives below is where implementations differ.

\section{The saddle-point system}

Recall the subproblem~\eqref{eq:sub}. Writing its stationarity and feasibility
conditions together gives the \emph{saddle-point} or \emph{bordered} system
\index{saddle-point system}
\begin{equation}
  \begin{bmatrix} G & -A \\ A\T & 0 \end{bmatrix}
  \begin{bmatrix} x \\ u \end{bmatrix}
  =
  \begin{bmatrix} a \\ b_\Active \end{bmatrix},
  \qquad A = C_{\cdot\Active} \in \R^{n\times k}.
  \label{eq:saddle}
\end{equation}
The coefficient matrix is symmetric but \emph{indefinite}: it has $n$ positive and
$k$ negative eigenvalues when $G \succ 0$ and $A$ has full column rank. That
indefiniteness is why one does not simply run Cholesky on it, and why the two
elimination strategies below exist.

\begin{proposition}[Nonsingularity]
\label{prop:saddlenonsing}
If $G \succ 0$ and $A$ has full column rank then the matrix in~\eqref{eq:saddle}
is nonsingular.
\end{proposition}

\begin{proof}
Suppose $(v,w)$ lies in the kernel: $Gv = Aw$ and $A\T v = 0$. Then
$v\T G v = v\T A w = (A\T v)\T w = 0$, so $v = 0$ by definiteness of $G$, whence
$Aw = 0$ and $w = 0$ by full column rank.
\end{proof}

The same three-line argument recurs in Part~III for the bordered system of the
parametric theory (Lemma~\ref{lem:affine}), where $G$ is replaced by the reduced
Hessian on the free block. It is worth learning once.

\section{Two eliminations}

\subsection{Eliminate $x$: the dual Hessian route}

Solve the first block row for $x$ and substitute:
\[
  x = G^{-1}(a + Au)
  \quad\Longrightarrow\quad
  \underbrace{(A\T G^{-1} A)}_{k \times k}\,u = b_\Active - A\T G^{-1}a .
\]
This is Proposition~\ref{prop:sub} again. Cost: one Cholesky of $G$ ($O(n^3)$,
once), $k$ triangular solves to form $G^{-1}A$ ($O(n^2 k)$), and a $k\times k$
Cholesky ($O(k^3)$). The route is natural when $k$ is small and $G$ is factored
once and reused across many working sets---which is exactly the situation of
Chapters~\ref{ch:guessing} and~\ref{ch:warmstart}.

\subsection{Eliminate $u$: the reduced Hessian route}

Alternatively, let the columns of $Z \in \R^{n \times (n-k)}$ be a basis of
$\nullsp(A\T)$, write $x = x_p + Z y$ for any particular $x_p$ with
$A\T x_p = b_\Active$, and substitute into the objective. The result is an
unconstrained problem in $y$ with Hessian $Z\T G Z$, the \emph{reduced Hessian},
positive definite because $G$ is. Cost: forming $Z$, then an $(n-k)$-dimensional
solve.
\index{reduced Hessian}

Which is better depends on $k$ versus $n - k$, and the answer flips between the
two regimes this book cares about. A general QP with few active constraints has
$k \ll n$: eliminate $x$. A long-only portfolio at the optimum has $k \approx n$,
almost every bound binding, so $n - k$ is tiny: eliminate $u$, or equivalently
work on the free block directly, which is what Chapters~\ref{ch:krylov}
and~\ref{ch:operator} do.

The method of Chapter~\ref{ch:walking} does something cleverer than either: it
carries a factorisation from which \emph{both} the reduced inverse Hessian and the
dual Hessian's Cholesky factor can be read off, and updates it in $O(n^2)$ per
step. Section~\ref{sec:updating} below prepares that.

\section{Block elimination in the parametric setting}

Part~III needs a variant worth recording now. When the right-hand side splits into
a constant part and a $\lambda$-linear part, the system is solved for \emph{two}
right-hand sides at once, against one factorisation. With
\[
  Y = G_{\Free\Free}^{-1} C_\Free\T, \qquad
  z = G_{\Free\Free}^{-1} r_1,
\]
obtained from a single multi-column solve, the multipliers follow from the
$k \times k$ Schur complement $S = C_\Free\,G_{\Free\Free}^{-1} C_\Free\T = C_\Free Y$:
\begin{equation}
  \nu = S^{-1}\bigl(C_\Free z - r_2\bigr),
  \qquad
  x_\Free = z - Y\nu .
  \label{eq:blockelim}
\end{equation}
\index{Schur complement}
The dominant cost is the multi-right-hand-side solve against the free block; the
Schur solve is $k\times k$, and for the single budget constraint of a portfolio
problem $k = 1$, so it is a scalar division. This is the inner loop of the
Critical Line Algorithm (Chapter~\ref{ch:parametric}) and of the equality-augmented
solver of Chapter~\ref{ch:krylov}, in both cases with the same structure: the
multipliers are recovered \emph{analytically} and never appear as unknowns in the
linear solve, so every system actually solved stays symmetric positive definite.

That last sentence is the point of the manoeuvre. An indefinite system needs
$LDL\T$ with pivoting; an SPD system takes Cholesky, or conjugate gradients, or a
Woodbury formula. Keeping the solves SPD is what lets Chapter~\ref{ch:krylov} use
a Krylov method at all.

\section{The Woodbury identity}

\begin{proposition}[Sherman--Morrison--Woodbury]
\label{prop:woodbury}
\index{Woodbury identity}
For invertible $D \in \R^{n\times n}$ and $\Delta \in \R^{k\times k}$ and
$U \in \R^{n\times k}$,
\[
  (D + U\Delta U\T)^{-1}
   = D^{-1} - D^{-1}U\bigl(\Delta^{-1} + U\T D^{-1} U\bigr)^{-1} U\T D^{-1},
\]
whenever the inverses on the right exist.
\end{proposition}

\begin{proof}
Multiply the right-hand side by $D + U\Delta U\T$ and simplify, using
$\Delta U\T D^{-1}U = (\Delta^{-1} + U\T D^{-1}U)\Delta - I$ to collapse the cross
terms.
\end{proof}

The identity converts an $n\times n$ inverse into a $k\times k$ one, and it is the
engine of Chapter~\ref{ch:operator}. Its importance here is that
\emph{diagonal-plus-low-rank} is the form that structured models actually take: a
factor risk model $\Sig = \diag(d) + U\Lambda U\T$ with $k \ll n$ factors; a
random-matrix-cleaned correlation $\bar\mu I + U_k\Delta_k U_k\T$; a Nyström kernel
approximation; a spiked covariance model. All four are the same object, and all
four invert in $O(nk^2 + k^3)$ instead of $O(n^3)$.

One subtlety that catches implementers. When the identity is applied to a
\emph{restricted} block $\Active$, the restricted factor $U_\Active = U[\Active,:]$
is generally \emph{not} orthonormal even if $U$ was:
$U_\Active\T U_\Active \ne I_k$. The correction matrix
$W = \Delta^{-1} + U_\Active\T D_\Active^{-1} U_\Active$ accounts for this, and it
is what must be formed and factored per solve. Its cost, $O(n_a k^2 + k^3)$, is
what Proposition~\ref{prop:crossover} in Chapter~\ref{ch:operator} weighs against
assembling and factoring the $n_a \times n_a$ block directly.

\section{Updating a factorisation}
\label{sec:updating}

An active-set method changes the working set by one index per step. Recomputing a
factorisation from scratch costs $O(nk^2)$; \emph{updating} it costs $O(n^2)$ or
$O(nk)$. Over the $O(n)$ steps of a run that is the difference between $O(n^4)$
and $O(n^3)$, and it is the entire reason factorisations are carried rather than
rebuilt.

Both updates below work the same way. Suppose we carry a matrix $J$ satisfying
$J\T G J = I$ together with an upper triangular $R$ with $J\T A = [R;0]$
(Chapter~\ref{ch:walking} explains why these are the right invariants). The first
condition is preserved by post-multiplying $J$ with \emph{any} orthogonal $W$,
since $(JW)\T G (JW) = W\T W = I$. So the update's only task is to choose $W$ so
that the triangular shape of the second condition is restored. Because
$J' = JW$ gives ${J'}\T A = W\T(J\T A)$, the same $W\T$ acts on the rows of $R$.

\subsection{Insertion}
\index{factorisation update!insertion}

Append a column $n_*$ to $A$. With $d = J\T n_*$ split as $(d_1,d_2)$ at index $k$,
\[
  J\T [\,A \ \ n_*\,] = \begin{bmatrix} R & d_1 \\ 0 & d_2 \end{bmatrix},
\]
which fails to be triangular only in the trailing $n-k-1$ entries of $d_2$. Choose
orthogonal $W_2$ with $W_2\T d_2 = \alpha e_1$ and set $W = \diag(I_k, W_2)$. Then
the new $R$ is $R$ bordered by $d_1$ and $\alpha$, and the leading $k$ columns of
$J$ are untouched.

Two implementations of $W_2$ compete. A \emph{chain of Givens rotations}, one per
trailing component, is $n-k-1$ operations of $O(n)$ each. A single
\emph{Householder reflection}
\begin{equation}
  W_2 = I - \frac{2}{\beta}\, v v\T,
  \qquad v = d_2 - \alpha e_1,
  \qquad \beta = v\T v,
  \qquad \alpha = \pm\norm{d_2},
  \label{eq:householder}
\end{equation}
achieves the same reduction in two operations of $O(n(n-k))$: a matrix--vector
product and a rank-one update. The arithmetic is comparable; the number of
operations is not, and on modern hardware the operation count is what one pays.

\begin{remark}[The sign of $\alpha$]
\label{rem:sign}
Numerical analysis prescribes $\alpha = -\sign((d_2)_1)\norm{d_2}$, making
$v_1 = (d_2)_1 - \alpha$ a sum of like-signed terms. A code that takes the opposite
sign---to match a Givens chain's sign convention, say---walks into the very
cancellation the prescription exists to avoid, severely when $d_2$ is already close
to a positive multiple of $e_1$. The cure is algebraic, not a change of convention:
with $\tau = \norm{(d_2)_{2:}}^2$ and $\nu = \norm{d_2}$,
\begin{equation}
  v_1 = (d_2)_1 - \sign((d_2)_1)\,\nu
      = \frac{(d_2)_1^2 - \nu^2}{(d_2)_1 + \sign((d_2)_1)\nu}
      = \frac{-\sign((d_2)_1)\,\tau}{|(d_2)_1| + \nu},
  \label{eq:nocancel}
\end{equation}
in which every operation combines like-signed quantities. Chapter~\ref{ch:numerics}
returns to this pattern of rewriting a difference as a quotient.
\end{remark}

\subsection{Deletion}
\index{factorisation update!deletion}

Delete column $p$ of $A$. This deletes column $p$ of $R$ and shifts the rest left,
leaving a matrix that is upper triangular except for one subdiagonal, with
offending entries $\tilde R_{j+1,j}$ for $j = p,\dots,k-1$. These lie in
\emph{distinct} columns, and no single reflection annihilates components of
distinct vectors. What restores the shape is a \emph{chase}: for $j = p,\dots,k-1$
in order, a $2\times2$ orthogonal acting on rows $j,j+1$ annihilates
$\tilde R_{j+1,j}$, and the same transformation is applied to columns $j, j+1$
of $J$.

A useful implementation trick: take the $2\times 2$ block to be the \emph{symmetric}
reflection
\begin{equation}
  W_j^{(2)} = \begin{bmatrix} c & s \\ s & -c \end{bmatrix},
  \qquad c = \frac{x}{h},\quad s = \frac{y}{h},
  \qquad h = \sign(x)\sqrt{x^2+y^2},
  \label{eq:refl2}
\end{equation}
rather than a Givens rotation. The update applies $W$ to the columns of $J$ and
$W\T$ to the rows of $R$; symmetry makes those the same matrix, so one $2\times2$
serves both sides. A Givens rotation is not symmetric and needs its transpose on
the $R$ side. The two differ by a sign, and Proposition~\ref{prop:invariance} in
Chapter~\ref{ch:walking} shows that difference is free.

\begin{remark}[Why the reflections cannot be blocked]
\label{rem:wy}
\index{WY representation}
A sequence of Householder reflections is normally applied as a block: the WY
representation \cite{bischof1987,schreiber1989} writes $W_1\cdots W_j$ as
$I + WY\T$, so $j$ reflections reach a matrix in one level-3 update rather than
$j$ level-2 ones. That is why blocked QR is fast, and it is the first thing to
reach for here.

It is not available, for a structural reason. Blocking requires the $j$
reflections to be known before any is applied. In an active-set iteration the
reflection that inserts a constraint is determined by $d = J\T n_*$, which depends
on $J$ \emph{after} the previous insertion; and between two insertions the solver
reads quantities off the updated factors to decide which constraint enters next,
and whether it enters at all. The updates are interleaved with reads of what they
produce, so there is no window in which two reflections are simultaneously known
and unapplied. The same obstruction governs the deletion chase, whose
transformations read their parameters from entries the previous one wrote.

\emph{The method therefore stays a level-2 computation no matter how it is coded.}
This is a real ceiling, and it is one reason Chapter~\ref{ch:measurement} argues
that iteration counts, not flop counts, are the currency at these sizes.
\end{remark}

\section{Conditioning, briefly}

The spectral condition number $\kappa(G) = \lambda_{\max}(G)/\lambda_{\min}(G)$
governs both the accuracy of a direct solve and the iteration count of an
iterative one. Two facts are used repeatedly later.

\begin{proposition}[Cauchy interlacing, the consequence we need]
\label{prop:interlace}
If $G \succ 0$ then $\lambda_{\min}(G_{\Free\Free}) \ge \lambda_{\min}(G)$ and
$\lambda_{\max}(G_{\Free\Free}) \le \lambda_{\max}(G)$ for every index set $\Free$.
Hence $\kappa(G_\Free) \le \kappa(G)$.
\end{proposition}

\begin{proof}
Both bounds follow from the Rayleigh quotient characterisation restricted to
vectors supported on $\Free$.
\end{proof}

So a free-block solve is never worse conditioned than the full problem---a small
mercy that Chapter~\ref{ch:krylov} leans on when it bounds an inner CG iteration
count uniformly over the working sets a loop might visit.

\begin{proposition}[Weyl, the consequence we need]
\label{prop:weyl}
\index{Weyl's inequality}
For $\alpha \in (0,1)$ and $R\T R \succ 0$, the convex combination
$G_\alpha = (1-\alpha)G + \alpha R\T R$ satisfies
\[
  \kappa(G_\alpha) \;\le\;
  \frac{(1-\alpha)\lambda_{\max}(G) + \alpha\,\lambda_{\max}(R\T R)}
       {\alpha\,\lambda_{\min}(R\T R)},
\]
which decreases monotonically in $\alpha$. For $R\T R = \bar\lambda I$ the bound is
$\tfrac{1-\alpha}{\alpha}\,\lambda_{\max}(G)/\bar\lambda + 1$.
\end{proposition}

\begin{proof}
Weyl's inequality \cite{hornjohnson2013} bounds $\lambda_{\max}(G_\alpha)$ above by
the sum of the two terms' maxima and $\lambda_{\min}(G_\alpha)$ below by
$\alpha\lambda_{\min}(R\T R)$, the first term contributing at least zero.
\end{proof}

This is the whole mathematical content of ``regularisation improves conditioning'',
and Chapter~\ref{ch:conditioning} spends it three times over: on Tikhonov
regularisation, on Ledoit--Wolf shrinkage, and on the observation that a
statistically motivated shrinkage intensity conditions the solver as a free
by-product.

\begin{notes}
Golub and Van Loan \cite{golub2013} covers saddle-point systems, Householder and
Givens transformations, and factorisation updating; Benzi, Golub and Liesen
\cite{benzi2005} is the survey of saddle-point problems specifically. Higham
\cite{higham2002} is the reference for the cancellation analysis behind
Remark~\ref{rem:sign}. The WY representation is \cite{bischof1987,schreiber1989};
the argument in Remark~\ref{rem:wy} that it is structurally unavailable here is
from the Goldfarb--Idnani note \cite{schmelzer2026quadprog}.
\end{notes}

\begin{exercises}

\begin{exercise}
Show that the saddle-point matrix in~\eqref{eq:saddle} has exactly $n$ positive and
$k$ negative eigenvalues when $G \succ 0$ and $A$ has full column rank. (Hint: use
the block $LDL\T$ factorisation with the Schur complement $-A\T G^{-1}A$.)
\end{exercise}

\begin{exercise}
Count the flops of the two eliminations of Section~2 for given $n$ and $k$,
assuming $G$ is already factored. At what $k/n$ do they cross?
\end{exercise}

\begin{exercise}
Prove Proposition~\ref{prop:woodbury} by verifying the product both ways. Then
specialise to $k = 1$ to recover the Sherman--Morrison formula.
\end{exercise}

\begin{exercise}
Let $U \in \R^{n\times k}$ have orthonormal columns and let $\Active$ index $n_a$
rows. Show $\norm{U_\Active\T U_\Active}_2 \le 1$, with equality iff some unit
vector in $\range(U\T)$ is supported on $\Active$. Explain why the correction
matrix $W$ of Section~4 is nevertheless always invertible when $\Delta \succ 0$.
\end{exercise}

\begin{exercise}
Implement the Householder insertion~\eqref{eq:householder} both with the naive
$v_1 = (d_2)_1 - \sign((d_2)_1)\norm{d_2}$ and with the rewriting~\eqref{eq:nocancel}.
Construct $d_2$ close to a positive multiple of $e_1$ and measure the relative
error in $\norm{W_2\T d_2 - \alpha e_1}$ for both.
\end{exercise}

\begin{exercise}
Show that the deletion chase of Section~5.2 requires exactly $k - p$ plane
transformations, and that each is determined by entries written by the previous
one. Conclude that no reordering makes it parallel.
\end{exercise}

\begin{exercise}[Implementation]
Write a routine that maintains a Cholesky factor of $A\T G^{-1}A$ under insertion
and deletion of columns of $A$, and a second routine that recomputes it from
scratch. On random data with $n = 200$, insert $100$ columns one at a time and
compare total runtime and final accuracy (measured against a fresh factorisation).
At what point does the accumulated error in the updated factor become visible?
\end{exercise}

\end{exercises}


\part{Finding the Active Set}
\chapter{Walking to It: The Dual Active-Set Method}
\label{ch:walking}

This chapter derives the method of Goldfarb and Idnani
\cite{goldfarb1982,goldfarb1983}, which for forty years has been the method of
choice for dense quadratic programs of small to moderate size---$n$ from a handful
to a few thousand. It is often presented as a sequence of updates, which makes it
look more intricate than it is. Read from its invariants it is short: one matrix
is carried, two identities define it, and \emph{everything} the iteration computes
is a consequence of those two lines.

The method's distinguishing property is that it needs no phase-one problem. A
primal active-set method maintains a feasible iterate and works toward
stationarity, so it must begin somewhere feasible, and finding a feasible point is
a linear program in its own right. The dual method reverses the roles: it maintains
stationarity and the multiplier signs throughout, and works toward feasibility.
Its starting point is free, because with no constraint active the stationarity
condition reads $Gx = a$ and the multiplier vector is empty, so
$x = G^{-1}a$ satisfies everything the method asks (Remark~\ref{rem:cold}). One
Cholesky factorisation and one triangular solve replace the phase-one problem
entirely.

The price is that the iterate is infeasible until the last step, so the method
cannot be stopped early for an approximate answer. In exchange it terminates at an
\emph{exactly} feasible point of the active-set lattice rather than approaching one
asymptotically, and it warm-starts almost perfectly---which is why
Chapter~\ref{ch:warmstart} is built on it.

\section{The matrix the solver carries}
\label{sec:factorisation}

Proposition~\ref{prop:sub} gives the subproblem solution in closed form, and one
could implement the method by evaluating it afresh every iteration. That costs a
$k\times k$ factorisation of the dual Hessian per step, $O(nk^2 + k^3)$, and over
the $O(n)$ steps a dual method takes it comes to $O(n^4)$. The whole point of the
Goldfarb--Idnani formulation is to replace this with $O(n^2)$ per step, and it does
so by carrying a single matrix and updating it.

\subsection{Two invariants}

Let $G = R_G\T R_G$ be the Cholesky factorisation with $R_G$ upper triangular, and
set $J_0 := R_G^{-1}$. Then $J_0 J_0\T = G^{-1}$ and $J_0\T G J_0 = I$. The
solver's state is a matrix $J \in \R^{n\times n}$ and an upper triangular
$R \in \R^{k\times k}$ satisfying
\index{invariants (Goldfarb--Idnani)}
\begin{align}
  J\T G J &= I_n, \label{eq:inv1}\\
  J\T A &= \begin{bmatrix} R \\ 0 \end{bmatrix}. \label{eq:inv2}
\end{align}

Read these slowly; everything below is a consequence. Equation~\eqref{eq:inv1} is
equivalent to $JJ\T = G^{-1}$ with $J$ nonsingular, and says the columns of $J$ are
a basis of $\R^n$ orthonormal in the inner product
$\inner{v}{w}_G = v\T G w$. Equation~\eqref{eq:inv2} says that in that basis the
working-set normals are triangular. Partition
\begin{equation}
  J = \begin{bmatrix} J_1 & J_2 \end{bmatrix},
  \qquad J_1 \in \R^{n\times k},\quad J_2 \in \R^{n\times(n-k)};
  \label{eq:split}
\end{equation}
the lower block of~\eqref{eq:inv2} reads $J_2\T A = 0$.

At the cold start $\Active = \emptyset$, so \eqref{eq:inv2} is vacuous and
$J = J_0$ serves. The updates of Section~\ref{sec:updates} maintain both invariants
by post-multiplying $J$ with orthogonal matrices, exactly as
Section~\ref{sec:updating} of Chapter~\ref{ch:freeblock} set up.

\begin{proposition}[What the blocks mean]
\label{prop:blocks}
Suppose \eqref{eq:inv1}--\eqref{eq:inv2} hold with $R$ nonsingular. Then
\begin{enumerate}
  \item[(i)] $R\T R = A\T G^{-1} A$, so $R$ is a Cholesky factor of the dual
        Hessian, unique up to the signs of its rows;
  \item[(ii)] $\range(J_2) = \nullsp(A\T)$, and the columns of $J_2$ are a
        $G$-orthonormal basis of it;
  \item[(iii)] $J_1 = G^{-1} A R^{-1}$ and $\range(J_1) = \range(G^{-1}A)$;
  \item[(iv)] $J_1 J_1\T = G^{-1}A(A\T G^{-1}A)^{-1}A\T G^{-1}$ and consequently
        \begin{equation}
          J_2 J_2\T = G^{-1} - G^{-1}A\bigl(A\T G^{-1}A\bigr)^{-1}A\T G^{-1}.
          \label{eq:reduced}
        \end{equation}
\end{enumerate}
\end{proposition}

\begin{proof}
(i) $A\T G^{-1} A = A\T J J\T A = R\T R$ by~\eqref{eq:inv2}, and two upper
triangular matrices with the same Gram matrix differ by a left factor
$\diag(\pm1)$.

(ii) $J_2\T A = 0$ places $\range(J_2) \subseteq \nullsp(A\T)$, with equality by
dimension since $J$ is nonsingular and $A$ has full column rank; the basis is
$G$-orthonormal because $J_2\T G J_2 = I_{n-k}$ is the trailing block
of~\eqref{eq:inv1}.

(iii) From~\eqref{eq:inv1}, $J^{-1} = J\T G$ and hence $J^{-\top} = G J$.
Inverting~\eqref{eq:inv2}, $A = J^{-\top}[R;0] = G J_1 R$, so
$J_1 = G^{-1}AR^{-1}$, whose range is that of $G^{-1}A$ since $R$ is invertible.

(iv) Substituting $J_1 = G^{-1}AR^{-1}$ and $R\T R = A\T G^{-1}A$,
\[
  J_1 J_1\T = G^{-1}A R^{-1} R^{-\top} A\T G^{-1}
            = G^{-1}A (A\T G^{-1}A)^{-1} A\T G^{-1},
\]
and \eqref{eq:reduced} follows from $J_1J_1\T + J_2J_2\T = JJ\T = G^{-1}$.
\end{proof}

The matrix in~\eqref{eq:reduced} is the \emph{reduced inverse Hessian} of the
working set: it inverts $G$ on the feasible subspace $\nullsp(A\T)$ and annihilates
the complement. So carrying $J$ means carrying that object in factored form, at no
cost beyond the $n\times n$ array---which turns the search directions of the next
section into single matrix--vector products.
\index{reduced inverse Hessian}

\begin{remark}[Why fold the two factorisations together]
An alternative state is $R_G$ together with an explicit orthogonal factor $Q$ of the
QR factorisation of $R_G^{-\top}A$. That holds the same information ($J = J_0Q$) and
stores no more numbers. But every use of $J$ below would become a triangular solve
followed by an orthogonal transformation---two operations where the folded form has
one, and the solve is the shape that does not reach a tuned kernel. The folded form
costs $J$ its triangularity after the first update, which is exactly the trade:
$n^2$ storage and $2n^2$ flops per product, in exchange for the products being a
plain \texttt{gemv} at every iteration whatever the working set does.
\end{remark}

\subsection{The representation is not unique, and need not be}

Invariants~\eqref{eq:inv1}--\eqref{eq:inv2} do not determine $J$ and $R$. By
Proposition~\ref{prop:blocks}, $R$ is pinned only up to a left factor
$S = \diag(\pm1)$, and $J_2$ may be any $G$-orthonormal basis of $\nullsp(A\T)$, so
it is pinned only up to a right factor $V$ with $V\T V = I$. The freedom is exactly
\begin{equation}
  (J_1, J_2, R) \;\longmapsto\; (J_1 S,\; J_2 V,\; S R).
  \label{eq:freedom}
\end{equation}
A Givens chain and a Householder reflection annihilating the same vector produce
factors differing in sign; a differently ordered sequence of insertions and
deletions produces a different $J_2$ outright. The following says none of it
matters, so an implementation may choose its reduction on grounds of speed alone.

\begin{proposition}[Invariance of the iterates]
\label{prop:invariance}
Let $(J,R)$ and $(\hat J,\hat R)$ both satisfy \eqref{eq:inv1}--\eqref{eq:inv2} for
the same $G$ and $A$. Then for every $n_* \in \R^n$ the quantities
\begin{equation}
  d = J\T n_*, \qquad z = J_2 d_2, \qquad r = R^{-1} d_1
  \label{eq:dzr}
\end{equation}
with $d = (d_1,d_2)$ split at $k$ satisfy $\hat z = z$ and $\hat r = r$. If the two
representations differ by~\eqref{eq:freedom} with $V = I$, the equalities hold
exactly in IEEE~754 arithmetic.
\end{proposition}

\begin{proof}
By Proposition~\ref{prop:blocks}(i) and $A\T J = [R\T\ 0]$ we have
$A\T G^{-1} n_* = R\T d_1$, so
\begin{equation}
  r = R^{-1}d_1 = (R\T R)^{-1} A\T G^{-1} n_*
    = \bigl(A\T G^{-1}A\bigr)^{-1} A\T G^{-1} n_* ,
  \label{eq:rclosed}
\end{equation}
which mentions only $A$, $G$, $n_*$. Likewise $z = J_2 J_2\T n_*$, and
by~\eqref{eq:reduced}
\begin{equation}
  z = \Bigl(G^{-1} - G^{-1}A\bigl(A\T G^{-1}A\bigr)^{-1}A\T G^{-1}\Bigr) n_*,
  \label{eq:zclosed}
\end{equation}
again free of the representation. For the last claim, $\hat R = SR$ gives
$\hat d_1 = S d_1$ and $\hat r = (SR)^{-1}Sd_1 = R^{-1}d_1$, in which the only
floating-point operations beyond those forming $r$ are multiplications by $\pm1$;
these are exact.
\end{proof}

\begin{remark}[What the proposition does and does not settle]
It settles that the iteration's \emph{decisions}---which constraint is most
violated, whether the step is full or partial, which constraint is dropped---are
functions of $z$ and $r$ alone, and those are functions of $(G,A,n_*)$ alone. It
does \emph{not} settle that two implementations follow the same trajectory in
finite precision, because the rounding \emph{within} the reductions differs and the
decisions are discontinuous functions of the computed values. That is an empirical
question. What is checkable against a single implementation is the weaker
statement: that the computed $z$ and $r$ agree with the closed
forms~\eqref{eq:zclosed} and~\eqref{eq:rclosed} to machine precision.
\end{remark}

\section{One iteration}
\label{sec:iteration}

The state at the top of an outer iteration is a working set $\Active$ of size $k$,
the factors $(J,R)$, the subproblem solution $x$ and multipliers $u$, and the
running objective $f(x)$. The invariant maintained throughout is
\begin{equation}
  \text{$(x,u)$ solves the subproblem~\eqref{eq:sub}, and}
  \quad u_i \ge 0 \ \text{ for every inequality } i \in \Active,
  \label{eq:invariant}
\end{equation}
which by the discussion after Proposition~\ref{prop:sub} is stationarity,
complementarity, and dual feasibility together. Only primal feasibility is missing,
and the iteration exists to obtain it.

\subsection{Choosing the entering constraint}

Compute the slacks $s_i = c_i\T x - b_i$. If they pass, $x$ is the solution by
Proposition~\ref{prop:sufficient}. Otherwise select the \emph{most violated}
constraint, measured relative to the norm of its own normal:
\begin{equation}
  v_i = \begin{cases} |s_i|, & i \le m_{\mathrm{eq}},\\ -s_i, & i > m_{\mathrm{eq}},\end{cases}
  \qquad
  i_* \in \argmax_i \; \frac{v_i}{\norm{c_i}},
  \qquad\text{terminating when } \max_i \frac{v_i}{\norm{c_i}} \le 0.
  \label{eq:select}
\end{equation}

The scaling is not cosmetic.

\begin{proposition}[Selection is scale invariant]
\label{prop:scale}
Replacing $(c_i,b_i)$ by $(\gamma c_i, \gamma b_i)$ for $\gamma > 0$ leaves the
feasible set, the solution, and the choice made by~\eqref{eq:select} unchanged.
\end{proposition}

\begin{proof}
The rescaled constraint has the same solution set. Its slack scales as
$s_i \mapsto \gamma s_i$ and its normal as $\norm{c_i}\mapsto\gamma\norm{c_i}$, so
the ratio is unchanged, as are the other constraints' ratios.
\end{proof}

Without the division, a constraint written in basis points rather than in units
would be selected first merely for having been written differently, and the
solver's trajectory---though not its answer---would depend on the caller's choice
of units.

\subsection{The two directions}

Write $n_* = c_{i_*}$, $b_* = b_{i_*}$, and form $d$, $z$, $r$ as
in~\eqref{eq:dzr}. The following lemma is the computational core of the method.

\begin{lemma}[Properties of the directions]
\label{lem:directions}
With $d$, $z$, $r$ as in~\eqref{eq:dzr},
\begin{enumerate}
  \item[(i)] $A\T z = 0$: moving along $z$ preserves the working-set equalities;
  \item[(ii)] $n_*\T z = \norm{d_2}^2 = z\T G z \ge 0$, with equality iff $z = 0$;
  \item[(iii)] $G z = n_* - A r$;
  \item[(iv)] $z = 0$ if and only if $n_* = A r$, i.e.\ iff $n_*$ lies in the span
        of the working-set normals, in which case $r$ holds its coordinates there.
\end{enumerate}
\end{lemma}

\begin{proof}
(i) $A\T z = (J_2\T A)\T d_2 = 0$ by~\eqref{eq:inv2}. (ii)
$n_*\T z = (J_2\T n_*)\T d_2 = d_2\T d_2$, and
$z\T G z = d_2\T(J_2\T G J_2)d_2 = d_2\T d_2$ by the trailing block
of~\eqref{eq:inv1}; $J_2$ has full column rank, so $z = 0$ iff $d_2 = 0$. (iii) By
\eqref{eq:reduced}, $Gz = n_* - A(A\T G^{-1}A)^{-1}A\T G^{-1}n_*$, and the second
term is $Ar$ by~\eqref{eq:rclosed}. (iv) If $d_2 = 0$ then
$J\T n_* = (d_1,0) = (J\T A)r$ and $J$ is nonsingular, so $n_* = Ar$; conversely
$n_* = Ar$ gives $Gz = 0$ by (iii).
\end{proof}

Part (ii) identifies $z$ as an ascent direction for the entering constraint's
slack, which grows at rate $n_*\T z > 0$ per unit step. Part (i) says the step
costs nothing in working-set feasibility. Together with the closed forms the
reading is:
\begin{itemize}
  \item $z$ is the $G$-orthogonal projection of the unconstrained direction
        $G^{-1}n_*$ onto the working set's feasible subspace $\nullsp(A\T)$;
  \item $r$ solves a least-squares problem in the $G^{-1}$ metric,
        \begin{equation}
          r = \argmin_{y \in \R^k} \; (Ay - n_*)\T G^{-1} (A y - n_*),
          \label{eq:lsq}
        \end{equation}
        the best approximation of the entering normal by the working-set normals,
        with residual $z = G^{-1}(n_* - Ar)$.
\end{itemize}

\begin{remark}[Why not call a least-squares routine]
Equation~\eqref{eq:lsq} is a genuine least-squares problem and it is tempting to
hand it to a library. Resist: the factor $R$ of its normal equations is
\emph{already held}, so \eqref{eq:dzr} costs one triangular solve, where a library
call would refactorise from scratch at $O(nk^2)$ every iteration. The same applies
to~\eqref{eq:subsol}. \emph{The closed form is the specification, not the
implementation}---a distinction worth carrying out of this book.
\end{remark}

\subsection{The affine path}

Let $u_*$ denote the entering constraint's own multiplier, starting at $0$.

\begin{proposition}[Dual feasibility along the path]
\label{prop:path}
For $t \in \R$ define
\begin{equation}
  x(t) = x + t z, \qquad u(t) = u - t r, \qquad u_*(t) = u_* + t .
  \label{eq:path}
\end{equation}
Then for all $t$: $A\T x(t) = b_\Active$, and stationarity holds for the working set
augmented by the entering constraint,
$G x(t) - a = A\, u(t) + u_*(t)\, n_*$.
\end{proposition}

\begin{proof}
The first claim is Lemma~\ref{lem:directions}(i). For the second, stationarity at
$t = 0$ reads $Gx - a = Au + u_* n_*$, and by Lemma~\ref{lem:directions}(iii),
$Gx(t) - a = Au + u_*n_* + t(n_* - Ar) = A(u-tr) + (u_*+t)n_*$.
\end{proof}

So the \emph{entire path is stationary}, with the entering multiplier rising at
unit rate and the working-set multipliers falling linearly along $-r$.
Complementarity holds for every constraint except the entering one, whose slack is
still nonzero while its multiplier is positive; that single violation is what the
step removes. Two limits on $t$ follow.

\paragraph{The dual limit $t_1$.} Dual feasibility requires $u_i(t) \ge 0$ for every
inequality $i \in \Active$. The binding indices are those with $r_i > 0$, and
\begin{equation}
  t_1 = \min\Bigl\{\, u_i/r_i \;:\; i \in \Active,\ i > m_{\mathrm{eq}},\ r_i > 0 \,\Bigr\},
  \label{eq:t1}
\end{equation}
with $t_1 = +\infty$ when the set is empty. Equalities are excluded: their
multipliers are unrestricted in sign. This is the classical ratio test.

\paragraph{The primal limit $t_2$.} The entering slack is $s_{i_*} + t\,n_*\T z$
with $s_{i_*} < 0$, reaching zero at
\begin{equation}
  t_2 = \frac{|s_{i_*}|}{n_*\T z} = \frac{|s_{i_*}|}{\norm{d_2}^2},
  \label{eq:t2}
\end{equation}
finite and positive whenever $z \ne 0$, and $t_2 = +\infty$ when $z = 0$, in which
case the primal iterate cannot move at all and only the multipliers change.

\paragraph{The step.} Take $t = \min(t_1,t_2)$. If $t_2 \le t_1$ the step is
\emph{full}: the entering constraint joins the working set with multiplier
$u_* + t_2$ and the factorisation is updated. If $t_1 < t_2$ the step is
\emph{partial}: the multiplier of the binding index reaches zero first, that
constraint leaves, the factorisation is downdated, $z$ and $r$ are recomputed, and
the inner loop repeats with the same entering constraint. If both limits are
infinite the problem is infeasible---see Proposition~\ref{prop:farkas}.

\begin{remark}[Equalities violated from above]
\label{rem:reverse}
An equality may be violated in either direction. When $s_{i_*} > 0$ the slack must
be \emph{decreased}, so the step is taken along $-z$: $t$ ranges over the negative
reals, $t_2 = -|s_{i_*}|/n_*\T z$, the ratio test runs against $-r$, and the
entering multiplier accumulates negatively. Nothing else changes, which is why
Proposition~\ref{prop:path} was stated for all $t \in \R$.
\end{remark}

\subsection{The objective, in closed form}

\begin{proposition}[Objective increment]
\label{prop:increment}
Along the path~\eqref{eq:path},
\begin{equation}
  f(x(t)) - f(x) = t\Bigl(\frac{t}{2} + u_*\Bigr)\, n_*\T z ,
  \label{eq:increment}
\end{equation}
strictly positive whenever $z \ne 0$, $t \ne 0$, and $t$, $u_*$ share a sign (in
particular whenever $t > 0$ and $u_* \ge 0$).
\end{proposition}

\begin{proof}
$f(x + tz) = f(x) + t\, z\T(Gx - a) + \tfrac12 t^2\, z\T G z$. Stationarity and
Lemma~\ref{lem:directions}(i) give $z\T(Gx-a) = u_*\,n_*\T z$;
Lemma~\ref{lem:directions}(ii) gives $z\T G z = n_*\T z$. Positivity follows since
$n_*\T z > 0$ and $t(t/2+u_*) > 0$ under the stated sign condition.
\end{proof}

The increment is exact, so the objective is carried as a running total at the cost
of one multiplication per step instead of an $O(n^2)$ re-evaluation, and it is
expressed in quantities the iteration already holds ($n_*\T z$ is needed for $t_2$
anyway). It handles the reversed step of Remark~\ref{rem:reverse} with no special
case: there $t < 0$ and $u_*$ accumulates negatively, so $t/2 + u_* < 0$ and the
product with $t$ is again positive.

\begin{algorithm}[ht]
\caption{The Goldfarb--Idnani dual active-set method}\label{alg:gi}
\begin{algorithmic}[1]
\Require $G \succ 0$, $a$, $C$, $b$, $m_{\mathrm{eq}}$
\State $J \leftarrow \mathrm{chol}(G)^{-1}$; \ $x \leftarrow G^{-1}a$; \
       $f \leftarrow -\tfrac12 a\T x$; \ $\Active \leftarrow \emptyset$; \
       $k \leftarrow 0$; \ $u \leftarrow ()$
\Loop
  \State $s \leftarrow C\T x - b$, with $s_i \leftarrow 0$ for $i \in \Active$;
         choose $i_*$ by~\eqref{eq:select}
  \State \textbf{if} none \textbf{then} \Return $(x,f,u)$
         \Comment{optimal, by Prop.~\ref{prop:sufficient}}
  \State $n_* \leftarrow c_{i_*}$; \ $u_* \leftarrow 0$
  \Loop
    \State $d \leftarrow J\T n_*$; \ $z \leftarrow J_2 d_2$; \
           $r \leftarrow R^{-1} d_1$
    \State $t_1, i_{\mathrm{del}} \leftarrow$ \eqref{eq:t1}; \
           $t_2 \leftarrow$ \eqref{eq:t2};
           \textbf{if} both infinite \textbf{then stop}: infeasible
    \State $t \leftarrow \min(t_1,t_2)$; \ $x \leftarrow x + tz$; \
           $f \leftarrow f + t(t/2 + u_*)\,n_*\T z$; \
           $u \leftarrow u - tr$; \ $u_* \leftarrow u_* + t$
    \If{$t_2 \le t_1$}
      \State $\Active \leftarrow \Active \cup \{i_*\}$, $k \leftarrow k+1$;
             insert; \textbf{break}
    \Else
      \State $\Active \leftarrow \Active \setminus \{i_{\mathrm{del}}\}$,
             $k \leftarrow k-1$; delete
    \EndIf
  \EndLoop
\EndLoop
\end{algorithmic}
\end{algorithm}

\section{Termination}
\label{sec:termination}

Three questions are separate and best kept so: does the inner loop stop, does the
outer loop stop, and what does it mean when the iteration gets stuck.

\begin{proposition}[The inner loop terminates]
\label{prop:inner}
For a fixed entering constraint with $n_* \ne 0$, the inner loop performs at most
$k+1$ passes, where $k$ is the working-set size on entry.
\end{proposition}

\begin{proof}
Every non-exiting pass removes one constraint, so after at most $k$ of them
$\Active = \emptyset$. Then $J_2 = J$ and $z = JJ\T n_* = G^{-1}n_*$, so
$n_*\T z = n_*\T G^{-1}n_* > 0$; thus $z \ne 0$, $t_2$ is finite, and
$t_1 = +\infty$ because the ratio test ranges over an empty set. The step is full
and the loop exits.
\end{proof}

The excluded case $n_* = 0$ is a column of $C$ that is identically zero. Such a
constraint reads $0 \ge b_*$: either vacuous ($b_* \le 0$) or rendering the problem
infeasible. An implementation scores it as infinitely violated when $b_* > 0$,
routing it to the infeasibility verdict rather than to a division by zero.

\begin{proposition}[Strict ascent and finiteness]
\label{prop:ascent}
Every outer iteration ending in a full step with $t_2 \ne 0$ strictly increases the
objective, hence by Proposition~\ref{prop:dual} the dual value. If every full step
has $t_2 > 0$, the method terminates after finitely many outer iterations.
\end{proposition}

\begin{proof}
The iteration's total increment is the sum of~\eqref{eq:increment} over its inner
passes, each positive since the accumulated $u_*$ and each $t$ share a sign
throughout, and the final term nonzero because $t_2 \ne 0$. A working set
determines $(x,u)$ uniquely by Proposition~\ref{prop:sub}, hence determines the
objective; that value strictly increases from one outer iteration to the next, so
no working set recurs. There are finitely many subsets of $\{1,\dots,m\}$.
\end{proof}

\begin{remark}[Degeneracy]
\label{rem:degeneracy}
The hypothesis fails exactly when a full step of length zero occurs, which requires
$s_{i_*} = 0$: more constraints pass through $x$ than the working set holds. A zero
step then adds one without changing $x$, $u$, or the objective, the ascent argument
goes silent, and a cycle is not excluded. This is primal degeneracy, exactly as
Chapter~\ref{ch:geometry} described it; the classical remedies are lexicographic or
least-index rules.
\end{remark}

\subsection{Infeasibility is a certificate, not a failure}

The interesting case is a stuck iteration: $z = 0$, so the primal cannot move, and
$t_1 = \infty$, so no multiplier limits the step. The classical reading is that the
dual is unbounded and therefore the primal infeasible. That argument needs care---by
Proposition~\ref{prop:path} the whole ray is dual feasible, and
by~\eqref{eq:increment} with $z = 0$ the value is constant. The cleanest statement
avoids the dual entirely.

\begin{proposition}[Farkas certificate]
\label{prop:farkas}
\index{Farkas certificate}
Suppose the iteration reaches a state in which the entering inequality $i_*$ is
violated, $s_{i_*} < 0$, and (a) $z = 0$, and (b) $r_i \le 0$ for every inequality
$i \in \Active$. Then $\Omega = \emptyset$.
\end{proposition}

\begin{proof}
By (a) and Lemma~\ref{lem:directions}(iv), $n_* = Ar = \sum_{i\in\Active} r_i c_i$.
Let $\tilde x \in \Omega$. Splitting $\Active$ into equalities $\mathcal{E}$ and
inequalities $\mathcal{I}$,
\[
  n_*\T \tilde x
  = \sum_{i \in \mathcal{E}} r_i\, c_i\T \tilde x
    + \sum_{i \in \mathcal{I}} r_i\, c_i\T \tilde x
  \;\le\; r\T b_\Active ,
\]
since equality members contribute $c_i\T\tilde x = b_i$ exactly and each inequality
member contributes $r_i c_i\T \tilde x \le r_i b_i$ because $c_i\T\tilde x \ge b_i$
and $r_i \le 0$. On the other hand $A\T x = b_\Active$, so
$r\T b_\Active = r\T A\T x = n_*\T x = s_{i_*} + b_* < b_*$. Combining,
$n_*\T\tilde x < b_*$, contradicting $\tilde x \in \Omega$.
\end{proof}

So the solver's infeasibility verdict is a \emph{proof} rather than an exhausted
search, and the proof is a vector the solver computed for its own purposes.
Hypothesis (a) is $t_2 = \infty$ and (b) is $t_1 = \infty$, so the condition an
implementation tests---neither step limit being finite---is literally the hypothesis
of the proposition. This is the second certificate of the book, and it deserves the
same emphasis as the first: an algorithm that returns ``infeasible'' should return
the reason.

\begin{remark}[The remaining hypothesis is not free in floating point]
\label{rem:spurious}
Proposition~\ref{prop:farkas} also needs the entering constraint to be
\emph{genuinely} violated. Exact arithmetic gives that; floating point does not,
and the gap is reachable by the very substitution Section~\ref{sec:updates} argues
for. A Householder reflection and a Givens chain leave the iterate at slightly
different points, so a constraint one implementation leaves just inside its
boundary another can leave the same distance outside. Reaching the stuck state with
such a constraint selected, a solver must neither enforce it, there being nothing
to enforce, nor conclude infeasibility, the conclusion not following; the remedy is
to set it aside for the current iterate and take the next candidate, discarding that
judgement as soon as the iterate moves. The tolerance separating rounding from a
real violation is comparatively easy to set, because infeasibility is macroscopic:
its size is set by the geometry of the constraints, not by the arithmetic.
\end{remark}

\section{Updating the factorisation}
\label{sec:updates}

The two updates are those of Section~\ref{sec:updating} in
Chapter~\ref{ch:freeblock}, applied to $(J,R)$. Insertion appends the entering
normal, costing $O(n^2)$; deletion chases a subdiagonal out of $R$, costing
$O(nk)$. What is new here is a consequence of the step rules.

\begin{proposition}[Full column rank is maintained]
\label{prop:rank}
The insertion is performed only after a full step, which requires $t_2 < \infty$
and hence $z \ne 0$. Then $\alpha \ne 0$ in~\eqref{eq:householder}, so $R'$ is
nonsingular and $A'$ has full column rank $k+1$.
\end{proposition}

\begin{proof}
$|\alpha| = \norm{d_2}$, nonzero iff $z \ne 0$ by Lemma~\ref{lem:directions}(ii).
$R'$ is triangular with diagonal that of $R$ extended by $\alpha$, hence
nonsingular, and ${R'}\T R' = {A'}\T G^{-1}A'$ is then positive definite, forcing
full column rank.
\end{proof}

\emph{The algorithm therefore never has to test for linear dependence among the
working-set normals: the step rules exclude it.} A constraint whose normal lies in
the span of those already held has $z = 0$ by Lemma~\ref{lem:directions}(iv), so its
step is limited by $t_1$ and it can only be added after enough partial steps have
removed the dependence.

This is the structural advantage that Chapter~\ref{ch:guessing} gives up, and it is
worth pausing on. Chapter~\ref{ch:geometry} identified the rank of $C_\Active$ as
the one thing that can go wrong with a general-constraint working-set system, and
observed that it is a property of the \emph{guess}. Here there is no guess: the
walk cannot reach a dependent set. Rank is protected not by a test but by the
geometry of the step rules---an elegance that is easy to overlook and expensive to
lose.

\section{Summary}

One matrix $J$ is carried, defined by $J\T G J = I$ and $J\T A = [R;0]$. Every
quantity the iteration needs is one matrix--vector product or one triangular solve
away, and every quantity it \emph{consumes} has a closed form in $(G,A,n_*)$ alone,
which is why the representation may be chosen for speed and why two implementations
holding different factors compute the same step. The iteration is one affine path
along which stationarity holds identically; two limits cut it short, and which binds
decides whether a constraint is added or dropped. The objective advances by an exact
closed form, positive in both step directions, so the method is a dual ascent and,
absent degeneracy, finite. When the path is blocked in both directions the vector
$r$ is a Farkas certificate.

\begin{notes}
The method is \cite{goldfarb1982,goldfarb1983}; this presentation, organised around
the two invariants, follows \cite{schmelzer2026quadprog}, the mathematics note
accompanying the \texttt{cvx-quadprog} package \cite{cvxquadprog}. The restriction to $G \succ 0$ was
lifted by Boland \cite{boland1997}, who proves finite termination in the
semidefinite case; Forsgren, Gill and Wong \cite{forsgren2016} prove finite
termination for paired primal and dual active-set methods more generally. The modern
successor is a different factorisation rather than a reimplementation: Arnström,
Bemporad and Axehill \cite{arnstrom2022daqp} carry the same dual walk on recursive
$LDL\T$ updates, handle semidefinite $G$ by proximal-point iterations, and ship it as
library-free C for embedded control. Arnström and Axehill
\cite{arnstrom2022certification} compute the \emph{exact} worst-case iteration count
of active-set methods on parametric QP families---stronger than
Proposition~\ref{prop:ascent} and different in kind, being computational and
per-family rather than general and conditional.
\end{notes}

\begin{exercises}

\begin{exercise}
Verify directly from~\eqref{eq:inv1} that $J^{-1} = J\T G$, and use it to give a
one-line proof that $J$ is nonsingular whenever $G$ is.
\end{exercise}

\begin{exercise}
Show that $\Pi := J_2 J_2\T G$ is idempotent, has range $\nullsp(A\T)$, and is
self-adjoint for $\inner{\cdot}{\cdot}_G$. Conclude that $z = \Pi\,G^{-1}n_*$ as
claimed after Lemma~\ref{lem:directions}.
\end{exercise}

\begin{exercise}
Prove that the multiplier of a rescaled constraint scales as
$\lambda_i \mapsto \lambda_i/\gamma$ under $(c_i,b_i)\mapsto(\gamma c_i,\gamma b_i)$,
and deduce that Proposition~\ref{prop:scale} extends to the sign conditions.
\end{exercise}

\begin{exercise}
Work Algorithm~\ref{alg:gi} by hand on $n = 2$, $G = I$, $a = (0,0)\T$, and the two
constraints $x_1 \ge 1$, $x_1 + x_2 \ge 3$. Record $\Active$, $x$, $u$, $t_1$, $t_2$,
and $f$ at every step. Which step is partial?
\end{exercise}

\begin{exercise}
Show that the total objective increase over a whole run equals
$f(x^\star) - f(G^{-1}a)$, and interpret this in terms of
Proposition~\ref{prop:dual}.
\end{exercise}

\begin{exercise}
Construct a two-variable infeasible instance and trace Algorithm~\ref{alg:gi} to the
stuck state. Exhibit the Farkas certificate $r$ explicitly and verify the inequality
in the proof of Proposition~\ref{prop:farkas}.
\end{exercise}

\begin{exercise}
Give an example in which a constraint is added to the working set and later removed.
Why does Proposition~\ref{prop:ascent} not forbid this?
\end{exercise}

\begin{exercise}[Implementation]
Implement Algorithm~\ref{alg:gi}. Verify Proposition~\ref{prop:invariance}
empirically: at each iteration compute $z$ and $r$ both from the carried factors and
from the closed forms~\eqref{eq:zclosed}--\eqref{eq:rclosed}, and report the largest
discrepancy over a run. Then verify Proposition~\ref{prop:increment} by comparing the
running objective against a fresh evaluation of $f(x)$ at every step.
\end{exercise}

\begin{exercise}[Implementation]
Modify your solver to rescale one constraint by $\gamma \in \{10^{-3},\dots,10^{3}\}$
and record the number of additions and removals for each $\gamma$.
Proposition~\ref{prop:scale} claims only that the first \emph{choice} is unchanged;
what do you observe about the whole trajectory?
\end{exercise}

\end{exercises}

\chapter{Guessing It: Primal--Dual Active Sets and Block Pivoting}
\label{ch:guessing}

The dual method of Chapter~\ref{ch:walking} reaches the active set one constraint
at a time. That is its defining feature and, at moderate $n$, its principal cost: a
budget-plus-bounds problem in a few hundred variables takes on the order of a
hundred outer iterations, each a handful of matrix--vector products.

The alternative is to guess the whole active set at once, solve the working-set
subproblem it induces, and repair the guess from the signs that come back. This
chapter develops that idea, proves it terminates finitely in the bound-constrained
case, and shows precisely why the proof does not survive the general constraint
class. The last point is the chapter's main contribution to the reader's judgement:
it is a case study in \emph{recognising that a theorem's hypothesis is structural},
not a technicality to be patched.

\section{The iteration}

Given a candidate working set $\Active$, Proposition~\ref{prop:sub} gives the
solution directly:
\begin{equation}
  \bigl(C_\Active\T G^{-1} C_\Active\bigr)\lambda_\Active
   = b_\Active - C_\Active\T x_u,
  \qquad
  x = x_u + G^{-1}C_\Active \lambda_\Active,
  \qquad x_u = G^{-1}a,
  \label{eq:pdassolve}
\end{equation}
with $\lambda_i = 0$ off $\Active$. Both solves go through one Cholesky
factorisation of $G$, computed once for the whole run, and one of the $k\times k$
dual Hessian, computed per repair.

The repair rule reads the two sign conditions
\eqref{eq:dfeas}--\eqref{eq:pfeas} as instructions:
\index{primal--dual active set}
\begin{equation}
  \Active' = \{1,\dots,m_{\mathrm{eq}}\} \;\cup\;
  \bigl\{\, i > m_{\mathrm{eq}} \;:\;
    (i \in \Active \text{ and } \lambda_i \ge 0)
    \ \text{ or }\
    (c_i\T x < b_i) \,\bigr\}.
  \label{eq:repair}
\end{equation}
In words: an active constraint whose multiplier has gone negative does not belong
and leaves; an inactive constraint that is violated at $x$ does belong and joins;
equalities are always in. The starting guess is the equalities together with
whatever the unconstrained minimiser $x_u$ violates---which is already the answer
when it violates nothing. If $\Active' = \Active$ the KKT conditions hold to the
tolerance used and~\eqref{eq:pdassolve} supplies stationarity by construction.

The whole set is exchanged at once, and in practice this converges in a handful of
repairs \emph{almost independently of $n$}. It trades arithmetic, which is abundant
at these sizes, for iterations, which are not---a trade Chapter~\ref{ch:measurement}
argues is the right one below a few hundred variables and the wrong one above a few
thousand.

\begin{example}[Why the guess is usually good]
For a long-only minimum-variance problem, $x_u = G^{-1}a$ is the unconstrained
minimum-variance portfolio, which typically has a few dozen negative weights out of
several hundred assets. Those are exactly the bounds the first guess activates,
and most of them are correct. The repairs correct the rest.
\end{example}

\section{Two readings of the same iteration}

The iteration above has two lineages, and each contributes something.

\paragraph{Semismooth Newton.} Hintermüller, Ito and Kunisch \cite{hintermuller2002}
show that the primal--dual active set strategy is a semismooth Newton method applied
to the complementarity condition written as a nonsmooth equation
$\lambda - \max(0, \lambda - \sigma(c_i\T x - b_i)) = 0$. This explains the observed
behaviour: Newton's method converges in few iterations and its iteration count does
not scale with the dimension. It also explains the failure mode---Newton's method
without globalisation may not converge at all.

\paragraph{Block principal pivoting.} Specialise to bound constraints, so that by
Chapter~\ref{ch:geometry} the problem is $\mathrm{LCP}(A,-b)$. Toggling one index
between free and bound is a \emph{principal pivot}; toggling several at once is a
\emph{block} principal pivot \cite{judice1994,kim2011}. The repair
rule~\eqref{eq:repair} is exactly a block pivot driven by signs.
\index{block principal pivoting}

The second reading is the one that supports a theorem, so we develop it.

\section{Finite termination for bound constraints}

Recall the bound-constrained problem~\eqref{eq:bqp} and its free set $\Free$. The
loop solves $A_\Free x_\Free = b_\Free$, drops free variables that came back
negative, admits bound variables whose reduced gradient is negative, and repeats.

Block pivots are fast, but---like any all-at-once exchange---they can cycle: a batch
drop can overshoot the support and a later batch add restore it, revisiting a free
set already seen. The remedy is due to Júdice and Pires \cite{judice1994}: run the
batch exchange as a \emph{fast path} guarded by an anti-cycling counter, and fall
back to a single lowest-index pivot when the fast path stalls. That single pivot is
Murty's least-index principal pivoting method \cite{murty1988}, the LCP analogue of
Bland's rule, and it cannot cycle.

\begin{algorithm}[ht]
\caption{Guarded block principal pivoting}\label{alg:elim}
\begin{algorithmic}[1]
\Require Inner solver returning $x_{\Free}$ for a free set $\Free$; SPD $A$ (or its
         mat-vec) and $b$; tolerance $\varepsilon$; patience $p_{\max}$ (default $3$)
\Ensure The minimiser $x$ of~\eqref{eq:bqp}
\State $\Free \leftarrow \{1,\ldots,n\}$;\; $\bar N \leftarrow n+1$;\; $p \leftarrow p_{\max}$
\Loop
  \State solve $A_\Free x_\Free = b_\Free$;\; $x_i \leftarrow 0$ for $i \notin \Free$
  \State $s \leftarrow Ax - b$
    \hfill\textit{(reduced gradient)}
  \State $\mathcal{D} \leftarrow \{i \in \Free : x_i < -\varepsilon\}$;\;
         $\mathcal{V} \leftarrow \{i \notin \Free : s_i < -\varepsilon\}$
    \hfill\textit{(primal, dual violators)}
  \State $N \leftarrow |\mathcal{D} \cup \mathcal{V}|$
  \If{$N = 0$}
    \State \textbf{break} \hfill\textit{(KKT~\eqref{eq:lcp} holds; globally optimal)}
  \ElsIf{$N < \bar N$ \textbf{ or } $p > 0$}
    \State \textbf{if} $N < \bar N$ \textbf{ then } $\bar N \leftarrow N$;\; $p \leftarrow p_{\max}$
           \textbf{ else } $p \leftarrow p - 1$
    \State $\Free \leftarrow (\Free \setminus \mathcal{D}) \cup \mathcal{V}$
      \hfill\textit{(batch exchange)}
  \Else
    \State $i^\star \leftarrow \min\{i : i \in \mathcal{D}\cup\mathcal{V}\}$;
           toggle $i^\star$
      \hfill\textit{(Bland least-index single pivot)}
  \EndIf
\EndLoop
\end{algorithmic}
\end{algorithm}

\begin{theorem}[Finite convergence]
\label{thm:termination}
\index{finite termination!block pivoting}
Let $f(x) = \tfrac12 x\T A x - b\T x$ with $A \succ 0$. For any patience budget
$p_{\max} \in \mathbb{N}$, Algorithm~\ref{alg:elim} terminates in finitely many outer
iterations and returns the unique global minimiser of $\min_{x\ge0} f(x)$,
\emph{even under degeneracy}.
\end{theorem}

\begin{proof}
\emph{Step 1 (each inner solve is well posed).} $A \succ 0$ makes every principal
submatrix $A_\Free$ SPD (Proposition~\ref{prop:spdp}), so the inner system has a
unique solution.

\emph{Step 2 (termination implies optimality).} Free-set stationarity gives
$s_i = 0$ on $\Free$. The loop exits only when all free $x_i \ge 0$ and all bound
$s_i \ge 0$, which with complementarity is exactly~\eqref{eq:lcp}. Strict convexity
makes these sufficient (Proposition~\ref{prop:sufficient}). It therefore suffices to
show the loop cannot run forever.

\emph{Step 3 (the fallback cannot cycle).} Call a maximal block of consecutive
\textbf{else} iterates a \emph{fallback run}. Within a run $\bar N$ is constant and
each iterate performs one principal pivot on the lowest-indexed violator. This is
Murty's least-index method for the $P$-matrix LCP of Step~2, which generates a
sequence of complementary bases in which none recurs and which reaches the solution
in finitely many pivots \cite{murty1988}. A fallback run is an initial segment of
that sequence, so each run is finite wherever it starts.

\emph{Step 4 (finitely many batch steps and runs).} $\bar N$ is non-increasing and
strictly decreases exactly when the progress branch fires, so that branch fires at
most $n+1$ times. Each remaining fast-path step decrements $p$, and $p$ resets only
on progress, so at most $p_{\max}$ such steps occur between progress events.
A fallback run ends at termination or at a progress step, and there are at most
$n+1$ progress steps. Hence at most $(n+1) + (n+2)p_{\max}$ fast-path steps and at
most $n+2$ fallback runs.

\emph{Step 5.} Their sum is finite. A crude ceiling is the number of complementary
bases, $2^n$.
\end{proof}

The $2^n$ ceiling establishes \emph{finiteness} only; it is not an efficiency
estimate. The guarantee is nevertheless unconditional---the least-index fallback is
what removes the non-degeneracy hypothesis that earlier active-set arguments need,
and it is the fallback, not the batch exchange, that the proof rests on.

In practice the fallback is rarely reached. A batch step fails to reduce $N$ only
when a dropped index is immediately re-added or vice versa, which forces $s_i = 0$
exactly for some toggled $i$---a codimension-one and hence non-generic condition
\cite{judice1994}. So block pivoting converges in a handful of outer steps on
typical data with the fast path alone certifying optimality, and the guarantee is
\emph{dormant} on generic data yet genuinely necessary: on adversarial
anti-correlated designs the unguarded batch path provably cycles and only the
fallback terminates it.

\section{Why the theorem does not survive general constraints}
\label{sec:no-theorem}

Now the negative result, which is the chapter's point.

Consider $\min_{x\ge0}\tfrac12 x\T G x - a\T x$, so $C = I$ and $m = n$. The system
solved on a candidate set is
\begin{equation}
  \bigl(I_{\cdot\Active}\T G^{-1} I_{\cdot\Active}\bigr)\lambda_\Active
  = \bigl(G^{-1}\bigr)_{\Active\Active}\lambda_\Active
  = b_\Active - x_{u,\Active},
  \label{eq:principal}
\end{equation}
whose coefficient matrix is a \emph{principal submatrix} of $G^{-1}$. Since
$G^{-1} \succ 0$, every principal submatrix of it is positive definite, so $G^{-1}$
is a $P$-matrix: every principal pivot is well defined, \emph{whatever the guess}.
That is the hypothesis of Theorem~\ref{thm:termination}.

None of it is available for general $C$. The matrix $C_\Active\T G^{-1} C_\Active$
is positive definite exactly when $C_\Active$ has full column rank, and that is a
property of the guess, not of the data: a guessed set may name $k > n$ constraints,
or $k \le n$ linearly dependent ones, and then the pivot is undefined. There is no
$P$-matrix to appeal to, and correspondingly no theorem to inherit.

Contrast Proposition~\ref{prop:rank}: the dual method's step rules make dependence
\emph{unreachable}, so no rank test is needed anywhere. Guessing forfeits that
protection. The exposure is structural, not a missing anti-cycling rule, and no
tie-breaking rule repairs it, because the failure is not a cycle but an undefined
step.

\begin{remark}[Is this a real hazard or a technicality?]
\label{rem:93}
It is easy to dismiss, and the evidence says one should not. Over six hundred
random instances across four constraint families---well-posed problems of the kind
this literature routinely tests on---the rank failure \emph{never occurs}, which is
precisely why it can be overlooked. Duplicating columns of $C$, as any
programmatically assembled model may do (the same sector cap entered twice by two
model components; a constraint generated once per group and once globally), raises
it to $93\%$ of instances, and the certified fraction then collapses from $100\%$ to
nothing.

The lesson generalises past this solver: \emph{a hazard that never fires on the
family you test on is not thereby absent}. It is invisible. Test families in a
literature tend to share the structure that makes the literature's methods look
good, and finding out which structure that is, is part of reading a paper.
\end{remark}

\section{The certificate, again}

Because the method may fail---and may fail by stabilising on a set that is simply
not optimal---the trade is sound only because the answer is checked. Here
Proposition~\ref{prop:sufficient} carries the argument: a candidate $(x,\lambda)$
satisfying the tolerance test~\eqref{eq:certificate} \emph{is} the answer to that
tolerance. A candidate that fails is discarded and the exact walk of
Algorithm~\ref{alg:gi} runs instead, so guessing can never degrade the answer: it
produces the minimiser or nothing.

\begin{remark}[Two tolerances with very different jobs]
\label{rem:twotol}
The tolerance in the repair rule~\eqref{eq:repair} decides which set is tried next.
A bad choice there costs a repair or a fallback and can \emph{never} cost
correctness, so it is not delicate.

The tolerance in the certificate~\eqref{eq:certificate} is the only thing standing
between a non-optimal point and the caller, and is therefore the one to reason
about. It is nonetheless not delicate either, for the bimodality reason of
Chapter~\ref{ch:certificate}: candidates built from well-conditioned working sets
satisfy it to a few multiples of machine epsilon, and candidates built on sets that
are not have no reason to come close.

Learning to tell these two kinds of tolerance apart is worth more than any
particular tolerance value. Chapter~\ref{ch:numerics} makes it a general principle.
\end{remark}

The implementation treats the Cholesky factorisation of
$C_\Active\T G^{-1}C_\Active$ as the rank test---a dependent guess makes it fail,
rather than returning something plausible---and the least-index rule as a way of
making progress where block exchange oscillates, not as a termination proof.
Neither substitutes for the guarantee. The certificate does.

\section{Inexact inner solves}

Theorem~\ref{thm:termination} assumes each inner call returns the \emph{exact}
minimiser on the free set. Chapter~\ref{ch:krylov} will replace that call with
conjugate gradients, which returns an iterate accurate only to its residual
tolerance. Does the guarantee survive?

The question is not rhetorical: a tolerance-level solve can flip a sign test, and a
flipped sign test reintroduces exactly the cycling the fallback exists to prevent.
The answer, developed as Lemma~\ref{lem:inexact} in Chapter~\ref{ch:krylov}, is that
it does survive, provided the inner residual is tied to a quantity called the
\emph{decision margin}: the smallest nonzero magnitude any primal or dual sign test
encounters along the exact trajectory. Once the residual is small against that
margin, the inexact loop makes the identical sign decisions as the exact one and
therefore visits the identical free sets, so the termination count transfers
decision for decision.

That is a good example of the right way to compose two guarantees. One does not
prove a new termination theorem for the inexact method; one proves that the inexact
method \emph{is} the exact method, on the trajectory in question.

\begin{notes}
The primal--dual active set strategy as a semismooth Newton method is
\cite{hintermuller2002}; block principal pivoting and the least-index guard are
\cite{judice1994,murty1988,kim2011}, with \cite{cottle1992} the reference for
linear complementarity generally. Theorem~\ref{thm:termination} is the Júdice--Pires
construction restated in active-set language, following \cite{schmelzer2026nncg};
no novelty is claimed for it. The rank obstruction of
Section~\ref{sec:no-theorem} and the $600$-instance and $93\%$ measurements of
Remark~\ref{rem:93} are from the Goldfarb--Idnani note \cite{schmelzer2026quadprog}.
\end{notes}

\begin{exercises}

\begin{exercise}
Show that the starting guess ``equalities plus whatever $x_u$ violates'' returns the
correct answer in one step whenever $x_u$ is feasible.
\end{exercise}

\begin{exercise}
Verify that the repair rule~\eqref{eq:repair} is a fixed point exactly at a KKT
point, i.e.\ $\Active' = \Active$ iff \eqref{eq:pfeas}--\eqref{eq:comp} hold for the
pair produced by~\eqref{eq:pdassolve}.
\end{exercise}

\begin{exercise}
In Algorithm~\ref{alg:elim}, show that $\bar N$ can decrease at most $n+1$ times,
and construct a run in which it decreases exactly twice.
\end{exercise}

\begin{exercise}
Prove that if a batch step fails to reduce $N$ then some toggled index $i$ has
$s_i = 0$ or $x_i = 0$ exactly. Explain why this makes stalling non-generic.
\end{exercise}

\begin{exercise}
Take $G = I$ and $C$ with $c_1 = c_2 = (1,1)\T$, $b_1 = b_2 = 1$, $n = 2$. Show that
the guess $\Active = \{1,2\}$ makes $C_\Active\T G^{-1} C_\Active$ singular, and that
the dual method of Chapter~\ref{ch:walking} never forms this set.
\end{exercise}

\begin{exercise}
Design an adversarial family for which the unguarded batch path cycles. (Hint: make
the dropped index's reduced gradient turn negative as soon as it is dropped, which
anti-correlated columns achieve.)
\end{exercise}

\begin{exercise}[Implementation]
Implement Algorithm~\ref{alg:elim} with a direct free-set solve. On random SPD
problems with a planted non-negative optimum, record the number of outer steps as
$n$ grows from $50$ to $2000$. Confirm it is essentially flat. Then disable the
fallback ($p_{\max} = \infty$) and run on the adversarial family of Exercise~6 to
observe a cycle.
\end{exercise}

\begin{exercise}[Implementation]
Instrument your solver from Chapter~\ref{ch:walking} to accept a guessed working set
as a fast path, falling back to the exact walk when the certificate fails. Measure
the fraction of random instances on which the fast path succeeds, and the speedup
when it does. Then duplicate a random column of $C$ and repeat, reproducing
Remark~\ref{rem:93}.
\end{exercise}

\end{exercises}

\chapter{Iterating Inside It: Krylov Methods on the Free Block}
\label{ch:krylov}

Chapters~\ref{ch:walking} and~\ref{ch:guessing} took the inner linear solve for
granted. This chapter takes it seriously.

The observation that organises the chapter is this. Once a free set $\Free$ is
fixed, the bound-constrained problem reduces to the unconstrained SPD system
$A_\Free x_\Free = b_\Free$---and the canonical solver for an SPD system is the
conjugate gradient method of Hestenes and Stiefel \cite{hestenes1952}. CG needs
only the action $v \mapsto A_\Free v$, never the matrix. So if $A$ arrives in
factored form---$A = M\T M$ for a data matrix $M$, which is exactly what a Gram
matrix or a sample covariance is---then the whole method can run without ever
assembling an $n \times n$ object.

CG solves linear systems; it does not, unaided, respect $x \ge 0$. Wrapping it in
the active-set loop of Chapter~\ref{ch:guessing} supplies the missing layer. The
composition is not free: the loop's termination theorem assumes exact inner solves
and CG returns an approximation. Section~\ref{sec:inexact} closes that gap, and it
is the most instructive argument in the chapter.

\section{Conjugate gradients, in one page}

For SPD $A$ and right-hand side $b$, CG builds its $k$-th iterate as the minimiser
of the energy $\tfrac12 x\T A x - b\T x$ over the Krylov subspace
$\Krylov_k(A,b) = \mathrm{span}\{b, Ab, \dots, A^{k-1}b\}$, and does so with one
matrix--vector product, a handful of vector operations, and $O(n)$ storage per
iteration. That it can optimise over a growing subspace at constant cost per step
is the content of the three-term recurrence, and any of
\cite{golub2013,greenbaum1997,trefethenbau1997} derives it.
\index{conjugate gradient method}

\begin{proposition}[CG convergence]
\label{prop:cg}
Let $A_\Free$ be SPD with condition number $\kappa = \kappa(A_\Free)$, and let $x_k$
be the $k$-th CG iterate for $A_\Free x_\Free = b_\Free$. Then
\[
  \frac{\norm{e_k}_{A_\Free}}{\norm{e_0}_{A_\Free}}
  \;\le\;
  2\left(\frac{\sqrt{\kappa}-1}{\sqrt{\kappa}+1}\right)^{\!k},
\]
where $\norm{e}_{A_\Free} = (e\T A_\Free e)^{1/2}$ and
$e = x_\Free - x_k$.
\end{proposition}

Reaching a fixed relative energy-norm tolerance therefore costs $O(\sqrt\kappa)$
iterations, and clustered spectra do better still---the superlinear phase that a
Chebyshev bound cannot see. Two consequences frame the rest of the chapter.

First, by Proposition~\ref{prop:interlace} the bound applies \emph{uniformly over
every free set the loop might visit}, since $\kappa(A_\Free) \le \kappa(A)$. One
does not need to know which free sets those are.

Second, $\sqrt\kappa$ is the whole reason to keep a Krylov method rather than a
simpler iteration. Section~\ref{sec:projgrad} makes the comparison precise.

\section{The matrix-free operator}
\label{sec:operator-two}

CG needs only $v \mapsto A_\Free v$. When $A = M\T M$ this is
\[
  A_\Free v = M_\Free\T(M_\Free\,v), \qquad M_\Free = M_{:,\Free},
\]
two products at $O(m\,|\Free|)$ with \emph{no} $O(n^2)$ storage---the standard
CG-on-the-normal-equations arrangement \cite{golub2013,paige1975}, on a column
subset that changes between outer steps.
\index{matrix-free}

The whole method needs exactly two operations of $A$, and it is worth naming them
because Chapter~\ref{ch:operator} builds an abstraction on precisely this list.

\begin{enumerate}
\item \textbf{The free-block application} $v \mapsto A_\Free v$, which is all CG
      needs to solve the inner system.
\item \textbf{The cross product}
      $x_\Free \mapsto A_{\Bound,\Free}\,x_\Free$, which supplies the reduced
      gradient $s_i = (Ax)_i - b_i$ at the bound indices (where $x_\Bound = 0$),
      and hence the dual-feasibility test.
\end{enumerate}

Any implementation realising these two runs the loop, its termination proof, and its
convergence analysis unchanged. A dense $A$ slices them directly. The Gram operator
$A = M\T M$ evaluates them from $M$ without forming $A$. A diagonal-plus-low-rank
operator $A = \diag(d) + U\Delta U\T$ inverts the free block by the Woodbury
identity---a \emph{direct} inner solve that bypasses CG entirely, which is what
Chapter~\ref{ch:operator} exploits.

\section{Adding a linear equality}

Many applications carry a linear balance system alongside the bounds:
\begin{equation}
  \min_{x}\; \tfrac12 x\T A x - b\T x
  \quad\text{subject to}\quad Bx = c,\; x \ge 0,
  \qquad B \in \R^{p\times n} \text{ of full row rank}.
  \label{eq:eqp}
\end{equation}
The budget constraint $\onevec\T x = \beta$ is the case $p = 1$; group budgets,
factor neutralities, and netting rules are $p > 1$ at the same cost structure.

Naively this makes the free-set system the indefinite saddle system
of~\eqref{eq:saddle}, which CG cannot touch. The remedy is the Schur elimination of
Chapter~\ref{ch:freeblock}, and it is worth writing out because it is what keeps
every solve SPD.

Stationarity gives $x_\Free = A_\Free^{-1}(b_\Free + B_\Free\T\lambda) = v_0 + V_1\lambda$
in terms of the SPD solves
\[
  v_0 = A_\Free^{-1} b_\Free,
  \qquad
  V_1 = A_\Free^{-1} B_\Free\T \quad(\text{its $p$ columns, one CG solve each}).
\]
Imposing $B_\Free x_\Free = c$ fixes the multiplier through the $p\times p$ system
\[
  S\,\lambda = c - B_\Free v_0,
  \qquad
  S = B_\Free V_1 = B_\Free A_\Free^{-1} B_\Free\T \succ 0,
  \qquad
  x_\Free = v_0 + V_1\lambda .
\]
So each outer step costs $p+1$ matrix-free CG solves sharing one operator, plus a
$p\times p$ dense Cholesky of negligible cost. \emph{The multipliers are recovered
analytically and never appear in the linear solve}, so every inner system stays SPD
and Proposition~\ref{prop:cg} applies to each.

For $p = 1$, $B = \onevec\T$, $c = \beta$ this collapses to two solves
$v_0, v_1 = A_\Free^{-1}\onevec$ with the scalar
$\lambda = (\beta - \onevec\T v_0)/(\onevec\T v_1)$; and the pure-normalisation case
$b = 0$, $\beta = 1$ collapses further to the single solve
$x_\Free = v_1/(\onevec\T v_1)$. That last formula is the entire long-only
minimum-variance portfolio, and Chapter~\ref{ch:operator} returns to it.

\begin{remark}[What this buys over feasible-point methods]
\label{rem:feasible}
There is a well-developed family of \emph{feasible-point} bound-constrained
solvers---Bertsekas' projected Newton method \cite{bertsekas1982}, Moré and
Toraldo's GPCG \cite{moretoraldo1991}, Dostál's proportioning and MPRGP algorithms
\cite{dostal1997,dostalschoberl2005}---in which every iterate satisfies the bounds
and the working face changes through gradient projections. They are excellent
methods and their rates are stated in $\kappa$ as well.

The distinction that matters here is the equality $Bx = c$. Those methods rely on a
closed-form projection onto the feasible set, which the orthant supplies but
$\{x\ge0 : Bx=c\}$ does not; they reach a general linear equality only by wrapping
an outer augmented-Lagrangian loop \cite{dostal2006smalbe}, at the price of an extra
level of iteration and of the exact certificate. The loop of this chapter takes
$Bx = c$ directly through the Schur elimination above.

The trade runs the other way too, and honesty requires stating it: this loop treats
only the non-negative orthant $x \ge 0$, not the general two-sided box
$\ell \le x \le u$ that those methods support. A two-state complementary basis
cannot express two-sided bounds; that would need a mixed-complementarity
formulation.
\end{remark}

\section{The projected-gradient reference}
\label{sec:projgrad}

To see what the Krylov subspace buys, compare against a method that forgoes it.
Projected gradient iterates
\[
  x_{k+1} = \Pi_\Omega\!\left(x_k - \tau\,\nabla f(x_k)\right),
  \qquad \nabla f(x) = A x - b,
\]
with $\tau = 1/L$, $L = \lambda_{\max}(A)$. For the orthant
$\Pi_\Omega(y) = \max(y,0)$ componentwise; for the simplex slice
$\{x \ge 0 : \onevec\T x = \beta\}$ the projection is $O(n\log n)$ by a single
sort-and-threshold pass \cite{duchi2008,condat2016}. There is no outer loop and no
linear solve.
\index{projected gradient}

\begin{proposition}[Projected-gradient convergence]
\label{prop:proj}
Let $A \succ 0$, so $f$ is $\mu$-strongly convex and $L$-smooth with
$\mu = \lambda_{\min}(A)$, $L = \lambda_{\max}(A)$. Projected gradient with
$\tau = 1/L$ satisfies
$\norm{x_k - x^\star}^2 \le (1 - 1/\kappa(A))^{k}\norm{x_0 - x^\star}^2$.
\end{proposition}

\begin{proof}[Sketch]
$x^\star$ is a fixed point; non-expansiveness of $\Pi_\Omega$ and co-coercivity of
the $L$-smooth gradient give
$\norm{x_{k+1}-x^\star}^2 \le \norm{x_k-x^\star}^2 -
\tau(2-\tau L)\inner{\nabla f(x_k)-\nabla f(x^\star)}{x_k-x^\star}$; setting
$\tau = 1/L$ and applying strong convexity gives the claim.
\end{proof}

So the iteration count is $O(\kappa)$ against $O(\sqrt\kappa)$ for CG, at the same
per-step cost of one operator application. On ill-conditioned problems that
$\sqrt\kappa$ gap is a runtime gap. Nesterov acceleration \cite{nesterov2004,beck2009}
would lift the projected-gradient rate to $O(\sqrt\kappa)$, but CG attains that rate
\emph{optimally} within the Krylov subspace: each CG iterate minimises $f$ over
$\Krylov_k$, a strictly richer search space than two-point momentum, and typically
with a smaller constant.

The trade is otherwise in the other direction, and the honest summary is that
projected gradient is loop-free, needs no complementarity certificate, and is the
simpler method when $\kappa$ is small or moderate accuracy suffices. It also loses
its per-step simplicity the moment $p > 1$: projecting onto
$\{x\ge0 : Bx=c\}$ for general $B$ is itself a non-negative quadratic program.

\section{Inexactness: composing the two guarantees}
\label{sec:inexact}

We now close the gap flagged at the end of Chapter~\ref{ch:guessing}.

\begin{lemma}[Inexact inner solves]
\label{lem:inexact}
\index{decision margin}
Fix the bound-constrained problem~\eqref{eq:bqp}. For a free set $\Free$ let
$\hat x(\Free)$ be the exact solution of the inner system extended by zeros, and
$\hat s(\Free) = A\hat x(\Free) - b$. Let $\Free_0 = \{1,\dots,n\}, \dots, \Free_T$
be the free-set trajectory of Algorithm~\ref{alg:elim} run with \emph{exact} inner
solves, finite by Theorem~\ref{thm:termination}. Define the \emph{decision margin}
\[
  \mu = \min_{0 \le t \le T} \min\Bigl(
    \{ |\hat x_i(\Free_t)| : i \in \Free_t,\ \hat x_i \ne 0 \}
    \cup
    \{ |\hat s_i(\Free_t)| : i \notin \Free_t,\ \hat s_i \ne 0 \}
  \Bigr) > 0,
\]
the smallest nonzero magnitude any primal or dual test encounters along the exact
trajectory. Suppose $0 < \varepsilon < \mu$ and every inner CG call is stopped at a
residual $\rho = \norm{b_\Free - A_\Free\tilde x_\Free}_2$ obeying
\begin{equation}
  \rho\left(\frac{1 + \lambda_{\max}(A)}{\lambda_{\min}(A)}\right)
  \;<\; \min\bigl(\varepsilon,\ \mu - \varepsilon\bigr).
  \label{eq:inexactcond}
\end{equation}
Then at every outer step the violator sets computed from the CG iterate coincide
with those computed from the exact solve on the same free set. The inexact loop
therefore traverses exactly the trajectory $\Free_0,\dots,\Free_T$,
Theorem~\ref{thm:termination} applies verbatim, and
$\norm{\tilde x - x^\star}_\infty \le \rho/\lambda_{\min}(A)$.
\end{lemma}

\begin{proof}
Write $r = b_\Free - A_\Free\tilde x_\Free$, so
$\tilde x_\Free - \hat x_\Free = -A_\Free^{-1}r$. By Cauchy interlacing
(Proposition~\ref{prop:interlace}), $\lambda_{\min}(A_\Free) \ge \lambda_{\min}(A)$,
hence
$\delta_x := \norm{\tilde x_\Free - \hat x_\Free}_2 \le \rho/\lambda_{\min}(A)$,
which also bounds each coordinate. For $i \notin \Free$,
$\tilde s_i - \hat s_i = A_{i,\Free}(\tilde x_\Free - \hat x_\Free)$, so by
Cauchy--Schwarz $\delta_s \le \lambda_{\max}(A)\delta_x$. Both are bounded by
$\delta := \rho(1+\lambda_{\max}(A))/\lambda_{\min}(A) < \min(\varepsilon,\mu-\varepsilon)$.

Fix $i \in \Free_t$. If $\hat x_i \ge 0$ then $\hat x_i = 0$ or $\hat x_i \ge \mu$;
either way $\tilde x_i \ge -\delta > -\varepsilon$ and the inexact test does not flag
$i$, matching the exact test. If $\hat x_i < 0$ then $\hat x_i \le -\mu$, so
$\tilde x_i \le -\mu + \delta < -\varepsilon$ and both tests flag $i$. The same
argument with $\hat s_i, \tilde s_i, \delta_s$ covers the dual test at every bound
index. Hence the violator sets agree at step $t$; since both loops start from the
same $\Free_0$ and the counters are deterministic functions of the violator
sequence, induction on $t$ gives the same trajectory and the same stopping step. At
termination the exact iterate is $x^\star$, so
$\norm{\tilde x - x^\star}_2 = \delta_x \le \rho/\lambda_{\min}(A)$.
\end{proof}

Read the structure of that argument, because it generalises. One does not prove a
new termination theorem for the approximate method. One proves that the approximate
method makes the \emph{same decisions} as the exact method, and then inherits the
exact method's theorem unchanged. Whenever an algorithm's control flow depends on
finitely many sign tests, and those tests have a margin, this is the argument
available.

The corresponding statement for the equality-augmented problem~\eqref{eq:eqp} needs
one extra step, because the multiplier $\lambda$ is itself recovered from $p+1$
inexact solves through the Schur complement rather than read off a single residual.
Propagating the perturbation through $\tilde S^{-1}$ by a Neumann series gives a
bound of the same shape, and the sign-test argument then runs verbatim.

\begin{remark}[The guarantee attaches to the outer loop]
Nothing in Theorem~\ref{thm:termination} mentions CG. A \emph{direct} factorisation
may replace the Krylov inner solve unchanged, and then no inexactness condition
needs verifying at all. Chapter~\ref{ch:operator} does exactly this with a Woodbury
solve. The separation is worth naming: the outer loop owns the combinatorics and the
certificate; the inner solver owns the linear algebra; and Lemma~\ref{lem:inexact} is
the interface contract between them.
\end{remark}

\section{Warm starts}

For a parametric family---a sequence of right-hand sides $b(\theta_1), b(\theta_2),
\dots$, or a swept normalisation $\beta$---consecutive minimisers usually differ in
only a few active constraints. Passing the previous problem's $(\Free, x)$ as the
initial state reduces both counts: the outer count, because the free set needs fewer
adjustments, and the inner count, because CG starts from a small initial residual.

\begin{proposition}[Warm starts across a support-stable step]
\label{prop:warm}
Let $b' = b + \delta$ and suppose the free set $\Free$ returned for $b$ is also
optimal for $b'$. Then Algorithm~\ref{alg:elim} started from $\Free$ terminates in a
\emph{single} outer step, and its one CG solve, warm-started at the previous
solution, reaches energy-norm accuracy $\tau$ in at most
\[
  k \;\le\; \left\lceil \frac{\sqrt{\kappa}+1}{2}\,
  \ln\!\frac{2\,\norm{\delta}_2}{\tau\sqrt{\lambda_{\min}(A)}} \right\rceil
  \text{ iterations.}
\]
\end{proposition}

\begin{proof}
The single solve returns $\hat x_\Free \ge 0$ with $s \ge 0$ on the bound set, so
both violator sets are empty and the loop exits, certified by Step~2 of
Theorem~\ref{thm:termination}'s proof. The warm start's initial error is
$e_0 = A_\Free^{-1}\delta_\Free$, whose energy norm is
$(\delta_\Free\T A_\Free^{-1}\delta_\Free)^{1/2} \le \norm{\delta}_2/\sqrt{\lambda_{\min}(A)}$
by interlacing. Proposition~\ref{prop:cg} with
$\rho_{\mathrm{cg}} = (\sqrt\kappa-1)/(\sqrt\kappa+1)$ and
$\ln(1/\rho_{\mathrm{cg}}) \ge 2/(\sqrt\kappa+1)$ gives the bound.
\end{proof}

The interesting feature is that the inner count scales with the logarithm of the
\emph{parameter step} $\norm{\delta}$, not of the solution norm: a finer grid is
cheaper per point. When the support drifts, no comparable a-priori bound is
available---the loop is not monotone in the free set---but empirically the
warm-started outer count tracks the support drift $|\Free \triangle \Free'|$ and
collapses to one whenever the drift is zero, exactly as the proposition prescribes.

Note also which solver profits from what. A \emph{direct} inner solver profits only
from the warm free set, having no analogue of an initial guess. A matrix-free CG
inner solver profits from both. That is where warm-started parametric sweeps obtain
their largest speed-ups, and Chapter~\ref{ch:warmstart} develops the theme.

\begin{notes}
The construction of this chapter is \cite{schmelzer2026nncg}, released as the
open-source \texttt{nncg} package; Lemma~\ref{lem:inexact} and its
equality-augmented corollary are its contribution, the exact-solve termination count
being Júdice and Pires' \cite{judice1994}. CG is \cite{hestenes1952};
\cite{greenbaum1997} and \cite{trefethenbau1997} are the textbook treatments of its
convergence, and \cite{golub2013} of the normal-equations arrangement. On the NNLS
side, Bro and De Jong's fast NNLS \cite{brodejong1997} accelerates Lawson--Hanson
\cite{lawson1974} by caching normal-equations cross-products, but still assembles
$M\T M$; the method here forms neither.
\end{notes}

\begin{exercises}

\begin{exercise}
Show that $A = M\T M$ is positive definite iff $M$ has full column rank, and that
the free-block operator $v \mapsto M_\Free\T(M_\Free v)$ is SPD for every $\Free$
when it is.
\end{exercise}

\begin{exercise}
Derive the $p = 1$ specialisation of the Schur elimination, and verify that
$b = 0$, $\beta = 1$ gives $x_\Free = v_1/(\onevec\T v_1)$.
\end{exercise}

\begin{exercise}
Show $S = B_\Free A_\Free^{-1} B_\Free\T \succ 0$ whenever $A_\Free \succ 0$ and
$B_\Free$ has full row rank. For $p = 1$, when can $B_\Free$ fail to have full row
rank?
\end{exercise}

\begin{exercise}
Prove that CG applied to $A_\Free x = b_\Free$ from $x_0 = 0$ produces iterates in
$\Krylov_k(A_\Free, b_\Free)$, and deduce that the exact solution is reached in at
most $|\Free|$ iterations in exact arithmetic. Why is this fact of limited practical
use?
\end{exercise}

\begin{exercise}
The decision margin $\mu$ of Lemma~\ref{lem:inexact} is defined along the exact
trajectory, which one does not know in advance. Explain why the lemma is
nevertheless useful, and propose a practical stopping rule that respects its spirit.
\end{exercise}

\begin{exercise}
Compare the total work of CG and projected gradient to reach relative accuracy
$10^{-8}$ on a problem with $\kappa = 10^{4}$, assuming one operator application per
step in both. Now repeat for $\kappa = 10^2$ and comment.
\end{exercise}

\begin{exercise}[Implementation]
Implement Algorithm~\ref{alg:elim} with a matrix-free CG inner solver on
$A = M\T M$. On a synthetic NNLS problem with a planted non-negative solution, plot
total CG iterations against the inner residual tolerance and locate the point at
which the outer trajectory changes---i.e.\ where~\eqref{eq:inexactcond} fails.
\end{exercise}

\begin{exercise}[Implementation]
Sweep the normalisation $\beta$ over a grid of $50$ values with and without warm
starting. Plot outer steps and total inner iterations for both, and check
Proposition~\ref{prop:warm}'s prediction that the inner count scales with
$\log\norm{\delta}$ by refining the grid.
\end{exercise}

\end{exercises}


\part{The Whole Path}
\chapter{Parametric Quadratic Programming}
\label{ch:parametric}

So far a QP has been a single problem with a single answer. This chapter changes
the question. Attach a scalar parameter to the linear term and ask how the solution
moves as the parameter varies:
\begin{equation}
\label{eq:pqp}
\begin{aligned}
x(\lambda) = \argmin_{x \in \R^n}\ &\tfrac12\, x\trans H x - \lambda\, d\trans x \\
\text{s.t.}\quad &A x = b,\ \ G x \le h,\ \ \ell \le x \le u .
\end{aligned}
\end{equation}
We have switched to the symbols the parametric literature uses---$H$ for the
quadratic form, $A$ and $G$ for equality and inequality systems, $\ell,u$ for
bounds---because Part~III is where this book meets four other literatures and
Appendix~\ref{app:notation} has to translate somewhere.

The answer is remarkable and is the foundation of everything in Part~III:
\emph{the solution map is piecewise linear}. On any interval of $\lambda$ where the
active set is constant, $x(\lambda)$ is an affine function of $\lambda$; the active
set changes at finitely many \emph{breakpoints}; and the entire family can therefore
be computed \emph{exactly}, corner to corner, by one finite algorithm. Not
approximated on a grid. Computed.

\section{The standing assumption}

We require only that $H \succeq 0$ and that the \emph{reduced} Hessian
$H_{\Free\Free}$ on the free block is positive definite at every active set the path
visits. This is weaker than the $H \succ 0$ used in Parts~I and~II, and deliberately
so: it is exactly what the regression instance of Chapter~\ref{ch:homotopy} needs.
There $H = X\T X$ is singular whenever the design has more columns than rows---the
very regime the LASSO exists for---yet the free block $X_\Free\T X_\Free$ is positive
definite as soon as the active columns are linearly independent, which holds in
general position. The finance instance has $H = \Sig \succ 0$ outright, so the
condition is automatic there.

\begin{remark}[Where the assumption bites]
\label{rem:rank}
This is an assumption and not a formality, and it is worth saying where it fails
rather than burying it. When $H_{\Free\Free}$ is singular at some visited active set,
the bordered system~\eqref{eq:border} below no longer determines the segment, and
what the path even \emph{is} becomes a question of which minimiser one selects.
Existing treatments divide on the answer. Tibshirani and Taylor \cite{tibshirani2011}
for the generalized lasso, and Gaines, Kim and Zhou \cite{gaines2018} for the
constrained LASSO, both restore strong convexity by adding a small ridge, so the
algorithms they analyse are those of a perturbed problem. Tibshirani
\cite{tibshirani2013}, by contrast, gives a LASSO path algorithm valid for
\emph{arbitrary} designs with no general-position hypothesis: it succeeds by tracing
the minimum-$\ell_2$-norm solution specifically, and turns on an
insertion--deletion property of that particular selection rather than on the segment
structure alone. That is a genuinely different argument.

No counterpart is known for the generalized or constrained problem. Whether the
bordered system admits a canonical selection playing the same role is, as far as we
know, open, and Chapter~\ref{ch:reading} returns to it as one of the subject's
genuine unsolved questions.
\end{remark}

\section{Affine segments}

\begin{lemma}[Affine segments and the bordered solve]
\label{lem:affine}
\index{affine segment}
On any maximal interval of $\lambda$ on which the active set is constant, the
optimiser of \eqref{eq:pqp} is affine,
\begin{equation}
\label{eq:affine}
x(\lambda) = \alpha + \lambda\,\beta ,
\end{equation}
and $(\alpha,\beta)$ come from a single bordered linear solve~\eqref{eq:border}
whose matrix is constant on the interval. An active inequality row enters that solve
identically to an equality row.
\end{lemma}

\begin{proof}
Fix $\lambda$ in the interior of the interval and name the active set: the
equalities $A$ (always active); the active inequality rows
$\mathcal{S} \subseteq \{1,\dots,p\}$, those held at $g_i\T x = h_i$, the remainder
carrying a zero multiplier by complementarity; and the blocked variables $\Bound$
pinned at a bound, with free set $\Free = \Bound^c$. Stack the active linear
constraints as $C = [A; G_{\mathcal{S}}]$ with right-hand side
$r = (b; h_{\mathcal{S}})$.

Because \eqref{eq:pqp} is convex with linear constraints, the KKT conditions are
necessary and sufficient. Stationarity is $Hx - \lambda d + C\T\zeta + \rho = 0$,
with $\zeta$ the equality and active-inequality multipliers and $\rho$ the bound
multipliers, supported on $\Bound$. Read the $\Free$ rows, where $\rho_\Free = 0$,
and split $(Hx)_\Free = H_{\Free\Free}x_\Free + H_{\Free\Bound}x_\Bound$ with
$x_\Bound$ fixed at its bound values. Primal feasibility of the active rows is
$C_\Free x_\Free + C_\Bound x_\Bound = r$. Collecting the unknowns
$(x_\Free,\zeta)$,
\begin{equation}
\label{eq:border}
\begin{bmatrix} H_{\Free\Free} & C_\Free\T \\[2pt] C_\Free & 0 \end{bmatrix}
\begin{bmatrix} x_\Free \\[2pt] \zeta \end{bmatrix}
=
\begin{bmatrix} \lambda\, d_\Free - H_{\Free\Bound}\,x_\Bound \\[2pt]
                r - C_\Bound\,x_\Bound \end{bmatrix}.
\end{equation}
The coefficient matrix is nonsingular by exactly the argument of
Proposition~\ref{prop:saddlenonsing}, using $H_{\Free\Free}\succ0$ and full row rank
of $C_\Free$. It depends only on $H$ and the active normals, hence is constant on
the interval, while the right-hand side is affine in $\lambda$---the only $\lambda$
term being $\lambda d_\Free$. So $(x_\Free(\lambda),\zeta(\lambda))$ is affine, and
appending the $\lambda$-independent blocked entries gives \eqref{eq:affine} with
$\beta_\Bound = 0$.

Finally, $G_{\mathcal{S}}$ enters~\eqref{eq:border} only as extra rows of $C$,
indistinguishable from the equality rows, and the box bounds are carried by the
partition $(\Free,\Bound)$ rather than folded into $C$. A polyhedral penalty
$\lambda\norm{x}_1$ is the same statement with $d$ replaced, on each segment, by the
constant subgradient sign pattern of the support---the only change needed to
read~\eqref{eq:border} as the LASSO segment solve.
\end{proof}

Each homotopy step is therefore \emph{one solve} of~\eqref{eq:border}: a
factorisation of $H_{\Free\Free}$ bordered by the active normals through its Schur
complement, which is exactly the block elimination~\eqref{eq:blockelim} of
Chapter~\ref{ch:freeblock}. In practice one solves it for two right-hand sides at
once against a single factorisation, obtaining $\alpha$ and $\beta$ together.

\section{How the active set evolves}
\label{sec:events}

A representation is not an algorithm. Lemma~\ref{lem:affine} says what a segment
\emph{is}; it does not say how one segment gives way to the next. That second half
is where each field's paper does its technical work, and where the instances of
Chapter~\ref{ch:homotopy} are least recognisable to one another. We write it once.

Fix the current breakpoint $\lambda_k$, the current active set, and a direction of
travel $\sigma \in \{+1,-1\}$; write $\lambda = \lambda_k + \sigma t$ and call
$t \ge 0$ the step. Solving \eqref{eq:border} twice against one factorisation gives
the segment $x_\Free(\lambda) = \alpha_\Free + \lambda\beta_\Free$ and the
multipliers $\zeta(\lambda) = \zeta^0 + \lambda\zeta^1$; the bound multipliers follow
from stationarity as $\rho(\lambda) = \lambda d - Hx(\lambda) - C\T\zeta(\lambda)$,
affine in $\lambda$ like everything else. Write $\dot\phi = \sigma\phi_1$ for the
rate of change of any affine $\phi(\lambda) = \phi_0 + \lambda\phi_1$ along the
direction of travel.

\emph{Four things, and only four, can end a segment}: a primal quantity can reach a
bound it must respect, or a dual quantity can reach the edge of the cone it must lie
in.
\index{events!the four families}

\paragraph{$\mathrm{P}_1$: a free variable reaches a bound.} For $i \in \Free$ the
coordinate must stay in $[\ell_i,u_i]$, so
\[
t^{\mathrm{P}_1}_i =
\begin{cases}
(u_i - x_i)/\dot x_i, & \dot x_i>0,\\
(\ell_i - x_i)/\dot x_i, & \dot x_i<0,\\
+\infty, & \dot x_i=0,
\end{cases}
\]
and the event moves $i$ from $\Free$ to $\Bound$, pinned at the bound it met.

\paragraph{$\mathrm{P}_2$: an inactive inequality becomes active.} For
$j \notin \mathcal{S}$ the slack $h_j - g_j\T x$ must stay non-negative, giving
$t^{\mathrm{P}_2}_j = (h_j - g_j\T x)/(g_j\T\dot x)$ when $g_j\T\dot x>0$ and
$+\infty$ otherwise; the event adjoins $j$ to $\mathcal{S}$, adding a row to $C$.

\paragraph{$\mathrm{D}_1$: a blocked variable releases.} For $i \in \Bound$ dual
feasibility confines $\rho_i(\lambda)$ to an interval
$[\underline\rho_i(\lambda), \overline\rho_i(\lambda)]$ whose endpoints are
themselves affine in $\lambda$. Both margins $e^-_i = \rho_i - \underline\rho_i$ and
$e^+_i = \overline\rho_i - \rho_i$ are affine and non-negative at $\lambda_k$, and
the step is the first zero of either,
\[
t^{\mathrm{D}_1}_i = \min\bigl\{ -e^-_i/\dot e^-_i \ \text{if}\ \dot e^-_i<0,\ \
-e^+_i/\dot e^+_i \ \text{if}\ \dot e^+_i<0 \bigr\},
\]
and the event moves $i$ from $\Bound$ to $\Free$.

This is the family worth pausing on. For a genuine box the interval is
$[0,\infty)$ at an upper bound and $(-\infty,0]$ at a lower one, and the event is
simply $\rho_i = 0$. For an $\ell_1$ penalty it is $[-\lambda,\lambda]$, and the
event is $|\rho_i| = \lambda$: an affine function meeting a \emph{moving} threshold
rather than a fixed one. \emph{That single generalisation---a feasibility boundary
that travels with $\lambda$---is the whole of what separates an ``induced'' active
set from an ``imposed'' one at the level of the algorithm.}

\paragraph{$\mathrm{D}_2$: an active inequality releases.} For $j \in \mathcal{S}$
the multiplier must satisfy $\zeta_j \ge 0$, so
$t^{\mathrm{D}_2}_j = -\zeta_j/\dot\zeta_j$ when $\dot\zeta_j<0$ and $+\infty$
otherwise; the event deletes $j$ from $\mathcal{S}$. Equality rows never generate
this event, their multipliers being unsigned.

\medskip
The step is $t^\star = \min t$ over the four families, the event attaining it is
applied, and the walk continues. This is the simplex method's ratio test, and it
inherits the simplex method's degeneracy hazard---ties at $t^\star = 0$ and
cycling---together with its remedy: a lowest-index (Bland) tie-break, which makes
the walk deterministic and finite. Because exactly one index changes status per
step, the factorisation of the bordered matrix is \emph{updated} rather than
rebuilt, by Section~\ref{sec:updating} of Chapter~\ref{ch:freeblock}.

\begin{table}[ht]
\centering
\footnotesize
\caption{The four events as two literatures name them. The algorithms differ in
which families are populated, not in the families.}
\label{tab:events}
\begin{tabular}{@{}lll@{}}
\toprule
Event & LASSO\,/\,LARS & Critical Line Algorithm \\
\midrule
$\mathrm{P}_1$ & a coefficient crosses zero (the ``leave'' move)
               & an asset reaches a bound and leaves the free set \\
$\mathrm{P}_2$ & a linear inequality on $\beta$ becomes tight
               & a group or turnover constraint becomes tight \\
$\mathrm{D}_1$ & an inactive coordinate's correlation reaches $\lambda$ (``enter'')
               & an asset's reduced cost changes sign and it re-enters \\
$\mathrm{D}_2$ & an active inequality's multiplier hits zero
               & a constraint stops binding at a turning point \\
\bottomrule
\end{tabular}
\end{table}

Table~\ref{tab:events} reads the catalogue back into two dialects. Forward least
angle regression has $\mathrm{D}_1$ only, which is exactly what makes its active set
monotone; the LASSO modification restores $\mathrm{P}_1$ and with it the leave move;
the Critical Line Algorithm runs $\mathrm{P}_1$ and $\mathrm{D}_1$ against genuine
box bounds, with fixed rather than moving feasibility intervals; the constrained
LASSO adds $\mathrm{P}_2$ and $\mathrm{D}_2$. Chapter~\ref{ch:homotopy} completes
the table.

\section{The homotopy}

Continuity of $x(\lambda)$ across breakpoints---the free block is uniquely
determined on each segment, and adjacent segments agree at the breakpoint---lets a
\emph{homotopy} recover the whole path: from a vertex with known active set, follow
\eqref{eq:affine} to the first event, update the active set by that one event,
repeat.
\index{homotopy}

\begin{algorithmbox}[The parametric active-set homotopy]
\label{alg:homotopy}
\begin{enumerate}
\item Compute a starting vertex and its active set (Section~\ref{sec:first} below).
\item While $\lambda$ has not reached its terminal value:
  \begin{enumerate}
  \item solve \eqref{eq:border} for two right-hand sides against one factorisation,
        obtaining $\alpha$, $\beta$, and the affine multipliers;
  \item compute every candidate step $t$ in the four families of
        Section~\ref{sec:events}, discarding inadmissible ones;
  \item take $t^\star = \min t$; if none is admissible, stop;
  \item apply the single event attaining $t^\star$ (Bland tie-break), update
        $(\Free,\Bound,\mathcal{S})$ and the factorisation, and record the
        breakpoint $x(\lambda_k + \sigma t^\star)$.
  \end{enumerate}
\end{enumerate}
\end{algorithmbox}

\subsection{The first turning point}
\label{sec:first}

The walk needs a starting vertex. For the portfolio instance, at $\lambda = +\infty$
variance is irrelevant and~\eqref{eq:pqp} degenerates into ``earn the most expected
return the bounds and budget allow''. With a single all-ones budget row that vertex
is computed greedily: initialise every weight at its lower bound, sort assets by
descending expected return, and walk down that order raising each weight toward its
upper bound until the budget is met. The last asset touched is only partially filled
and is the single \emph{free} asset of the first turning point.

Two details of that construction repay attention, because both are the kind of thing
a derivation omits and an implementation cannot.

First, the budget test needs a tolerance. The accumulated increment can only reach
$1$ up to rounding, and without slack the loop skips past the genuinely interior
asset and marks the next, zero-weight, asset as free.

Second, marking the last asset \emph{free} even when it lands exactly on a bound is
deliberate: it guarantees the first turning point has at least one free asset, so the
reduced system is nonsingular from the outset. Otherwise every asset sits on a bound
and the reduced solve degenerates.

For a general equality system the greedy fill no longer applies, and the
maximum-return vertex is instead the solution of a linear program,
$\max_x d\T x$ subject to $Ax = b$, $\ell \le x \le u$, which a simplex code
returns as a vertex. Only the first turning point needs the linear program; the loop
that follows is unchanged.

\subsection{Termination}

\begin{proposition}[Finite termination of the homotopy]
\label{prop:homotopy-finite}
Under the standing assumption, with a lowest-index tie-break, the homotopy visits
each active-set state at most once and terminates in finitely many steps.
\end{proposition}

\begin{proof}
There are finitely many active-set states, at most $2^{n+p}$. The parameter
$\lambda$ is monotone along the walk. The only failure mode is revisiting a state at
the same $t^\star = 0$, and Bland's lowest-index rule excludes that exactly as it
does for the simplex method.
\end{proof}

The bound $2^{n+p}$ is a finiteness statement, not an efficiency estimate. In
practice the number of breakpoints is empirically $O(n)$, and
Chapter~\ref{ch:reading} discusses what is and is not known about the gap.

\section{Two readings of the parameter}

Before Chapter~\ref{ch:homotopy} identifies the instances, one conceptual point.
The parameter $\lambda$ in~\eqref{eq:pqp} multiplies the \emph{linear} term while
the quadratic term carries no coefficient of its own. In the portfolio reading it is
therefore a \emph{return tilt}, not a risk aversion; sweeping it from $0$ (minimum
variance) to $\infty$ (maximum return) traces the whole efficient frontier. In the
regression reading it is a penalty, and the sweep runs the other way, from
$\lambda_{\max}$ (everything zero) to $0$ (the unpenalised fit).

Either way it plays the role a statistician reserves for a regularisation
parameter: a single knob weighing two competing goals, whose whole trajectory---not
any one point of it---carries the information. That the trajectory happens to be
computable exactly, rather than sampled, is what makes the subject worth a part of
this book.

\begin{remark}[Vector parameters]
\label{rem:vector}
Nothing forces $\lambda$ to be a scalar. If the parameter is a \emph{vector}---the
state of a plant, in the control instance---then the interval of breakpoints becomes
a polyhedral partition of parameter space into \emph{critical regions}, and the
homotopy becomes region traversal. That is multi-parametric quadratic programming
\cite{bemporad2002,tondel2003}. It is a genuine generalisation of the scalar device
rather than another instance of it, and this book treats only its scalar slice;
Chapter~\ref{ch:homotopy} says where the boundary falls.
\end{remark}

\begin{notes}
The piecewise-linearity of the solution path under a quadratic loss and a polyhedral
penalty was characterised within statistics by Rosset and Zhu \cite{rosset2007}; the
same fact is the piecewise-affinity of the mpQP solution over critical regions in
control \cite{bemporad2002}, and Roll \cite{roll2008} traces single-parameter paths
in that setting, explicitly generalising the LARS construction. The bordered-solve
presentation here follows \cite{schmelzer2026perspective}, and the first-turning-point
construction and its two details follow the \texttt{cvxcla} software paper
\cite{cvxcla}; the degenerate case in which every starting asset sits on a bound is
not addressed by the otherwise thorough step-by-step review of Markowitz and
Todd \cite{markowitz2021}.
\end{notes}

\begin{exercises}

\begin{exercise}
Verify that the coefficient matrix of~\eqref{eq:border} is nonsingular under the
standing assumption, adapting the proof of Proposition~\ref{prop:saddlenonsing}.
Where exactly is $H_{\Free\Free} \succ 0$ used, and where is full row rank of
$C_\Free$ used?
\end{exercise}

\begin{exercise}
Show that $\beta_\Bound = 0$ on every segment, and interpret this: what does it say
about a blocked variable as $\lambda$ moves?
\end{exercise}

\begin{exercise}
Take $n = 2$, $H = I$, $d = (1,0)\T$, bounds $0 \le x \le 1$, and the single equality
$x_1 + x_2 = 1$. Trace the path from $\lambda = \infty$ to $\lambda = 0$ by hand,
identifying every breakpoint and which of the four event families produced it.
\end{exercise}

\begin{exercise}
Show that for a genuine box bound the $\mathrm{D}_1$ event reduces to $\rho_i = 0$,
and that for the $\ell_1$ penalty it is $|\rho_i| = \lambda$. Why does the second
case make the active set ``induced'' rather than ``imposed''?
\end{exercise}

\begin{exercise}
Prove that between consecutive breakpoints the objective value
$\tfrac12 x(\lambda)\T H x(\lambda)$ is a quadratic function of $\lambda$, and that
the parametric objective $\tfrac12 x\T Hx - \lambda d\T x$ is concave in $\lambda$
along the path.
\end{exercise}

\begin{exercise}
Explain why the first turning point of the greedy construction has at least one free
asset, and construct data for which the naive version (no tolerance on the budget
test) marks the wrong asset free.
\end{exercise}

\begin{exercise}
Show that the number of breakpoints is at least the number of distinct active sets
visited, and give an instance with $n = 3$ where an asset leaves the free set and
later re-enters.
\end{exercise}

\begin{exercise}[Implementation]
Implement Algorithm~\ref{alg:homotopy} for the box-and-budget case (events
$\mathrm{P}_1$ and $\mathrm{D}_1$ only). Validate it against a general-purpose QP
solver called at $50$ values of $\lambda$: the traced segment, evaluated at each
value, should agree to solver precision. Report the largest discrepancy.
\end{exercise}

\begin{exercise}[Implementation]
Extend your tracer to inequality rows ($\mathrm{P}_2$, $\mathrm{D}_2$) and validate a
group-exposure cap the same way. How many extra breakpoints does one cap introduce?
\end{exercise}

\end{exercises}

\chapter{One Homotopy, Many Names}
\label{ch:homotopy}

Chapter~\ref{ch:parametric} developed the parametric active-set homotopy as a piece
of optimization. This chapter makes the claim that gives Part~III its point: the
homotopy is not an implementation detail of any one method but a \emph{computational
primitive}, and at least four literatures arrived at it independently. In one case
two of the instances are not analogous but \emph{identical}, and
Theorem~\ref{thm:identity} proves it by a one-line substitution.

\section{The instances}

Read through Chapter~\ref{ch:parametric}, each member is the common scheme with its
quadratic form and its face-mechanism filled in.

\paragraph{The LASSO and least angle regression.} $H = X\T X$, $d = X\T y$; the
active set is \emph{induced} by the $\ell_1$ subgradient and is non-monotone. A
coordinate enters when its correlation with the residual reaches $\lambda$ (event
$\mathrm{D}_1$ with a moving threshold) and leaves when its coefficient crosses zero
(event $\mathrm{P}_1$). Forward least angle regression is the entering-only
ancestor; the LASSO modification restores the leave move
\cite{efron2004,osborne2000,tibshirani1996}.

\paragraph{The Critical Line Algorithm.} $H = \Sig$, $d = \mu$; the active set is
\emph{imposed} by the box bounds and is non-monotone, with assets leaving and
re-entering the free set as $\lambda$ decreases, tracing the \emph{corner
portfolios} of the efficient frontier. It is the oldest instance by four decades
\cite{markowitz1956,markowitz1959}.
\index{Critical Line Algorithm}

\paragraph{The support-vector-machine path.} Hastie et al.\ \cite{hastie2004svm}
trace the SVM solution as the cost parameter varies; $H$ is a \emph{kernel} Gram
matrix, the active set is the set of margin support vectors, and the events are
points crossing the margin. Induced and non-monotone, and the sharpest illustration
of the operator reading of Chapter~\ref{ch:operator}: the path touches the kernel
only through products, never forming it densely.

\paragraph{The elastic net and the generalized lasso.} The elastic net
\cite{zou2005} adds a ridge, sending $H \to X\T X + \eta I$ and otherwise leaving the
scheme untouched: the regularisation merely conditions the quadratic form. The
generalized lasso \cite{tibshirani2011}, with penalty $\lambda\norm{D\beta}_1$ for a
structured difference operator---the fused lasso, trend filtering, and graph
penalties as special cases---traces a piecewise-linear \emph{dual} path by the same
homotopy; the face mechanism acts on $D\beta$ rather than $\beta$, a change of
coordinates, not of scheme.

\paragraph{Explicit model predictive control.} The constrained linear-quadratic
regulator solved as a function of the state \cite{bemporad2002,tondel2003} is the
multidimensional cousin of Remark~\ref{rem:vector}: the parameter is a
\emph{vector}, the interval of breakpoints becomes a polyhedral partition into
critical regions, and the homotopy becomes region traversal. Only its one-parameter
slice is literally \eqref{eq:pqp}, so explicit MPC is a \emph{generalisation} of the
common scheme rather than another instance of the scalar identity below.

\medskip
None of these needs machinery beyond Chapter~\ref{ch:parametric}. What looks like
five literatures is one algorithm with five fillings.

\section{The efficient frontier, for the reader who has not met it}
\label{sec:frontier-primer}

Since the identity below pairs finance with regression, a paragraph of each for the
reader who knows only one.

Mean--variance portfolio selection is a bi-objective problem, and the statistician
already knows its shape. An investor spreads weights $w$ over $n$ assets, the budget
$\onevec\T w = 1$ fixing the scale and a long-only mandate adding $w \ge 0$, and
wants two incompatible things: high expected return $\mu\T w$ and low risk
$w\T\Sig w$. No portfolio optimises both, so the object of study is the
\emph{efficient frontier}: the portfolios that are not dominated. That is precisely a
\emph{Pareto front}, and ``Markowitz-efficient'' is Pareto's nondomination in another
vocabulary.
\index{efficient frontier}

A convex Pareto front is recovered by scalarisation: weight the two objectives with a
single $\lambda \ge 0$ and minimise $\tfrac12 w\T\Sig w - \lambda\mu\T w$, which is
\eqref{eq:pqp}. Convex risk and polyhedral constraints make the sweep \emph{exact},
with nothing skipped: the front is an exact parabola in the
(variance, return) plane between corner portfolios, and those corners---where the set
of held assets changes---are its breakpoints.

So the efficient frontier is a Pareto front of the same \emph{kind} as a
regularisation path: a convex bi-objective traded off by one weight and traced corner
to corner. Which front, and traded against what, differs by instance: return against
variance here, fit against sparsity for the LASSO.

\section{The exact identity}
\label{sec:identity}

Now the sharpest special case: two of the most dissimilar members are, under one
substitution, the same instance.

For weights $w \in \R^n$, expected returns $\mu$, and covariance $\Sig \succ 0$, the
gross-exposure-constrained mean--variance program is
\begin{equation}
\label{eq:mark}
\begin{aligned}
\mathrm{M}_c:\quad \min_{w}\ &\tfrac12\, w\trans \Sig\, w - \mu\trans w \\
\text{s.t.}\quad &\norm{w}_1 \le c,\ \ A w = b,\ \ G w \le h .
\end{aligned}
\end{equation}
For a response $y \in \R^m$ and design $X \in \R^{m\times n}$, the constrained LASSO
is
\begin{equation}
\label{eq:lasso}
\begin{aligned}
\mathrm{L}_\lambda:\quad \min_{\beta}\ &\tfrac12\,\norm{y - X\beta}_2^2 + \lambda\,\norm{\beta}_1 \\
\text{s.t.}\quad &A\beta = b,\ \ G\beta \le h .
\end{aligned}
\end{equation}

\begin{theorem}[CLA--LASSO identity under general constraints]
\label{thm:identity}
\index{identity!CLA and LASSO}
Suppose $\Sig = X\T X$ and $\mu = X\T y$. Then for every $\lambda \ge 0$ there is a
$c(\lambda) = \norm{\beta(\lambda)}_1$, nonincreasing in $\lambda$, such that
$\mathrm{M}_{c(\lambda)}$ and $\mathrm{L}_\lambda$ have the same minimiser. The
solution path of $\mathrm{M}_c$ as $c$ varies and the regularization path of
$\mathrm{L}_\lambda$ as $\lambda$ varies are therefore the \emph{same}
piecewise-linear curve in $\R^n$, traced by the same active-set homotopy and sharing
the same breakpoints as points of $\R^n$.
\end{theorem}

\begin{proof}
Completing the square,
\[
  \tfrac12\norm{y-Xw}_2^2 = \tfrac12\, w\T (X\T X)\, w - (X\T y)\T w + \tfrac12\norm{y}_2^2 .
\]
With $\Sig = X\T X$ and $\mu = X\T y$ the objective of $\mathrm{M}_c$ equals
$\tfrac12\norm{y-Xw}_2^2$ up to the additive constant $\tfrac12\norm{y}_2^2$, which
does not affect the minimiser. Hence, under the same linear constraints,
$\mathrm{M}_c$ is exactly the constrained least-squares problem subject to
$\norm{w}_1 \le c$. Its Lagrangian relaxation of that single constraint is
$\mathrm{L}_\lambda$. Because the objective is convex and $\norm{\cdot}_1$ is convex,
strong duality holds for the norm-ball constraint, so the constrained and penalised
forms have identical solution sets as $c$ and $\lambda$ range over $[0,\infty)$, with
$c(\lambda) = \norm{\beta(\lambda)}_1$ the active constraint level at the dual value
$\lambda$; $c$ is nonincreasing because larger penalties shrink $\norm{\beta}_1$.
Both objectives are quadratic and both feasible sets polyhedral, so by
Lemma~\ref{lem:affine} each path is piecewise linear; sharing a minimiser at every
parameter value, the two paths are the same set of points with the same corners, and
tracing either by the homotopy visits the same vertices in the same order.
\end{proof}

\begin{remark}[When the substitution is available]
The hypothesis is mild and runs in both directions. From the regression side it holds
by definition. From the portfolio side, where one starts with $(\Sig,\mu)$, any
$\Sig \succeq 0$ factors as $X\T X$ (a symmetric square root or Cholesky factor
serves), and $\mu = X\T y$ is solvable for $y$ precisely when
$\mu \in \range(X\T)$. When $\Sig \succ 0$---the finance regime---$X$ has full column
rank and a representing $y$ always exists, so the theorem applies to the
mean--variance problem for every $(\Sig,\mu)$. The only obstruction can occur for
singular $\Sig$, and bites in neither instance.
\end{remark}

\subsection{The parameter map}

The map $c(\lambda)$ can be pinned down, which both sharpens the theorem and isolates
the only way the parameter bijection can fail.

\begin{proposition}[Regularity of the $c$--$\lambda$ correspondence]
\label{prop:reg}
Let $\beta(\lambda)$ be the constrained LASSO path and
$c(\lambda) = \norm{\beta(\lambda)}_1$. Then $c$ is continuous and piecewise linear,
affine on each segment with $c'(\lambda) = -s_\Active\T\beta_{\mathrm{slope},\Active}$
over the current support $\Active$ with signs $s_\Active$. In the unconstrained case
$\beta_{\mathrm{slope},\Active} = (X_\Active\T X_\Active)^{-1}s_\Active$, so
$c'(\lambda) < 0$ on every segment with nonempty support and $c$ is a strictly
decreasing, piecewise-linear bijection. In the constrained case
$s_\Active\T\beta_{\mathrm{slope},\Active} \ge 0$, with strict inequality---and hence
the bijection---whenever $s_\Active$ is not annihilated by the reduced Hessian on the
active constraint tangent space.
\end{proposition}

\begin{proof}
On a segment the support and signs are fixed, so
$\norm{\beta(\lambda)}_1 = s_\Active\T(\alpha_\Active - \lambda\beta_{\mathrm{slope},\Active})$,
affine with the stated slope; continuity follows from continuity of $\beta(\lambda)$.
Unconstrained, segment stationarity gives
$\beta_{\mathrm{slope},\Active} = (X_\Active\T X_\Active)^{-1}s_\Active$ and the
quadratic form is positive on a nonempty support. In the constrained case,
eliminating $\zeta$ from~\eqref{eq:border} by its Schur complement expresses the free
block's response to a right-hand side $(g;0)$ as $Z(Z\T H_{\Free\Free}Z)^{-1}Z\T g$,
where the columns of $Z$ span $\nullsp(C_\Free)$. With $g = s_\Active$,
\[
  s_\Active\T\beta_{\mathrm{slope},\Active}
  = \bigl\|(Z\T H_{\Free\Free}Z)^{-1/2}Z\T s_\Active\bigr\|_2^2 \ge 0,
\]
with equality only when $Z\T s_\Active = 0$, i.e.\ when $s_\Active$ lies in the span
of the active normals.
\end{proof}

\subsection{The gross-exposure hinge}

The identity has a structural reading. In $\mathrm{M}_c$ the sparsity is governed by
an \emph{explicit} constraint, the gross-exposure ball; in $\mathrm{L}_\lambda$ it is
\emph{induced} by a penalty through its subgradient. The next proposition records
when the explicit one becomes vacuous, and it is the reason the long-only frontier
looks nothing like a LASSO path even though they are the same object.

\begin{proposition}[The gross-exposure hinge]
\label{prop:hinge}
For a long-only, fully invested portfolio ($w \ge 0$, $\onevec\T w = 1$) the
gross-exposure constraint is vacuous, since $\norm{w}_1 = \onevec\T w = 1$ for all
feasible $w$. The long-only mean--variance problem is therefore a
\emph{non-negative} LASSO with the equality $\onevec\T\beta = 1$, and it traces a
frontier rather than a sparsity path. Short positions admitted under a binding cap
$\norm{w}_1 \le c$ with $c > 1$ make the constraint active, and then $\mathrm{M}_c$
and $\mathrm{L}_\lambda$ coincide as in Theorem~\ref{thm:identity}.
\end{proposition}

\begin{proof}
Immediate from $\norm{w}_1 = \sum_i w_i = \onevec\T w$ when $w \ge 0$; the budget
fixes $\norm{w}_1 = 1$ so $\norm{w}_1 \le c$ cannot bind for $c \ge 1$, and for
$c < 1$ the long-only budget is infeasible.
\end{proof}

\subsection{The third axis}

The two-dimensional risk--return picture has a \emph{third} axis, and drawing it
heads off a likely misreading. The gross-exposure formulation adds
$\norm{w}_1$, a measure of leverage; the efficient set is then a \emph{surface} in
(return, risk, leverage). The frontier and the LASSO path are perpendicular cuts of
that surface. The Critical Line Algorithm fixes the leverage structure and sweeps
the return; Brodie et al.\ \cite{brodie2009} fix the return and sweep the leverage
penalty, following the fit-against-sparsity path rather than the risk--return
frontier. Their path meets the frontier at a single point, its zero-penalty end;
raising the penalty walks off the frontier toward sparser, less-levered, slightly
riskier portfolios.

The third axis collapses when the weights are pinned, which is
Proposition~\ref{prop:hinge}: long-only and fully invested freezes leverage, and one
is back to the two-dimensional frontier with the active set driven by box bounds.
The leverage axis is genuinely live only for long--short books.

\section{A taxonomy}
\label{sec:taxonomy}

With the identity in hand, the right way to see the family is not as a list of
methods from different fields but as points in a small design space. What the members
share is the common scheme---parametric QP, affine segment, homotopy, simplex-inherited
anti-cycling---and that part is fixed, and is most of the algorithm. What varies is
captured by four features. Read them as \emph{coordinates}, not as a checklist: a
choice of all four specifies a method, and every member of the family, including
combinations no field has yet had reason to name, is some such choice.
\index{taxonomy of path methods}

\emph{Imposed versus induced active sets.} Constraints impose faces; penalties induce
them through a subgradient. Theorem~\ref{thm:identity} and
Proposition~\ref{prop:hinge} show the two are interconvertible, with the
gross-exposure ball the hinge, but their computational ergonomics differ, which is
why each field built its own version.

\emph{Monotone versus non-monotone active sets.} Forward selection and forward LARS
only \emph{add} variables; the LASSO modification and the CLA both add and drop.
The CLA corresponds to the full LASSO homotopy, not to forward LARS.

\emph{The structure of $H$.} A covariance in finance, a Gram matrix in regression, a
kernel in learning. What $H$ is fixes both interpretation and cost, and---crucially---
the homotopy never needs it explicitly (Chapter~\ref{ch:operator}).

\emph{Degeneracy.} Ties and near-rank-deficiency are the recurring practical hazard;
the inherited Bland rule guarantees determinism but, as in the simplex method, does
not tame the adversarial worst case.

Most cells of that space are occupied, several of them more than once by different
fields---the LASSO and the CLA sitting at nearly the same point and recognised as one
only through Theorem~\ref{thm:identity}. Some are barely explored: kernelised problems
under hard linear constraints, traced exactly rather than on a grid, are an imposed,
non-monotone, kernel-$H$ combination that no field has had much occasion to build, yet
the scheme constructs it with no new machinery. And because the axes are shared, a
result, an algorithmic refinement, or a piece of software attached to one cell carries
to its neighbours along them.

\section{What the identity buys, and what it does not}
\label{sec:transfer}

The force of Theorem~\ref{thm:identity} is not that two programs happen to share a
minimiser at each parameter value; it is that the two solution paths are the
\emph{same curve}. What that licenses is narrower than it first appears, and the
boundary is worth drawing precisely, because it is a good lesson in what a
mathematical identity does and does not transfer.

\paragraph{Statements about the curve transfer verbatim.} That the solution is unique
at a given point, that the two homotopies visit the same active sets in the same
order, that a segment is affine, that the path has at most so many breakpoints: each
is a geometric fact about a subset of $\R^n$, and a bound on the number of LASSO
breakpoints is at once a bound on the number of turning points of the frontier.
Algorithms and code transfer for the same reason, which is why one engine traces both.

\paragraph{Statements about a fixed value of the parameter do not.} Here the
reparametrisation $\lambda \mapsto c(\lambda)$ stops being a convenience and becomes
an obstruction. The map is $c(\lambda) = \norm{\beta(\lambda)}_1$: it is computed
\emph{from the data}, so it is not a fixed relabelling of an axis but a \emph{random}
one. Holding $\lambda$ fixed while the response varies and holding $c$ fixed while
the response varies are two different experiments, and a quantity defined by averaging
over the response is attached to the experiment, not to the curve.

Degrees of freedom is the clearest case. The general results for the lasso and
generalized lasso \cite{tibshirani2012df,zou2007df} concern the Lagrange form at a
fixed penalty; they do not by themselves yield the corresponding statement for the
constrained form at a fixed bound, because the estimator whose degrees of freedom one
would be computing is a different function of the response. The same caution applies
to post-selection inference \cite{lee2016,tibshirani2016psi} and to any risk or
model-selection criterion that integrates over the response.

\emph{The identity is an identity of curves; it is not an identity of statistical
experiments.} Results of the second kind must be proved in the parametrisation where
they are wanted. Chapter~\ref{ch:reading} takes up what survives.

\begin{notes}
This chapter follows \cite{schmelzer2026perspective}. The unconstrained
correspondence between the LASSO path and least angle regression is classical
\cite{efron2004}; the single gross-exposure case is Brodie et al.\ \cite{brodie2009};
the constrained-LASSO path under general linear equalities and inequalities is Gaines,
Kim and Zhou \cite{gaines2018}. What Theorem~\ref{thm:identity} adds is the
identification of the bound-driven frontier tracer as Markowitz's Critical Line
Algorithm---one breakpoint sequence read in both directions---together with
Proposition~\ref{prop:reg}. Rosset and Zhu \cite{rosset2007} give the general
characterisation of piecewise-linear regularised paths, and generalise in a direction
not pursued here, to piecewise-quadratic and piecewise-linear losses such as the
Huber, hinge, and quantile losses.
\end{notes}

\begin{exercises}

\begin{exercise}
Verify Theorem~\ref{thm:identity} numerically on a small design: trace the LASSO path
by any method, then solve $\mathrm{M}_c$ as a quadratic program at each
$c = \norm{\beta(\lambda)}_1$ and compare the minimisers.
\end{exercise}

\begin{exercise}
Show that $\mu \in \range(X\T)$ is necessary and sufficient for $\mu = X\T y$ to be
solvable, and give a singular $\Sig$ and $\mu$ for which it fails.
\end{exercise}

\begin{exercise}
Prove Proposition~\ref{prop:hinge} and then explain, in one sentence each, why a
long-only frontier still has breakpoints and why they are not sparsity events in the
LASSO sense.
\end{exercise}

\begin{exercise}
Give an example of a constraint system for which $Z\T s_\Active = 0$ in
Proposition~\ref{prop:reg}, so that $c$ has a flat segment. Why is this non-generic?
\end{exercise}

\begin{exercise}
Locate each of the five instances of Section~1 in the four-coordinate taxonomy of
Section~\ref{sec:taxonomy}. Which cells are unoccupied, and can you name a problem
that would occupy one?
\end{exercise}

\begin{exercise}
The forward LARS active set is monotone. Which of the four event families of
Chapter~\ref{ch:parametric} does forward LARS populate, and which does it suppress?
What is lost by suppressing it?
\end{exercise}

\begin{exercise}
Explain, using Section~\ref{sec:transfer}, why ``the number of breakpoints of the
LASSO path is at most $K$'' transfers to the frontier but ``the expected number of
nonzero coefficients at $\lambda$ equals the degrees of freedom'' does not.
\end{exercise}

\begin{exercise}[Implementation]
Take your path tracer from Chapter~\ref{ch:parametric} and drive it twice: once with
$H = \Sig$, $d = \mu$ for a small portfolio problem, and once with $H = X\T X$,
$d = X\T y$ for a small regression. Confirm that a single implementation serves both,
and compare the second against a reference LARS implementation breakpoint by
breakpoint.
\end{exercise}

\begin{exercise}[Implementation]
Trace the leverage axis: fix a return target, add the constraint $\mu\T w = \rho$,
and sweep an $\ell_1$ penalty. Plot the resulting path in (risk, leverage) and confirm
it meets the frontier only at the zero-penalty end, as Section~3.3 claims.
\end{exercise}

\end{exercises}

\chapter{Reading the Path}
\label{ch:reading}

A single quadratic program returns one answer. Re-solving on a grid returns a
scatter of them. The homotopy returns the whole curve exactly, and the curve carries
readouts that no single optimum can. This chapter is about those readouts: what to
plot, what it means, what a statistician already knows about it that a portfolio
manager does not, and---the part that requires discipline---which of those meanings
survive the translation of Section~\ref{sec:transfer} and which do not.

It closes with the open problems, because a graduate reader is entitled to know
where the subject stops.

\section{The coefficient profile}

Plot each coordinate of $x(\lambda)$ against $\lambda$. In regression this is the
LARS/LASSO coefficient profile, the picture a statistician reads to see which
variables a model admits and in what order. In finance the same plot, drawn as the
frontier is traced from minimum variance to maximum return, shows exactly how each
position enters, grows, and leaves, with the turning points as entry and exit events:
a membership and turnover map of the efficient set.
\index{coefficient profile}

It is the same plot. Under Theorem~\ref{thm:identity} it \emph{is} the coefficient
profile when $H = X\T X$. The kinks are breakpoints; a coordinate's segment slope is
$\beta_i$ from Lemma~\ref{lem:affine}; a coordinate pinned at a bound is flat with
$\beta_i = 0$.

Two practical readings that the exactness makes trustworthy. First, an investor
operates in a band of volatility rather than over the whole frontier, and the exact
path lets one trace only that band and read the precise weights at every point inside
it, paying only for the turning points the band contains. Second, because segments
are affine in \emph{weight} space, any intermediate portfolio is an exact convex
combination of its two bracketing turning points---no interpolation error, and no
re-solving.

\section{The support size}

Count the coordinates held strictly between their bounds at each point of the path.
In regression this is the number of nonzero coefficients; in finance it is the
number of positions held. Again the same quantity.
\index{support size}

The count falls in a characteristic knee shape: many holdings shed for little added
risk up to a corner, then a long flat tail of one or two concentrated names. An
investor can read off, at any target volatility, how many names the optimum will
hold, or invert the reading and ask what risk a given cardinality buys. That is a
complexity profile of the efficient set, free from the exact path, and the portfolio
literature had no occasion to define it because it had no exact path to read it from.

Here the discipline of Section~\ref{sec:transfer} is required, and the temptation to
be resisted is specific and attractive.

\begin{remark}[Degrees of freedom: what transfers and what does not]
\label{rem:df}
\index{degrees of freedom}
As a \emph{descriptive} complexity measure the count transfers immediately: the
number of positions a frontier portfolio holds strictly between its bounds is a
well-defined effective dimension for that portfolio, computed along the path at no
cost.

The \emph{distributional} statement is another matter. That the support size is an
unbiased estimate of the degrees of freedom of the fit
\cite{zou2007df,tibshirani2012df} is a theorem about the Lagrange-form estimator at a
\emph{fixed penalty} $\lambda$, with the response Gaussian about $X\beta_0$. The
efficient frontier as the Critical Line Algorithm traces it is indexed by the
gross-exposure bound $c$ (equivalently by target return), which is the
\emph{constrained} form; and $c(\lambda)$ depends on the response. So the unbiasedness
does \emph{not} transfer to a frontier point at fixed $c$.

Where it should transfer is the \emph{penalised} portfolio problem,
$\min \tfrac12 w\T\Sig w - \mu\T w + \tau\norm{w}_1$ at fixed leverage penalty
$\tau$---the Brodie et al.\ slice of Section~3.3 rather than the frontier itself:
there the tuning parameter is held fixed as the data vary and the experiment is the
right one. Even there the extension is not immediate, since the general-constraint
case is outside what \cite{tibshirani2012df} establishes. \emph{The leverage axis, not
the frontier, is where to look.}
\end{remark}

The same caution applies to model selection. Because the count is the LASSO's support
size, its knee is recognisably a model-selection elbow, and the devices statistics
uses to choose a penalty by the corner of a complexity-versus-fit curve---the L-curve
criterion \cite{hansen1992}, information criteria on the degrees of freedom---suggest
a principled rule for choosing where to sit on the frontier. \emph{Suggest} is the
right verb: an information criterion needs the degrees-of-freedom statement, which
holds for the penalised parametrisation and not for the bound-indexed frontier.
Turning the elbow into a rule, in the parametrisation where the statistics is valid,
is real work and not bookkeeping.

\section{Selection as the choice of a corner}

Post-selection inference conditions on the event that a particular model was
selected, which for the LASSO is a polyhedron in the response space
\cite{lee2016,tibshirani2016psi}. Under Theorem~\ref{thm:identity} that polyhedron is
precisely the set of data for which a given corner portfolio is optimal, so each
breakpoint of the frontier is a selection event of exactly the form the
selective-inference machinery treats. This much is a statement about the geometry of
the path, which the identity does license.

The caveat is the same one and it is not cosmetic: the conditioning event must be
described in the parametrisation in which the analysis is carried out, so conditioning
on the selected support at fixed $\lambda$ and conditioning on it at a fixed
gross-exposure cap are different events, and the second is the one a portfolio manager
is in. The apparatus should carry over; the polyhedron it conditions on has to be
rewritten first.

We flag this as the most concrete of the transfers the unified view makes available,
and as sketch rather than finished argument.

\section{From one path to a cloud of them}

Because a single path is now cheap---an exact whole curve from one homotopy, and
cheaper still under the operator reading of Chapter~\ref{ch:operator}, where a factor
or kernel form costs $O(nk)$ rather than $O(n^2)$ per step---one can afford to compute
not one but hundreds.

Resample the inputs: draw $(\mu,\Sig)$ from a bootstrap, a Bayesian posterior, or a
rolling estimation window; trace the exact path for each. The ensemble is a
\emph{cloud} of paths rather than a single curve.
\index{resampled frontier}

The exactness is what lets the cloud be read coherently. Each path is affine between
its turning points, so evaluating all of them at a shared abscissa---a target
volatility, or the $\ell_1$ budget of Theorem~\ref{thm:identity}---is \emph{exact}
interpolation, and at each level one obtains a genuine cross-sectional distribution of
weights and values across the draws, not the scatter a grid of separately solved
programs would leave.

Every readout of this chapter then lifts from a point to a distribution. The
coefficient profile acquires a band showing how stable each holding is under
estimation error. The support-size knee becomes a \emph{distribution} of elbows, a
direct picture of how firmly the model-selection operating point is pinned. And the
resampled median path sits systematically \emph{above} the true one in the finance
instance, which is the in-sample optimism a single estimated frontier hides and the
cloud exposes.

The ensemble view of estimation error is itself old---it is Michaud's resampled
efficiency \cite{michaud1999}. What the common scheme adds is that each member is
exact and nearly free, so the cloud is affordable at factor-model scale, and that its
readouts carry the regression meanings above through the identity: an estimation-error
band on a frontier is a confidence statement about a LASSO path read in portfolio
coordinates.

\section{Why exactness is not a numerical nicety}

It is tempting to regard the exact path as a refinement of a grid: nicer, but the
grid would do. It would not, and the reason is specific.

Gridding the parameter and interpolating between grid points cuts \emph{the very
corners} that carry the model-selection information. The breakpoints are where the
active set changes---where a variable enters, where a position is dropped---and they
are precisely what a grid can miss. Coordinate-descent and proximal solvers on a
$\lambda$-grid approximate the path in exactly this way; the homotopy returns the
corners themselves.

There is a numerical detail that makes the difference real rather than notional. The
breakpoints of a large path crowd together, so the ratio test that orders them must
use a tolerance \emph{relative} to the scale of $\lambda$. With it, the path comes back
strictly monotone to floating-point precision; with a fixed absolute tolerance,
distinct breakpoints eventually merge and the parameter steps backwards.
Chapter~\ref{ch:numerics} generalises the lesson.

\section{Open problems}

The family poses its open problems once, which is one of the benefits of seeing it as
one family.

\paragraph{Path length.} The number of breakpoints is near-linear in the problem size
in practice, which is the homotopy's reason for being. Yet Mairal and Yu
\cite{mairal2012} construct LASSO paths whose length is \emph{exponential} in the
dimension, the statistical mirror of the long-known exponential bounds on the critical
regions of a multi-parametric QP \cite{tondel2003}. What separates the benign average
from this worst case, when tracing the whole path beats solving at a single parameter,
and whether the Bland tie-break can be replaced by a rule with a worst-case guarantee,
are open in the shared form. One resolution would serve every instance.

\paragraph{Rank deficiency.} Remark~\ref{rem:rank} set this up. The scheme assumes a
positive-definite reduced Hessian at every visited active set, and for the plain LASSO
that assumption can be removed \cite{tibshirani2013} by tracing the
minimum-$\ell_2$-norm solution, on the strength of an insertion--deletion property of
that selection rather than of the segment structure. No counterpart is known for the
generalized or constrained problem, where the penalty matrix and the constraint system
may be arbitrary but the design may not. Whether the bordered system admits a canonical
selection playing the same role, or whether the arbitrary-design LASSO path is
genuinely a different construction that the common scheme does not contain, is the most
interesting question this material leaves open---and a caution against reading the
scheme as more general than it is.

\paragraph{Constraints and breakpoint counts.} Proposition~\ref{prop:reg} shows the
penalty and constraint parametrisations agree except on flat segments, which are
non-generic. Which constraint systems force them on a positive-measure set, and
whether $p$ inequality rows can \emph{multiply} rather than merely add to the
breakpoint count, are not known.

These are questions about the shared object rather than about portfolios or
regressions, which is the point of having identified it.

\begin{notes}
The readouts of Sections~1--4 follow \cite{schmelzer2026perspective}; the corrections
of Remark~\ref{rem:df} are its own, and are worth reading as a case study in
withdrawing an overreaching claim precisely rather than quietly. Resampled efficiency
is \cite{michaud1999}; the L-curve criterion \cite{hansen1992}; degrees of freedom for
the lasso \cite{zou2007df,tibshirani2012df}; selective inference
\cite{lee2016,tibshirani2016psi}. The exponential path-length construction is
\cite{mairal2012}.
\end{notes}

\begin{exercises}

\begin{exercise}
Show that on a segment the objective value is quadratic in $\lambda$ but the weight
vector is affine, and explain why plotting the frontier in
(variance, return) rather than (volatility, return) coordinates makes each segment an
exact parabolic arc.
\end{exercise}

\begin{exercise}
Prove that any portfolio between two adjacent turning points is an exact convex
combination of them, and give the formula for its variance in terms of the two
endpoints' variances and their cross term.
\end{exercise}

\begin{exercise}
Explain in your own words why $c(\lambda) = \norm{\beta(\lambda)}_1$ being
data-dependent blocks the transfer of a degrees-of-freedom statement, but not the
transfer of a breakpoint-count bound.
\end{exercise}

\begin{exercise}
The maximum-Sharpe portfolio maximises $\mu\T w / \sqrt{w\T\Sig w}$ along the
frontier. Show that on an affine segment $w(t) = (1-t)w_0 + tw_1$ this is an affine
numerator over the square root of a quadratic, and derive the single stationary point
$t^\star$ in closed form. Why does unimodality along the frontier mean that checking
the two segments bracketing the best turning point suffices?
\end{exercise}

\begin{exercise}
Give an example where a grid of $50$ parameter values misses a breakpoint entirely,
and quantify the resulting error in the support size.
\end{exercise}

\begin{exercise}
Explain why a fixed absolute tolerance in the ratio test can make the traced parameter
sequence non-monotone, and design a relative tolerance that prevents it.
\end{exercise}

\begin{exercise}[Implementation]
Trace an exact frontier for $n = 30$ and plot the coefficient profile and the support
size against volatility. Mark the knee. Then re-solve at $50$ grid points with a
general QP solver and overlay: how many corners does the grid miss?
\end{exercise}

\begin{exercise}[Implementation]
Build a cloud: from a fixed population factor model, draw $300$ samples, re-estimate
$(\mu,\Sig)$ from each, trace the exact frontier for every draw, and plot the
percentile envelope of returns at each risk level together with the true frontier.
Confirm that the resampled median sits above the truth, and explain why.
\end{exercise}

\end{exercises}


\part{Structure and Conditioning}
\chapter{The Quadratic Form as an Operator}
\label{ch:operator}

The point that does the most work in practice is the easiest to miss, because every
textbook presentation hides it. The reader is trained to see a \emph{matrix}: a
covariance in finance, a Gram matrix in regression, a kernel in learning---formed,
stored, and factored.

None of the algorithms in this book sees that matrix. Inspecting them,
the quadratic form enters only through a short list of linear \emph{actions}. So the
quadratic form is, for the algorithm, an \emph{operator}: defined by what it does to a
vector, not by a table of $n^2$ numbers. That change of model is what turns one
algorithm into a method for many problems, because the operator may be supplied in
whatever form is cheapest.

This chapter states the contract, gives the four backends that recur, works out the
cost of the most important one, and says where the crossover falls.

\section{The contract}

\begin{definition}[Quadratic-form operator]
\label{def:qform}
\index{operator!quadratic-form contract}
An implementation of the quadratic form $H$ supplies three linear actions:
\begin{enumerate}
\item \texttt{matvec}: $v \mapsto Hv$;
\item \texttt{solve\_free}: $R \mapsto H_{\Free\Free}^{-1}R$ for a multi-column
      right-hand side $R$ and a given free set $\Free$;
\item \texttt{cross}: $x_\Bound \mapsto H_{\Free\Bound}\,x_\Bound$.
\end{enumerate}
\end{definition}

That is the whole interface, and it is worth checking against the algorithms already
seen. The dual method of Chapter~\ref{ch:walking} carries $J$ with $JJ\T = G^{-1}$ and
consumes it only through products---so its needs are a superset of \texttt{matvec} but
the same in kind. The block-pivoting loop of Chapter~\ref{ch:guessing} needs
\texttt{solve\_free} for the inner system and \texttt{cross} for the reduced gradient
at bound indices; Section~\ref{sec:operator-two} of Chapter~\ref{ch:krylov} named
exactly those two. The homotopy of Chapter~\ref{ch:parametric} needs all three: the
bordered solve is \texttt{solve\_free}, the blocked contribution
$H_{\Free\Bound}x_\Bound$ on the right-hand side is \texttt{cross}, and the reduced
gradients that drive the $\mathrm{D}_1$ events are \texttt{matvec}.

\emph{One interface serves the whole book.} That is not a coincidence; it is what the
active-set structure permits, and it is the technical content of the claim in
Chapter~\ref{ch:homotopy} that one engine traces five literatures' paths.

\section{Four backends}

\subsection{Dense}

An explicit symmetric $n\times n$ array. \texttt{matvec} is a matrix--vector product;
\texttt{cross} is an off-diagonal block product; \texttt{solve\_free} is a Cholesky
factorisation of the principal submatrix $H_{\Free\Free}$, falling back to a general
solve on the rare block that is not numerically positive definite. This is the
reference backend and the one against which the others are validated.

Cost: $O(n^2)$ storage, $O(n_\Free^3)$ per free-block solve.

\subsection{Incremental dense}

The same explicit matrix, but \emph{maintaining} $H_{\Free\Free}^{-1}$ rather than
refactorising at every step. Because an active-set method changes the free set by one
index between consecutive solves, a rank-one bordered update (on entry) or a symmetric
deletion (on exit) modifies the stored inverse in place at $O(n_\Free^2)$, against
$O(n_\Free^3)$ for a fresh solve.

The trade is numerical, and it is worth stating plainly because it is the recurring
shape of this decision. A maintained inverse accumulates floating-point error over the
$O(n)$ updates of a trace, whereas a fresh solve is clean at every step. A guard is
needed---a non-positive or non-finite pivot, or any free-set change that is not a
single-index flip, triggers a full refactorisation---and even with it, on
ill-conditioned or near-degenerate inputs the plain dense backend is the safer choice.
So this is an opt-in fast path, not a default. Roughly a factor of two on a
well-conditioned dense problem of a few hundred variables.

\subsection{Factor: diagonal-plus-low-rank}
\label{sec:factor}

\begin{equation}
\label{eq:factor}
H \;=\; \diag(d) \;+\; U\,\Delta\,U\T,
\qquad d \in \R^{n}_{>0},\; U \in \R^{n\times k},\; \Delta \in \R^{k\times k},
\end{equation}
with $k \ll n$. This is the form taken by a factor risk model \cite{sharpe1963}, by a
random-matrix-cleaned covariance (Chapter~\ref{ch:conditioning}), by an approximate
factor model \cite{chamberlain1983} and its modern estimators
\cite{fan2013poet}, and by a low-rank kernel approximation. One object, arrived at
independently by finance, statistics, and machine learning.
\index{factor model}

The free-block solve never forms $H_{\Free\Free}$. By Proposition~\ref{prop:woodbury},
\begin{equation}
\label{eq:woodbury}
H_{\Free\Free}^{-1} \;=\; D_\Free^{-1} - D_\Free^{-1} U_\Free\, W^{-1}\, U_\Free\T D_\Free^{-1},
\qquad
W \;=\; \Delta^{-1} + U_\Free\T D_\Free^{-1} U_\Free \;\in\;\R^{k\times k},
\end{equation}
so a solve costs $O(n_\Free k^2 + k^3)$ rather than $O(n_\Free^3)$, and no $n\times n$
matrix is ever allocated: memory and per-solve cost scale as $O(nk)$ instead of
$O(n^2)$. The cross product contributes only through the low-rank part, the diagonal
having no off-diagonal block:
$H_{\Free\Bound}x_\Bound = U_\Free\Delta(U_\Bound\T x_\Bound)$.

\subsection{Gram: straight from the data matrix}
\label{sec:gram}

A sample covariance \emph{is} a Gram matrix. If $X_c \in \R^{T\times n}$ is the
column-centred data,
\begin{equation}
\label{eq:gram}
H \;=\; \frac{1}{T-1}\,X_c\T X_c,
\end{equation}
and the contract is implemented straight from $X_c$: \texttt{matvec} is the pair of
thin products $X_c\T(X_c v)/(T-1)$, \texttt{cross} restricts those products to free
and blocked columns, and \texttt{solve\_free} works on $X_{c,\Free}$ alone. Memory is
$O(Tn)$ rather than $O(n^2)$. This is the backend Chapter~\ref{ch:krylov} runs CG
against.

The catch is rank: $X_c\T X_c$ has rank at most $T-1$. When $T \ge n$ with full column
rank, $H \succ 0$ and this is the dense reference at lower memory. In the
short-sample regime $T < n$---a covariance estimated from far fewer periods than
assets---$H$ is singular and the free block becomes unsolvable once the free set
outgrows the data rank. An optional ridge $\delta I$ restores positive definiteness,
and the free-block solve is then \eqref{eq:woodbury} carried out in the
$T$-dimensional observation space, at $O(n_\Free T^2 + T^3)$. \emph{This is precisely a
factor model whose factors are the observations themselves}, with $U = X_c\T$ and
$k = T$.

\begin{remark}[A memory benefit, not an unconditional speed-up]
When $T \ge n$ and the dense matrix fits in memory, the dense backend is faster per
step: it slices a precomputed $H_{\Free\Free}$, whereas the Gram backend re-forms
$X_{c,\Free}\T X_{c,\Free}$ at each solve. The same caveat governs the factor backend:
a structured solve pays only while $k$ (or here $T$) stays well below $n$. At full
rank the structured route is \emph{slower} than the plain dense matrix. State this
clearly in any benchmark you report.
\end{remark}

\subsection{Kernels, and sparse designs}

Two more, briefly, because they cost nothing to add. A kernel Gram matrix is given
implicitly by its kernel function; the support-vector-machine path of
Chapter~\ref{ch:homotopy} touches it only through the same actions, so the kernel trick
enters unchanged and the feature space may be infinite-dimensional without the
algorithm noticing. And when $X$ is sparse, the action $v \mapsto X\T(Xv)$ costs the
number of nonzeros rather than $n^2$; the operator inherits the data's sparsity with no
special handling.

\section{The crossover}

When is the structured solve actually cheaper? The question has a clean answer for the
scalar-identity-plus-low-rank case, which is the one Chapter~\ref{ch:conditioning}
produces.

\begin{proposition}[Woodbury versus assemble-and-factor]
\label{prop:crossover}
\index{crossover!Woodbury versus Cholesky}
Assembling $H_{\Free\Free}$ by a rank-$k$ update and factoring it costs
$n_a^2 k + n_a^3/3$ flops, where $n_a = |\Free|$; forming and factoring the $k\times k$
correction $W$ of~\eqref{eq:woodbury} costs $n_a k^2 + k^3/3$. Writing $x = k/n_a$, the
Woodbury solve is cheaper per step iff
$n_a k^2 + \tfrac{k^3}{3} < n_a^2 k + \tfrac{n_a^3}{3}$, which on dividing by $n_a^3$
factors as $(x-1)(x^2 + 4x + 1) < 0$, i.e.\ iff $x < 1$, that is $k < n_a$.
\end{proposition}

The condition $k < n_a$ is embarrassingly mild, and it holds across essentially the
whole of a long-only frontier. On S\&P~500 data, $k = 17$ signal factors against
$n_a = 46$ assets held at the minimum-variance solution, so $x \approx 0.37$. On a
synthetic frontier $x$ runs from about $0.33$ at the high-return end down to about $1$
at the tightest point, where the number of holdings dips to roughly $k$ itself. Only at
that extreme are the two comparable; elsewhere the flop advantage applies, and the
measured gap is larger still---consistent with the $k\times k$ working set fitting in
lower cache levels than an assembled $n_a \times n_a$ system.

Two practical notes. When $n_a \le k$ a solver should fall back to the assemble branch.
And the Woodbury solve itself remains \emph{valid} for all $n_a \ge 1$, since
$W = \Delta^{-1} + \bar\mu^{-1}U_\Free\T U_\Free$ is a positive definite diagonal plus
a positive semidefinite matrix, hence always invertible; only its cost advantage
disappears.

\begin{remark}[Sharing a factorisation across right-hand sides]
The $p+1$ right-hand sides of an equality-augmented solve
(Chapter~\ref{ch:krylov}, Section~3) all share the single factorisation of $W$, so each
right-hand side beyond the first costs only $O(n_a k)$. Structure the code so that
$W$ is factored once per outer step and the solves are applied to a block, not looped.
This is the difference between $p+1$ solves and one solve with $p+1$ columns, and on a
factor model it is the whole of the cost.
\end{remark}

\section{Regularisation through the operator}
\label{sec:reg-operator}

A regularised objective is supported whenever the penalty is \emph{quadratic}, and for
exactly the reason the bare algorithms work: a quadratic penalty changes only the
Hessian, so the systems stay linear and---in the parametric case---the path stays
piecewise linear. Ridge is the canonical case,
\begin{equation}
\label{eq:ridge}
\tfrac12\, w\T H\, w - \lambda\,\mu\T w + \tfrac{\gamma}{2}\,\norm{w}_2^2
\;=\;
\tfrac12\, w\T(H + \gamma I)\, w - \lambda\,\mu\T w,
\end{equation}
which is the substitution $H \mapsto H + \gamma I$; more generally $H + \gamma M$ for
any $M \succeq 0$---a factor-exposure penalty, a deviation-from-benchmark term, a
quadratic turnover penalty $\norm{w - w_{\mathrm{prev}}}_2^2$.

Because the algorithms reach the quadratic form only through
Definition~\ref{def:qform}, the regularised form is just another operator: the penalty
is supplied by the \emph{backend} and the loop is untouched. The result is the exact
solution, or the exact path, of the regularised objective. Two backends already realise
this: the Gram backend takes a ridge directly, and the factor
form~\eqref{eq:factor} is itself a structured ridge, positive definite by the
idiosyncratic floor $d > 0$.

The boundary follows from the same requirement, and is worth knowing.

\begin{itemize}
\item Shrinkage toward a \emph{nonzero} target $w_0$ is still quadratic but adds a
      $\lambda$-independent linear term $-\gamma w_0\T w$, which an interface whose
      linear term is exactly $-\lambda\mu$ does not expose. Fixable, but not free.
\item A nonsmooth $\ell_1$ penalty is not a term in this objective at all. For the
      long-only fully invested book it is vacuous
      (Proposition~\ref{prop:hinge}); where it bites, it is traced by the LASSO
      homotopy of Chapter~\ref{ch:homotopy}; and a gross-exposure \emph{constraint}
      $\norm{w}_1 \le c$ is polyhedral and expressible as inequality rows.
\item Genuinely non-quadratic penalties---entropy, a log barrier---break the structure
      and lie outside the method.
\end{itemize}

\section{What the operator reading is for}

Two claims, and they are the reason this chapter exists.

\paragraph{Cost.} In every case the algorithm is the same loop and returns the
identical answer; only the cost of a step changes, from $O(n^2)$ to $O(nk)$ or to the
number of nonzeros. That is how a method designed in 1956 for a dense covariance
reaches universe sizes its dense formulation could not.

\paragraph{Transfer.} Seeing the members as one says that the factor risk model and the
kernel trick are the \emph{same device}---an operator standing in for a matrix---so
either field's machinery for it is immediately available to the other. And the traffic
is asymmetric, because each field has matured a tool the others under-use. In finance
the quadratic form is almost always a \emph{modelled} object, a diagonal-plus-low-rank
risk model being the industry default, since forming and trusting a dense covariance at
scale is infeasible; sparse regression, kernels aside, takes the Gram matrix as whatever
the design delivers and seldom models it as structured. The habit of replacing the
quadratic form by a cheap structured operator is one that regression can import
wholesale from finance. The traffic runs the other way too: inequality constraints and
the path's inferential apparatus are routine in statistics and absent from the canonical
portfolio codes.

Handed a factor covariance together with $X\T y$, the same engine traces a least angle
regression path in $O(nk)$ per step; under a linear inequality it traces a constrained
path that standard LARS implementations do not provide. These are the under-occupied
cells of the design space of Section~\ref{sec:taxonomy}: the same engine with a
different operator and one extra event type, not new algorithms.

\begin{notes}
The operator contract of Definition~\ref{def:qform} is the \texttt{QuadraticForm}
protocol of the \texttt{cvxcla} package \cite{cvxcla}, whose four backends are those of
Section~2; the reading of it as a cross-field device is
\cite{schmelzer2026perspective}. Proposition~\ref{prop:crossover} and its empirical
crossover values are from \cite{schmelzer2026rmt}. The Woodbury identity is standard
\cite{golub2013}; the diagonal-plus-low-rank covariance goes back to Sharpe's
single-index model \cite{sharpe1963} and the approximate factor model of Chamberlain
and Rothschild \cite{chamberlain1983}, with POET \cite{fan2013poet} a modern estimator
returning the same form.
\end{notes}

\begin{exercises}

\begin{exercise}
For each algorithm in Chapters~\ref{ch:walking}, \ref{ch:guessing}, \ref{ch:krylov},
and~\ref{ch:parametric}, list exactly which of the three actions of
Definition~\ref{def:qform} it uses and where.
\end{exercise}

\begin{exercise}
Verify~\eqref{eq:woodbury} directly by multiplying out
$H_{\Free\Free}\,H_{\Free\Free}^{-1}$, and check that $W \succ 0$ whenever
$\Delta \succ 0$ and $d > 0$.
\end{exercise}

\begin{exercise}
Show that the Gram backend with a ridge is the factor form~\eqref{eq:factor} with
$U = X_c\T$, $k = T$, $\Delta = I/(T-1)$, $d = \delta\onevec$. What does
Proposition~\ref{prop:crossover} then say about when it pays?
\end{exercise}

\begin{exercise}
Count the flops of the two branches in Proposition~\ref{prop:crossover} for
$n_a = 500$ and $k \in \{5, 50, 200, 500, 800\}$, and confirm the crossover at $k = n_a$.
\end{exercise}

\begin{exercise}
Show that $(H_{\Free\Free})^{-1} \ne (H^{-1})_{\Free\Free}$ in general, and give a
$2\times2$ counterexample. Which of the three actions does this fact complicate, and
why does it matter for preconditioning (Chapter~\ref{ch:conditioning})?
\end{exercise}

\begin{exercise}
Explain why a quadratic turnover penalty $\norm{w - w_{\mathrm{prev}}}_2^2$ is
expressible through the operator but shrinkage to a nonzero target is not, given the
interface's fixed linear term.
\end{exercise}

\begin{exercise}[Implementation]
Implement Definition~\ref{def:qform} as an interface with the dense and factor
backends. Verify that a solver written against the interface returns identical answers
from both when the factor backend is handed $\diag(d) + U\Delta U\T$ and the dense
backend is handed the same matrix formed explicitly. Report the agreement.
\end{exercise}

\begin{exercise}[Implementation]
Time both backends over a path trace for $n \in \{200, 500, 1000, 2000\}$ with $k = 20$
fixed, and plot per-step cost against $n$. Then hold $n = 500$ and vary $k$ to locate
the crossover empirically. Does it fall where Proposition~\ref{prop:crossover} predicts,
and if not, why not?
\end{exercise}

\end{exercises}

\chapter{Conditioning the Problem}
\label{ch:conditioning}

Every iteration count in this book is governed by a condition number. CG needs
$O(\sqrt\kappa)$ iterations (Proposition~\ref{prop:cg}); projected gradient needs
$O(\kappa)$ (Proposition~\ref{prop:proj}); a direct solve loses digits in proportion
to $\kappa$. So anything that compresses the spectrum shortens the solve, and this
chapter is about the mechanisms that do.

There are three, and they are usually taught in three different courses. \emph{
Regularisation} moves the operator toward a target the application deems preferable.
\emph{Preconditioning} accelerates the solve of the operator as given, leaving the
answer untouched. \emph{Spectral cleaning} replaces the part of the spectrum that is
estimation noise. The chapter's thesis is that on the problems of this book the three
coincide in a useful way, and that a practitioner applying the first for statistical
reasons is getting the second for free---a connection that appears not to have been
made explicit in the application literature.

The running example is a sample covariance matrix, because it is the most
spectacularly ill-conditioned operator a reader is likely to meet: on five years of
S\&P~500 daily returns, $\kappa(\hat\Sig) \approx 10^5$.

\section{Regularisation}

Replace $H$ by the convex combination
\begin{equation}\label{eq:regsplit}
  H_\alpha = (1-\alpha)\,H + \alpha\,R\T R, \qquad \alpha \in (0,1),\; R\T R \succ 0,
\end{equation}
a ridge/Tikhonov regularisation toward the SPD target $R\T R$.
Proposition~\ref{prop:weyl} already gave the effect on the condition number, and for
the scaled-identity target it is $\kappa(H_\alpha) = O((1-\alpha)/\alpha)$: doubling
$\alpha$ roughly halves the bound, \emph{independently of how small the smallest
eigenvalue of $H$ was}.
\index{regularisation}

Two features of the split matter beyond the bound.

\paragraph{It is free in factored form.} When $H = M\T M$ the regularised operator is
again a Gram matrix, of the stacked operator
\[
  \tilde{M} = \begin{pmatrix} \sqrt{1-\alpha}\,M \\ \sqrt{\alpha}\,R \end{pmatrix},
  \qquad \tilde{M}\T\tilde{M} = H_\alpha,
\]
so it is simply the Gram backend of Chapter~\ref{ch:operator} applied to $\tilde M$,
and every solver runs on it with no change. The augmented-matrix device is standard in
ridge regression and Tikhonov regularisation \cite{golub2013,lawson1974}; what is worth
noticing is that it keeps the method \emph{matrix-free}.

\paragraph{It restores the $P$-matrix property.} Because $H_\alpha \succ 0$ strictly
whenever $\alpha > 0$, every principal submatrix is invertible even when $H$ is only
positive semidefinite---a rank-deficient least-squares operator with fewer rows than
columns, say. So the split simultaneously conditions the inner solve and
\emph{licenses the finite-termination proof} of Theorem~\ref{thm:termination}, which
needs the $P$-matrix property of Chapter~\ref{ch:geometry}. A strictly positive
$\alpha$ is therefore the natural setting for the active-set loop, and the two
motivations for it are independent.

\section{Preconditioning}

Regularisation lowers $\kappa$ by \emph{changing the operator}. When no modelling
change is admissible and $H$ must be solved as given, the same reduction in iteration
count is available through preconditioning, which leaves the minimiser untouched.

An SPD preconditioner $P_\Free \approx H_{\Free\Free}$ turns the inner solve into
preconditioned CG, equivalently CG on
$P_\Free^{-1/2} H_{\Free\Free} P_\Free^{-1/2}$, whose spectrum coincides with that of
$P_\Free^{-1}H_{\Free\Free}$. PCG realises this without ever forming
$P_\Free^{-1/2}$: each iteration adds a single application of $P_\Free^{-1}$
\cite{golub2013,benzi2005}. Proposition~\ref{prop:cg} then holds with $\kappa$
replaced by $\kappa_P = \kappa(P_\Free^{-1}H_{\Free\Free})$, so a preconditioner helps
exactly when $\kappa_P < \kappa(H_{\Free\Free})$, and is ideal when $P_\Free =
H_{\Free\Free}$.
\index{preconditioning}

\subsection{Preconditioning on a \emph{changing} free set}

Here the active-set setting imposes a constraint that a general treatment of
preconditioning does not, and it sorts the standard choices sharply. The free set
changes between outer steps, so a preconditioner is cheap only if it \emph{restricts}
to the new free set at negligible cost.

\paragraph{Jacobi restricts exactly.} $P = \diag(H)$ satisfies
$(\diag H)_\Free = \diag(H_\Free)$, so $P_\Free$ is the free-set slice of a single
vector computed once. For $H = M\T M$ that vector is the squared column norms of $M$,
available without forming $H$, and its application is $O(|\Free|)$. It carries no
per-step setup and is valid on every free set the loop visits---the natural default
when $H$ has a widely varying diagonal.

\paragraph{Incomplete Cholesky does not restrict.} In general
$(P_\Free)^{-1} \ne (P^{-1})_\Free$, and the factor of the full $P$ is not a factor of
$P_\Free$, so honouring $H_{\Free\Free}$ would require refactorising on each free set,
reintroducing exactly the assembly the method avoids. Such preconditioners pay off here
only once the free set stabilises---as it does near convergence, and across the
support-stable warm starts of Proposition~\ref{prop:warm}---where one factor serves many
outer steps.

\paragraph{Low-rank-plus-diagonal sits between.} $P = \diag(d) + U\Delta U\T$ built
from a few dominant eigenpairs \cite{saad2003iterative} restricts consistently in its
\emph{application}, $P_\Free = \diag(d)_\Free + U_\Free\Delta U_\Free\T$, and is
inverted matrix-free in $O(|\Free|r^2 + r^3)$ by Woodbury---precisely the direct
free-block solve of Chapter~\ref{ch:operator}. It captures the leading spectral
structure a diagonal misses. And when $H$ is \emph{itself} diagonal-plus-low-rank the
preconditioner \emph{is} $H$, giving $\kappa_P = 1$ and a direct inner solve that
bypasses CG altogether. Randomised Nyström sketching is the practical way to build one,
in a resketch-per-block variant and a sketch-once-and-mask variant that trade setup
cost against reuse.

\begin{remark}[The stopping criterion must not drift]
PCG naturally monitors the \emph{preconditioned} residual norm
$(r\T P_\Free^{-1}r)^{1/2}$, whereas Lemma~\ref{lem:inexact}---the transfer of finite
termination to inexact solves---is stated for the ordinary residual. The two differ by
a factor between $\lambda_{\min}(P_\Free)^{1/2}$ and $\lambda_{\max}(P_\Free)^{1/2}$.
To keep the guarantee intact, evaluate the inner stop on the unpreconditioned residual
(or tighten the preconditioned tolerance by $\kappa(P_\Free)^{1/2}$). The lemma then
applies word for word, because preconditioning changes only how fast the residual falls,
leaving the operator, the residual itself, and the decision margin untouched.

This is a small point with a general moral: when two guarantees are composed, check
that they are stated in the same units.
\end{remark}

\begin{remark}[Preconditioning widens the Krylov advantage]
A scaled gradient step minimises $f$ in the $P$-metric, but the orthant and simplex
projections that make projected gradient cheap are closed-form only in the Euclidean
norm, and become non-negative quadratic programs under a general $P$. So
preconditioning is a lever the Krylov inner solver has that the first-order reference
effectively lacks. It \emph{widens} the $\sqrt\kappa$ gap of
Chapter~\ref{ch:krylov} rather than closing it.
\end{remark}

\section{Linear shrinkage, and conditioning as a by-product}
\label{sec:shrinkage}

Now the statistical instance, which is where the chapter's thesis lives.

The sample covariance $\hat\Sig = X\T X/T$ is a poor estimator of $\Sig$ when $n$ is
not small relative to $T$: its smallest eigenvalues are severely downward-biased and
its condition number far exceeds that of $\Sig$. Ledoit and Wolf \cite{ledoit2004}
propose \emph{linear shrinkage}---replace $\hat\Sig$ by a convex combination with an
SPD target,
\[
  \hat\Sig_{\mathrm{LW}} = (1-\alpha)\,\hat\Sig + \alpha\, R\T R,
\]
which is exactly~\eqref{eq:regsplit}. The intensity $\alpha$ is chosen analytically by
minimising a consistent estimator of the expected \emph{Frobenius} loss; the classical
target is the scaled identity $\bar\lambda I$ with
$\bar\lambda = \tr(\hat\Sig)/n$.
\index{Ledoit--Wolf shrinkage}

Here is the tension, and its resolution.

\paragraph{The statistically optimal intensity barely helps conditioning.} On S\&P~500
data with $n/T \approx 0.41$, the Ledoit--Wolf oracle gives
$\alpha_{\mathrm{LW}} \approx 0.017$, and the Oracle Approximating Shrinkage
estimator \cite{chen2010} gives $\approx 0.010$. Both are small, and the mechanism is
structural rather than accidental. The oracle is $\alpha^\star = \beta^2/\delta^2$ with
$\delta^2 = \norm{\hat\Sig - \bar\lambda I}_F^2$, and $\delta^2$ is \emph{dominated by
the large, well-estimated eigenvalues}: the market factor alone contributes
$(\lambda_1 - \bar\lambda)^2 \approx (0.067)^2$, while each near-zero noise eigenvalue
contributes only $\bar\lambda^2 \approx 2\times10^{-7}$. The near-zero eigenvalues---the
sole cause of $\kappa \approx 109{,}000$---receive negligible correction, because the
Frobenius loss barely notices them. Corollary to Proposition~\ref{prop:weyl}: at
$\alpha^\star \approx 0.017$ the bound gives $\kappa \lesssim 8{,}500$, barely below
the unshrunk value.

\paragraph{The intensity practitioners actually use does help, decisively.}
Practitioners routinely apply $\alpha \approx 0.3$--$0.5$ for out-of-sample stability
\cite{demiguel2009,ledoit2003jpm}. At $\alpha = 0.5$ the bound gives
$\kappa \le \lambda_{\max}(\hat\Sig)/\bar\lambda + 1 \approx 147$, which matches the
empirical value---the bound is near-tight because
$\lambda_{\min}(\hat\Sig) \ll \alpha\bar\lambda$, so
$\lambda_{\min}(\hat\Sig_{\mathrm{LW}}) \approx \alpha\bar\lambda$, precisely the
denominator Weyl's inequality assumes. The measured effect on the solver is
proportionate: CG iterations fall from $684$ to $82$, and the outer active-set loop
contracts from $18$ steps to $6$, because a lower condition number moves the
unconstrained portfolio toward equal weights so fewer assets receive negative weights
in the primal step.

\paragraph{Is there a statistics-versus-speed trade-off?} The Frobenius surrogate and
the condition number are minimised at different $\alpha$, which looks like one. A
rolling-window out-of-sample study says there is none: realised out-of-sample variance
is minimised at $\alpha \approx 0.3$, inside the heavy-shrinkage range that gives the
conditioning benefit, while turnover falls monotonically in $\alpha$. \emph{The
intensity that makes the solver fast is also the intensity a risk-conscious
practitioner would choose}, and the apparent tension is an artefact of the Frobenius
surrogate rather than a property of the problem.

\paragraph{Factor targets.} The scaled identity is not the only choice. A structured
target $T_0 = BB\T + D$, with $B$ holding the loadings of $k \ll n$ risk factors and
$D$ a diagonal idiosyncratic variance, is a drop-in replacement in the stacked-matrix
construction: take $R = (B\T; D^{1/2})$. Because the factor structure clusters
eigenvalues more tightly than a scaled identity, Proposition~\ref{prop:weyl} yields a
smaller bound at the same $\alpha$. The product cost is unchanged, the factor and
diagonal terms costing $O(nk)$ and $O(n)$ against the data term's $O(nT)$.

\section{Spectral cleaning: Marchenko--Pastur}

Shrinkage moves every eigenvalue toward the target. Random-matrix theory offers
something sharper: identify \emph{which} eigenvalues are noise and replace only those.
\index{Marchenko--Pastur law}

Work with the sample \emph{correlation} $C = D^{-1/2}\hat\Sig D^{-1/2}$, where
$D = \diag(\hat\Sig)$. This matters: individual variances are well estimated even in
high dimensions, and the noise that destabilises the problem lives in the correlation
structure. The covariance mixes the well-estimated variances into the spectrum, so its
bulk is not a clean band and its edge has no scale-free form.

The Marchenko--Pastur law \cite{marchenko1967} describes the limiting spectral
distribution of a sample correlation of standardised series when $n, T \to \infty$
with $n/T \to c \in (0,1)$: under the null of no correlation structure the bulk is
supported on $[\lambda_-,\lambda_+]$ with $\lambda_\pm = (1 \pm \sqrt{c})^2$. In
finite samples the empirical upper edge is $\hat\lambda_+ = (1 + \sqrt{n/T})^2$.
Correlation eigenvalues above it carry genuine factor structure; those below are
consistent with noise. Write $k$ for the number above. On S\&P~500 data with $n = 494$,
$T = 1213$: $\hat\lambda_+ \approx 2.28$ and $k = 17$.

\subsection{The constant-residual-eigenvalue estimator}

Let $C = U\Lambda U\T$, partition by the edge into signal $U_k, \Lambda_k$ and noise
$U_0$, and replace the bulk eigenvalues by their average
\[
  \bar\mu = \frac{1}{n-k}\sum_{j>k}\lambda_j = \frac{n - \sum_{j\le k}\lambda_j}{n-k},
\]
which preserves the trace $\tr(C) = n$. The resolution of the identity
$U_kU_k\T + U_0U_0\T = I$ then gives the \emph{constant residual eigenvalue} (CRE), or
eigenvalue-clipping, cleaned correlation \cite{laloux1999,bun2017,lopezdeprado2018}
\begin{equation}
  C_0 = U_k\Lambda_k U_k\T + \bar\mu\,U_0U_0\T = \bar\mu\,I + U_k\Delta_k U_k\T,
  \qquad \Delta_k = \Lambda_k - \bar\mu I,
  \label{eq:cre}
\end{equation}
which is \emph{scalar-identity-plus-low-rank}, with $\delta_j = \lambda_j - \bar\mu > 0$.

\begin{proposition}[Spectral properties of $C_0$]
\label{prop:rmt_spectrum}
$C_0 \succ 0$ with eigenvalues $\{\lambda_1,\dots,\lambda_k,\bar\mu,\dots,\bar\mu\}$
and $\tr(C_0) = n$; $\kappa(C_0) = \lambda_{\max}(C)/\bar\mu$; $C_0$ differs from $C$
only in the noise subspace; and $C_0 v = \bar\mu v + U_k(\Delta_k(U_k\T v))$ costs
$O(nk)$ with $O(nk)$ storage.
\end{proposition}

\begin{proof}
All four from the eigendecomposition
$C_0 = [U_k\ U_0]\diag(\Lambda_k,\bar\mu I)[U_k\ U_0]\T$, the positivity of
$\delta_j$, and the definition of $\bar\mu$.
\end{proof}

Cleaning lifts the noise floor to $\bar\mu$ while preserving the signal, so
$\kappa(C_0) \approx 313$ against $\kappa(C) \approx 91{,}000$: an improvement of
roughly $290\times$, and decisively better than what heavy shrinkage achieves.

\subsection{Two design decisions worth understanding}

\begin{remark}[Why solve in standardised coordinates]
\label{rem:standardised}
Solve the problem in $y = D^{1/2}w$, posing it against $C_0$ rather than rebuilding the
covariance $\Sig_0 = D^{1/2}C_0D^{1/2}$. The reason is arithmetic, not aesthetic: the
rebuilt covariance is \emph{far worse conditioned} ($\kappa(\Sig_0) \approx 1{,}400$ on
the same data), inflated by the spread of asset variances---a factor of about $49$
between largest and smallest. The change of variables is exact, and it also puts the
tolerance on an $O(1)$ scale with noise floor $\bar\mu \approx 0.5$, removing the
dependence on the raw data scale that a covariance formulation carries.
\end{remark}

\begin{remark}[Why clipping, and not a better estimator]
\label{rem:rie}
CRE clipping is the simplest rotationally invariant estimator and is \emph{not}
loss-optimal: nonlinear shrinkage \cite{ledoit2012,ledoit2020} and the oracle RIE
\cite{ledoitpeche2011,bun2017}, which reshape \emph{every} eigenvalue, dominate it in
Frobenius loss. We nevertheless use it for a \emph{structural} reason: clipping is the
only member of the family whose cleaned correlation is
scalar-identity-plus-low-rank---the spiked form \cite{johnstone2001} that admits the
closed-form Woodbury solve of Chapter~\ref{ch:operator}. A full-rank RIE modification
would require a general factorisation and forfeit it.

Whether the statistical gap matters for the application is an empirical question, and
on long-only minimum-variance data it is small. But note the shape of the reasoning:
a slightly worse estimator that is $O(nk)$ to invert may dominate a slightly better one
that is $O(n^3)$, once one is computing hundreds of paths rather than one point.
\emph{Statistical optimality and computational structure are both objectives, and they
are traded, not ranked.}
\end{remark}

Two further notes. A blend $(1-\alpha)C + \alpha C_0$ would re-introduce a fraction of
the bulk eigenvalues, and---the algebraic point---only at $\alpha = 1$ does the
restricted system retain the scalar-identity-plus-low-rank form; for
$\alpha \in (0,1)$ the residual structure is full rank and no $k\times k$ correction
exists. And the solver applies for \emph{any} choice of $k$: a practitioner with an
external risk model, or a PCA with $k$ fixed by cross-validation, can supply the
eigenpairs directly. What the solve requires is not the source of $k$ but the
replacement of the remaining $n-k$ eigenvalues by a single scalar.

\subsection{Is $k$ a delicate hyperparameter?}

No, and it is worth knowing why one can say so. Two biases perturb the threshold:
finite-$T$ Tracy--Widom fluctuations of the largest noise eigenvalue, at scale
$O(T^{-2/3})$ \cite{johnstone2001,elkaroui2008}, and diffuse sector structure just below
the edge. Both move $k$ by at most an eigenvalue or two. An audit at $k-1$, $k$, $k+1$
costs one extra solve per audit point and reveals the sensitivity directly, which is
small over the bulk of a frontier.

That is the right way to treat a threshold in general: rather than argue that it is
correct, measure what changes when it moves by the amount it could plausibly be wrong
by.

\begin{notes}
Ledoit and Wolf's linear shrinkage is \cite{ledoit2004}, with OAS \cite{chen2010} and
nonlinear shrinkage \cite{ledoit2020}. Section~\ref{sec:shrinkage} follows
\cite{schmelzer2026matrixfree}, whose contribution is the identification of the
shrinkage intensity as a preconditioning parameter through the $\alpha$-dependent
bound, and the out-of-sample study reconciling the two objectives. The RMT material is
\cite{schmelzer2026rmt}; Marchenko--Pastur is \cite{marchenko1967}, eigenvalue clipping
\cite{laloux1999}, and the modern treatment \cite{bun2017}. Preconditioning on a
changing free set, including the Nyström constructions, is from
\cite{schmelzer2026nncg}; Benzi \cite{benzi2005} and Saad \cite{saad2003iterative} are
the general references.
\end{notes}

\begin{exercises}

\begin{exercise}
Prove Proposition~\ref{prop:weyl} from Weyl's inequality, and verify the
scaled-identity corollary.
\end{exercise}

\begin{exercise}
Verify that $\tilde M\T\tilde M = H_\alpha$ for the stacked matrix of Section~1, and
explain why this keeps the method matrix-free where forming $H_\alpha$ would not.
\end{exercise}

\begin{exercise}
Show that Jacobi preconditioning restricts exactly, $(\diag H)_\Free = \diag(H_\Free)$,
and that incomplete Cholesky does not, by exhibiting a $3\times3$ example.
\end{exercise}

\begin{exercise}
Given $\lambda_{\max}(\hat\Sig) = 0.07$, $\bar\lambda = 4.6\times10^{-4}$, and
$\lambda_{\min}(\hat\Sig) = 6\times10^{-7}$, evaluate the bound of
Proposition~\ref{prop:weyl} at $\alpha \in \{0.017, 0.1, 0.3, 0.5\}$ and confirm the
numbers quoted in Section~\ref{sec:shrinkage}.
\end{exercise}

\begin{exercise}
Explain, using $\delta^2 = \norm{\hat\Sig - \bar\lambda I}_F^2$, why the Ledoit--Wolf
oracle intensity is insensitive to the near-zero eigenvalues. Construct a spectrum for
which it would not be.
\end{exercise}

\begin{exercise}
Verify Proposition~\ref{prop:rmt_spectrum}, and compute $\bar\mu$ and $\kappa(C_0)$ for
a spectrum with $n = 100$, $k = 3$, $\lambda_{1,2,3} = (20, 8, 4)$ and the remaining
trace spread over the bulk.
\end{exercise}

\begin{exercise}
Show that a blend $(1-\alpha)C + \alpha C_0$ is scalar-identity-plus-low-rank only for
$\alpha \in \{0,1\}$, confirming Section~5.2.
\end{exercise}

\begin{exercise}[Implementation]
Download or simulate a $T\times n$ return matrix with $n/T \approx 0.4$. Compute
$\kappa$ of the sample covariance, of the correlation, of the CRE-cleaned correlation,
and of the rebuilt cleaned covariance. Confirm the ordering claimed in
Remark~\ref{rem:standardised}.
\end{exercise}

\begin{exercise}[Implementation]
Solve the long-only minimum-variance problem with matrix-free CG under
$\alpha \in \{0, 0.017, 0.1, 0.3, 0.5\}$ and record CG iterations and outer active-set
steps for each. Reproduce the qualitative shape of the $684 \to 82$ collapse, and
explain why the outer count falls too.
\end{exercise}

\end{exercises}

\chapter{Warm Starts and Families of Problems}
\label{ch:warmstart}

Solvers are usually studied one problem at a time. Applications almost never present
them that way. An efficient frontier is a family of QPs differing only in a return
tilt. A rolling rebalance is a family differing only in the data. A scenario grid, a
cross-validation sweep, a sensitivity study, a resampled cloud: each is dozens or
hundreds of instances that are nearly the same.

Solved independently, each instance rediscovers an active set it almost always
already had. This chapter is about not doing that. It has three parts: what can be
\emph{reused} exactly, what must be \emph{repaired}, and why neither can change the
answer.

\section{What does not depend on the linear term}

Start with an observation about the invariants of Chapter~\ref{ch:walking} that is
easy to read past.
\begin{equation*}
  J\T G J = I_n, \qquad J\T A = \begin{bmatrix} R \\ 0 \end{bmatrix}.
\end{equation*}
\emph{Neither factor depends on $a$.} $J$ depends on $G$ alone, through
$J\T GJ = I$, and on the working set through the orthogonal transformations
accumulated into it; $R$ depends on $G$ and the working set. The linear term enters
only through $x_u = G^{-1}a$ and the right-hand sides.

So across a family of programs in which $G$, $C$, $b$ are fixed and only $a$ varies,
\emph{both factors are reusable verbatim}, and the entire solution can be recovered
from them.

\begin{proposition}[Recovery from cached factors]
\label{prop:recover}
\index{warm start!recovery}
Let $(J,R)$ satisfy the invariants for a working set $\Active$ of size $k$ with
normals $A$. For a new linear term $a$, set
\begin{equation}
  x_u = J\,(J\T a),
  \qquad
  \rho = b_\Active - A\T x_u,
  \qquad
  R\T y = \rho,
  \qquad
  x = x_u + J_1 y,
  \qquad
  R\,u = y .
  \label{eq:recover}
\end{equation}
Then $(x,u)$ is the solution and multiplier pair of the working-set
subproblem~\eqref{eq:sub}.
\end{proposition}

\begin{proof}
$x_u = JJ\T a = G^{-1}a$ by the first invariant. By
Proposition~\ref{prop:blocks}(i), $u = R^{-1}y = (R\T R)^{-1}\rho = (A\T G^{-1}A)^{-1}\rho$,
the multiplier of~\eqref{eq:subsol}. For the iterate, $Ru = y$ gives
$J_1 y = J_1 R u = G^{-1}Au$ by Proposition~\ref{prop:blocks}(iii), so
$x = x_u + G^{-1}Au$, again as in~\eqref{eq:subsol}.
\end{proof}

\subsection{What the recovery costs, and a correction worth making}

The cost is $O(n^2)$, \emph{all of it in $x_u = J(J\T a)$}: two matrix--vector products
with a dense $n\times n$ matrix, $J$ having lost its triangularity at the first update.
Everything after that is $O(nk + k^2)$. Re-deriving the factors for the same working
set instead costs $k$ insertions at $O(n^2)$ each, so the saving is a factor of order
$k$---largest exactly where it matters, since a long-only optimum is a vertex at which
most bounds bind.

The exponent is easy to get wrong, and the reason is instructive.

\begin{remark}[The natural experiment cannot see the exponent]
\label{rem:exponent}
On a box-constrained family the active set grows with $n$, so $k = \Theta(n)$ and the
predictions $O(nk)$ and $O(n^2)$ describe the \emph{same curve}. Either reading fits
the data, and a paper can report ``linear in $nk$'' without anything looking wrong.

Separating them requires holding $k$ \emph{fixed} while $n$ grows. Doing so, the
measured local exponent $d\log t/d\log n$ runs from well below one at small $n$ to
close to two at large $n$---and neither end is $1$. The flat small-$n$ end is not a
different asymptotic; it is an \emph{overhead floor}, where a solve costs what its
dozen array operations cost to dispatch, and it ends where the $n^2$ arithmetic
overtakes dispatch.

The general lesson is one Chapter~\ref{ch:measurement} develops: \emph{a benchmark
family in which two candidate exponents coincide cannot distinguish them}, and
designing the family that breaks the tie is part of the work.
\end{remark}

There is also an asymmetry in cost between recovering and \emph{checking}. The recovery
is $O(n^2)$, but the certificate must look at every constraint, so it is $O(nm)$ for a
dense $C$. On a dense problem with many constraints, \emph{verification and not
recovery dominates a successful reuse}. That is a slightly deflating fact and a useful
one: it tells you where to spend optimisation effort.

\section{Repair rather than restart}

When the certificate fails, the cached working set is stale---but it remains a far
better starting point than the unconstrained minimum. What blocks resuming from it is
the dual method's invariant~\eqref{eq:invariant}: the iteration requires dual
feasibility, and a stale set typically has one or more negative multipliers. Those are
precisely the constraints that no longer belong.

\begin{proposition}[Repair terminates in a resumable state]
\label{prop:repair}
\index{warm start!repair}
Iterate the following from the cached set: recover $(x,u)$ by~\eqref{eq:recover}; if
every inequality multiplier is non-negative, stop; otherwise drop one constraint with a
negative multiplier, downdate $(J,R)$, and repeat. The loop stops after at most $k$
passes, in a state satisfying~\eqref{eq:invariant}.
\end{proposition}

\begin{proof}
Each pass that does not stop shrinks the working set by one, so there are at most $k$;
with an empty set the sign condition is vacuous, which is the cold start of
Remark~\ref{rem:cold}. On exit $(x,u)$ solves the subproblem by
Proposition~\ref{prop:recover} and the sign conditions hold, which
is~\eqref{eq:invariant}; full column rank is inherited, deleting columns preserving it.
\end{proof}

From such a state \emph{the iteration cannot tell a resumed solve from a cold one}, the
invariant being all it reads, so the resumed walk needs only to drive the remaining
primal infeasibility to zero. In the worst case everything is dropped and the resumed
walk is a cold solve, so repair is never worse than restarting up to the cost of the
drops themselves.

\section{Warm starts in the other two strategies}

The same idea appears in each of the book's strategies, with a different mechanism.

\paragraph{Block pivoting (Chapter~\ref{ch:guessing}).} The warm start is the previous
problem's \emph{free set} as the initial guess. Since the repair rule needs no dual
feasibility to begin---it reads signs and toggles---no repair loop is required at all;
a stale set simply costs an extra exchange or two.

\paragraph{Krylov inner solves (Chapter~\ref{ch:krylov}).} Two distinct warm starts are
available and they are worth distinguishing, because only one is available to a direct
solver. The \emph{outer} warm start is the free set, which cuts the number of outer
steps. The \emph{inner} warm start is the previous solution as CG's initial iterate,
which cuts inner iterations---and by Proposition~\ref{prop:warm} the inner count then
scales with $\log\norm{\delta}$, the logarithm of the \emph{parameter step}, so a finer
grid is cheaper per point. A direct inner solver profits only from the first, having no
analogue of an initial guess. That is where warm-started parametric sweeps obtain their
largest speed-ups.

\paragraph{The homotopy (Chapter~\ref{ch:parametric}).} Here warm starting is not an
optimisation but the algorithm itself: every step is a warm start from the previous
breakpoint. Which is one way of saying that the parametric method is what you get when
you take warm starting to its limit and let the parameter vary continuously.

\begin{remark}[Why active-set methods warm-start and interior-point methods do not]
An active-set method terminates \emph{at} a face of the polyhedron, with the face
recorded as a discrete object---a set of indices. That object is exactly what transfers
to a neighbouring problem. An interior-point method approaches a face asymptotically
from the interior and holds no such object; restarting it near a previous solution puts
the iterate close to the boundary, where the barrier is badly behaved, and the standard
remedy is to back off toward the centre---discarding most of the information one hoped
to reuse. This asymmetry, not raw speed at a single solve, is the main reason
parametric and rolling applications are built on active-set methods.
\end{remark}

\section{Structural continuity: why the active set barely moves}

Warm starting works because consecutive solutions share almost all of their active set.
For a parametric family that is not a hope but a theorem, and it is
Chapter~\ref{ch:parametric}'s.

Along a frontier the solution is piecewise linear in $\lambda$ with generically
\emph{single-index} breakpoint changes. So adjacent solves on a fine grid share all but
a few active constraints, and a warm solver typically terminates in one to two outer
steps. \emph{This is a structural property of parametric quadratic programming, not an
artefact of the data.}

When the family varies in the data rather than in a parameter---a rolling window, a
resampled draw---no such theorem is available, and the empirical picture takes over: the
warm-started outer count tracks the support drift $|\Free \triangle \Free'|$ and
collapses to one whenever the drift is zero. That is the honest statement, and it is
worth keeping the two cases separate rather than claiming the theorem covers both.

\section{Why none of this can change the answer}

Neither reuse nor repair introduces an approximation, and the argument is the
certificate principle of Chapter~\ref{ch:certificate} once more.

Reuse returns a point \emph{only when it carries a certificate}: the working-set
inequalities' multipliers must be non-negative and no constraint outside the working set
may be violated, which by Proposition~\ref{prop:sufficient} is a proof. Repair merely
chooses a different starting state for an iteration whose output is determined
independently of where it started.

What \emph{does} change is the reported iteration count, which counts working-set edits
actually performed and is therefore not comparable with a cold solve's. A benchmark
table mixing warm and cold iteration counts in one column is meaningless, and
Chapter~\ref{ch:measurement} says what to report instead.

\begin{example}[A frontier sweep, end to end]
Fifty points on a long-only efficient frontier for a few hundred assets. Cold, each
point is an independent solve: fifty Cholesky factorisations, fifty walks of tens of
steps. Warm, the first point is a cold solve and each subsequent point recovers from the
cached factors, verifies, and---on the minority of points where the active set has
moved---repairs and resumes for one or two steps. The measured advantage of the warm
sweep over an assemble-and-factor baseline is roughly fivefold, and roughly fifteenfold
against a cold baseline, growing with $n$. Against a general-purpose modelling layer
called once per point, the end-to-end gap is three orders of magnitude---but that figure
mostly measures the modelling layer, and Chapter~\ref{ch:measurement} explains why one
should say so.
\end{example}

\begin{notes}
Propositions~\ref{prop:recover} and~\ref{prop:repair}, and the corrected cost exponent
of Remark~\ref{rem:exponent}, are from the Goldfarb--Idnani note \cite{schmelzer2026quadprog}; the warm-start analysis for
the block-pivoting loop with Krylov inner solves is \cite{schmelzer2026nncg}
(Proposition~\ref{prop:warm}); the frontier-sweep measurements are from
\cite{schmelzer2026rmt} and \cite{schmelzer2026matrixfree}. Parametric warm starting of
active-set QP solvers is treated generally in \cite{nocedal2006}; combining it with a
matrix-free Krylov inner solver on a changing active set is more recent.
\end{notes}

\begin{exercises}

\begin{exercise}
Verify Proposition~\ref{prop:recover} by substituting~\eqref{eq:recover}
into~\eqref{eq:subsol} directly, without using Proposition~\ref{prop:blocks}.
\end{exercise}

\begin{exercise}
Show that the recovery~\eqref{eq:recover} is $O(n^2)$ and identify precisely which two
operations carry the $n^2$. Why does $J$ lose its triangularity after the first update?
\end{exercise}

\begin{exercise}
Prove that the repair loop of Proposition~\ref{prop:repair} preserves full column rank
of the working-set normals, and explain why the corresponding claim would fail for an
\emph{addition} loop.
\end{exercise}

\begin{exercise}
Suppose the cached working set is optimal for the new problem. How many operations does
a successful reuse cost, in terms of $n$, $m$, $k$? Which term dominates for
$m \approx n$, and for $m \approx n^2$?
\end{exercise}

\begin{exercise}
Give a family of problems in which the cached active set is \emph{always} stale, so that
repair drops everything and the warm start is pure overhead. What does this say about
when to attempt reuse?
\end{exercise}

\begin{exercise}
Explain why warm-started and cold iteration counts are not comparable, and propose two
quantities that \emph{are} comparable across the two regimes.
\end{exercise}

\begin{exercise}[Implementation]
Take your solver from Chapter~\ref{ch:walking} and add factor caching plus the recovery
and repair of this chapter. Sweep the linear term over $50$ values and record, per
point: whether reuse succeeded, how many constraints repair dropped, and how many steps
the resumed walk took. Plot all three against the sweep index.
\end{exercise}

\begin{exercise}[Implementation]
Reproduce Remark~\ref{rem:exponent}. Time a warm recovery on a box family where $k$
grows with $n$, fit a power law, and record the exponent. Then hold $k$ fixed at a small
value while $n$ grows and refit. Report both exponents and the $n$ at which the second
family leaves its overhead floor.
\end{exercise}

\end{exercises}


\part{Practice}
\chapter{Numerics: Tolerances, Ties, and Degeneracy}
\label{ch:numerics}

Every chapter so far has ended a proof and then quietly deferred a difficulty. This
chapter collects them.

The material here is usually treated as implementation folklore---the constants a code
carries that no paper explains. It is not folklore. Each of the decisions below has a
reason that can be stated, and stating them is worth a chapter because they are where
correct mathematics turns into incorrect software. A graduate course that stops at the
theorem leaves students to rediscover this layer on their own, generally under
deadline.

\section{The two kinds of tolerance}

The single most useful thing to know about tolerances in this subject is that there are
two kinds and they are not equally delicate.
\index{tolerance!two kinds}

\begin{quote}
\textbf{A steering tolerance} decides what the algorithm tries next: which constraint
enters, which index toggles, which set is guessed. A bad value costs an extra iteration
or a fallback. \emph{It can never cost correctness.}

\textbf{A verdict tolerance} decides what the algorithm returns: the certificate test,
the feasibility check, the infeasibility declaration. It is the only thing standing
between a wrong answer and the caller.
\end{quote}

Almost every tolerance in an active-set code is of the first kind, and the anxiety
practitioners feel about tolerances is largely misdirected onto them. The repair rule
of Chapter~\ref{ch:guessing}, the entering-constraint rule of
Chapter~\ref{ch:walking}, the violator sets of Algorithm~\ref{alg:elim}: all steering.
Set them roughly and move on.

The verdict tolerances deserve the attention---and even they are less delicate than they
look, for the reason given in Chapter~\ref{ch:certificate}: the quantity being
thresholded is \emph{bimodal}. A candidate built from a well-conditioned working set
satisfies the KKT test to a few multiples of machine epsilon; one built on a set that
is not has no reason to come close. The two populations sit orders of magnitude apart
and the threshold lies in empty space between them. The same is true of the
infeasibility verdict of Proposition~\ref{prop:farkas}: infeasibility is
\emph{macroscopic}, its size set by the geometry of the constraints rather than by the
arithmetic, so separating it from rounding is easy.

\emph{When a threshold sits in a gap, its exact value does not matter; when it sits in a
continuum, no value is right.} The useful engineering question is therefore not ``what
tolerance?'' but ``is this quantity bimodal, and if not, can I make it so?''

\begin{example}[Making a threshold non-delicate by choice of coordinates]
Chapter~\ref{ch:conditioning} solved the RMT problem in standardised coordinates
$y = D^{1/2}w$ partly for conditioning. There is a second benefit. In those coordinates
both the primal quantity (a standardised weight) and the dual quantity (a reduced
gradient) live on the $O(1)$ correlation scale, with noise floor $\bar\mu \approx 0.5$,
so a default $\varepsilon = 10^{-6}$ is a standard KKT tolerance with no data-dependent
scaling. The change of variables removes the dependence on the raw return scale that a
covariance formulation would carry. Choosing coordinates so that the tolerances become
absolute is a real design move.
\end{example}

\section{Scale invariance}

A tolerance is meaningless without a scale, and the cheapest way to get one is to make
the algorithm invariant instead.

Chapter~\ref{ch:walking} gave the canonical instance: dividing the violation by
$\norm{c_i}$ in the entering rule~\eqref{eq:select} makes the selection invariant under
$(c_i,b_i)\mapsto(\gamma c_i,\gamma b_i)$ (Proposition~\ref{prop:scale}). Without it, a
constraint written in basis points rather than in units is selected first merely for
having been written differently. The proposition claims only that the first
\emph{choice} is unchanged; measured over three decades of $\gamma$ either side of
unity, the observed behaviour is stronger---an identical number of additions and
removals, so the whole trajectory is unchanged, with the returned minimiser agreeing to
about $3\times10^{-16}$.

The general rule: \emph{before choosing a tolerance for a quantity, ask what the
quantity's units are, and whether the algorithm can be rewritten so that it has none}.

\section{Slope floors and window guards}

Two guards recur in path-tracing codes and they are guarding different things.

\paragraph{The slope floor.} In the event rules of Chapter~\ref{ch:parametric} every
candidate step is a ratio $t = (\text{gap})/(\text{slope})$. A slope of exactly zero
means the event cannot occur; a slope of $10^{-17}$ means it occurs at
$t \approx 10^{17}$, which is noise amplified. But a \emph{tiny but genuine} slope still
crosses a bound given a long enough parameter range, so filtering at the user tolerance
would let variables drift out of bounds.

The resolution is to test slopes against $\sqrt{\varepsilon_{\mathrm{mach}}}$---not
against the user tolerance---discarding only slopes at floating-point noise level. And
the exact level is not delicate, for the bimodality reason again: genuine slopes sit
orders of magnitude above $\sqrt{\varepsilon_{\mathrm{mach}}}$ and round-off noise
orders below, so nothing lands near the threshold.

\paragraph{The window guard.} The segment is valid only for parameter values at or below
the current one, because the path is traced monotonically. Candidate ratios above the
current value are spurious crossings outside the segment and must be discarded. This is
not a numerical guard at all; it is the algorithm's contract, and a code that omits it
will occasionally step backwards and then loop.

\paragraph{Relative versus absolute in the ratio test.} As Chapter~\ref{ch:reading}
noted, the breakpoints of a large path crowd together. The tolerance that orders them
must therefore be \emph{relative} to the scale of the parameter. With it the path comes
back strictly monotone to floating-point precision; with a fixed absolute tolerance,
distinct breakpoints eventually merge and the parameter steps backwards. Same lesson as
Section~2: get the units right first.

\section{Ties: anti-cycling in practice}

On degenerate problems---tied expected returns, duplicated columns, several constraints
active at once---multiple events occur at the same parameter value. A naive ``pick any
maximiser'' rule can cycle.
\index{anti-cycling}

The remedy throughout this book is a \textbf{Bland-style rule}: among all events within
tolerance of the best ratio, select the one with the lowest index, and within an index
the lowest event-type. This makes the choice deterministic and resolves degeneracies one
index per iteration. Termination then follows by the standard argument: finitely many
active-set states, a monotone parameter, and the lowest-index rule excluding the only
failure mode, which is revisiting a state without progress.

Three practical notes.

First, \emph{determinism is itself the point}, quite apart from termination. A solver
that makes an arbitrary choice among tied events is not reproducible across platforms,
BLAS versions, or compiler flags, and debugging it is miserable.

Second, the rule must be applied to events \emph{within a tolerance} of the best ratio,
not to exact ties, since exact ties are not what one observes.

Third, an iteration cap is a \emph{numerical safety net}, not the termination
guarantee. A correct trace empirically runs $O(n)$ iterations, so a cap of, say,
$100(n+p+1)$ is never approached; exceeding it signals a floating-point failure such as
a tolerance-induced cycle, which should be raised loudly rather than allowed to hang.
Do not confuse the cap with the theorem.

\section{Cancellation, and rewriting differences as quotients}

Remark~\ref{rem:sign} of Chapter~\ref{ch:freeblock} gave the pattern in its cleanest
form. A Householder reflection needs $v_1 = (d_2)_1 - \alpha$ with
$\alpha = \pm\norm{d_2}$; choosing the sign that makes this a difference of nearly equal
quantities loses digits catastrophically when $d_2$ is close to a positive multiple of
$e_1$. The cure is algebraic:
\[
  v_1 = \frac{(d_2)_1^2 - \nu^2}{(d_2)_1 + \sign((d_2)_1)\nu}
      = \frac{-\sign((d_2)_1)\,\tau}{|(d_2)_1| + \nu},
  \qquad \tau = \norm{(d_2)_{2:}}^2, \ \nu = \norm{d_2},
\]
in which every operation combines like-signed quantities.

Two observations generalise. First, \emph{a subtraction of nearly equal quantities can
usually be rewritten as a quotient} by multiplying by the conjugate---the same device
that fixes the quadratic formula and $\sqrt{x+1}-\sqrt{x}$. Look for it whenever a
formula contains a difference whose operands you expect to be close.

Second, note the interaction with Proposition~\ref{prop:invariance}: the sign convention
is \emph{immaterial to the iterates}, so one is free to choose it. The rewriting is what
makes it immaterial to the \emph{accuracy} as well. Mathematical freedom and numerical
freedom are not the same freedom, and having established the first does not discharge
the second.

\section{When the iterate leaves the feasible set}

Here is a failure mode that is neither a conditioning failure nor a rank failure in the
usual sense, and that earlier versions of real codes aborted on.

On a near-rank-deficient quadratic form the smallest eigenvalues span near-flat
directions of the objective. Accumulated round-off over the hundreds of steps of a large
trace nudges a free variable a hair---typically $10^{-5}$ to $10^{-3}$---outside its box
along one of those directions. The candidate is therefore \emph{optimal to solver
precision but not exactly feasible}, and a strict feasibility check rejects it.

Because the deviation lies in the flat directions, projecting the candidate back onto
the feasible box, and rescaling to restore the budget, changes the objective only at the
level of the round-off itself. Doing so lets the trace complete with points that remain
optimal: across a controlled rank-deficiency sweep every completed point matches a
reference QP solve in objective to about $10^{-8}$, and a well-posed trace is unchanged
to the last digit, the projection being a no-op there.

\paragraph{The repair must be bounded, and by the right quantity.} The two regimes are
told apart by the \emph{conditioning of the free block}, not by the size of the violation
it happens to produce. While that block is numerically full rank its solve is reliable
and any infeasibility is round-off, which is projected away; once the free set outgrows
the rank of the quadratic form the block is numerically singular and its solve is
unreliable, so projecting whatever it returns could silently yield a suboptimal answer.
A guard therefore declines as soon as the free block's reciprocal condition number falls
below about $10^{-12}$, raising an actionable diagnosis instead of a plausible number.

Why key on conditioning rather than on the violation? \emph{Because the violation is the
residual of a singular solve, and its magnitude varies with the BLAS and LAPACK build,
whereas the conditioning is read from the block's symmetric eigenvalues and is identical
on every platform.} Guard on the reproducible quantity.

This gives a precise guarantee, and it is worth stating in the form a user can rely on:
while the free block stays numerically full rank the trace completes and every point is
optimal; when it does not, the algorithm \emph{declines rather than mislead}. Note also
what the guarantee does not cover---it speaks to the linear algebra, not to event ties,
which are the separate combinatorial degeneracy of Section~4.

\section{A short checklist}

For a reader who will implement one of these methods:

\begin{enumerate}
\item Classify every tolerance as steering or verdict. Spend your attention on the
      verdicts.
\item For each verdict tolerance, plot the distribution of the quantity it thresholds
      over a realistic test set. If it is bimodal, you are done. If it is not, change
      coordinates or change the test.
\item Make the algorithm scale-invariant wherever a normalisation can do it, so that
      fewer tolerances need scales at all.
\item Use a lowest-index tie-break, applied within tolerance, and make the solver
      deterministic.
\item Carry an iteration cap, and treat hitting it as a bug report rather than an
      answer.
\item Check the certificate against the original data, never against the construction
      that produced the candidate.
\item Guard on conditioning, which is reproducible, rather than on residual magnitudes,
      which are not.
\item Decline loudly rather than return quietly when the guard fires.
\end{enumerate}

\begin{notes}
The two-tolerance distinction and the bimodality argument are from the Goldfarb--Idnani note \cite{schmelzer2026quadprog};
the slope floor, window guard, relative ratio tolerance, and feasible-set projection
with its conditioning guard are from the \texttt{cvxcla} software paper \cite{cvxcla};
the standardised-coordinate tolerance argument is \cite{schmelzer2026rmt}. Higham
\cite{higham2002} is the reference for cancellation and for accuracy analysis
generally; Golub and Van Loan \cite{golub2013} for the transformations themselves.
\end{notes}

\begin{exercises}

\begin{exercise}
Classify every tolerance appearing in Chapters~\ref{ch:walking}--\ref{ch:parametric} as
steering or verdict, and justify each classification in one sentence.
\end{exercise}

\begin{exercise}
Show that $\sqrt{x+1}-\sqrt{x}$ loses roughly half the available digits for large $x$,
and rewrite it as a quotient. Relate the rewriting to~\eqref{eq:nocancel}.
\end{exercise}

\begin{exercise}
Explain why $\sqrt{\varepsilon_{\mathrm{mach}}}$ rather than
$\varepsilon_{\mathrm{mach}}$ is the right scale for a slope floor. (Hint: consider the
error in a computed slope that is the difference of two $O(1)$ quantities.)
\end{exercise}

\begin{exercise}
Construct a two-variable example in which a fixed absolute tolerance in the ratio test
merges two distinct breakpoints, and show that a relative tolerance separates them.
\end{exercise}

\begin{exercise}
Give an example of a solver decision that is deterministic in exact arithmetic but
platform-dependent in floating point, and propose a rule that removes the dependence.
\end{exercise}

\begin{exercise}
Argue that projecting a slightly infeasible iterate back onto its box changes the
objective by $O(\epsilon^2\lambda_{\min})$ when the deviation $\epsilon$ lies in a flat
direction, and explain why this justifies the repair of Section~6.
\end{exercise}

\begin{exercise}[Implementation]
Instrument your solver from Chapter~\ref{ch:walking} to record, over a few hundred
random instances, the maximum KKT residual of every accepted candidate. Plot the
histogram on a log scale and confirm the bimodality claimed in
Chapter~\ref{ch:certificate}. Then deliberately feed it a rank-deficient working set and
see where the second mode lands.
\end{exercise}

\begin{exercise}[Implementation]
Build a controlled rank-deficiency sweep: generate covariance matrices of decreasing
numerical rank, trace a frontier for each, and record for every point whether the free
block passed the conditioning guard, whether the projection fired, and the objective gap
against a reference QP solve. Reproduce the two regimes of Section~6.
\end{exercise}

\end{exercises}

\chapter{Measuring What Matters}
\label{ch:measurement}

The last chapter of a book on algorithms should be about how to tell whether one is
faster than another, because that turns out to be harder than either the mathematics or
the implementation.

The chapter has two halves. The first is methodological: which regime you are in
decides what to optimise, and most benchmark tables in this literature do not say which
regime they measured. The second is a worked case study---the long-only portfolio, end
to end---which puts every part of the book together and reports the numbers honestly,
including the ones that make the method look less impressive than a headline figure
would.

\section{Two regimes}

\begin{quote}
Below a few hundred variables, a solve costs what its array operations cost to
\emph{dispatch}, not what its arithmetic costs. Above a few thousand, the reverse.
\end{quote}

This is not a subtlety; it inverts the design objective.

In the dispatch-bound regime, halving the flop count buys nothing and halving the
\emph{iteration} count buys everything. A method that performs more arithmetic in fewer
steps wins. This is why the guessing strategy of Chapter~\ref{ch:guessing} beats the
exact walk of Chapter~\ref{ch:walking} at moderate $n$ by \emph{more} than the gap
between two implementations of that walk in different languages---a striking measurement,
because it says the algorithmic choice dominates the language choice, which is the
opposite of the usual folklore.

In the arithmetic-bound regime, the opposite: iteration counts matter only through their
product with per-step cost, and structure in the operator (Chapter~\ref{ch:operator})
is what pays.

\emph{Knowing which regime you are in is the whole of the decision.} Effort spent on
kernels is wasted where the bottleneck is the interpreter, and effort spent on iteration
counts is wasted where it is not. So the first thing a benchmark should report is not a
speedup but a regime---and the cheapest diagnostic is to plot the local exponent
$d\log t/d\log n$ and look for where the curve leaves its flat overhead floor.

\begin{remark}[The overhead floor is not an asymptotic]
Remark~\ref{rem:exponent} of Chapter~\ref{ch:warmstart} made this point about a specific
measurement and it is worth generalising. A timing curve that is flat at small $n$ is
not exhibiting a better complexity there; it is exhibiting a constant. Fitting a single
power law across both regimes produces an exponent that describes neither. Fit them
separately, and say where you cut.
\end{remark}

\section{Designing a benchmark family that can distinguish}

Two failure modes recur, and both are avoidable.

\paragraph{The family that hides the exponent.} If two candidate complexities coincide
on your test family, no amount of data distinguishes them. On a box-constrained family
the active set grows with $n$, so $O(nk)$ and $O(n^2)$ predict the same curve; holding
$k$ fixed while $n$ grows separates them. Before running a scaling study, write down the
two hypotheses and check that the family can tell them apart.

\paragraph{The family that hides the hazard.} Chapter~\ref{ch:guessing} gave the sharp
case. Over six hundred random instances across four well-posed constraint families, the
rank failure that destroys the finite-termination guarantee \emph{never occurs};
duplicating columns of the constraint matrix raises it to $93\%$. A literature testing
only on the first kind of family would reasonably conclude the hazard is theoretical.

The general point is uncomfortable and worth stating plainly: \emph{test families in a
literature tend to share the structure that makes the literature's methods look good.}
Finding out which structure that is, is part of reading a paper. The synthetic families
common in the QP literature produce small active sets; a long-only portfolio produces a
huge one---on S\&P~500 data roughly $450$ of $495$ constraints active. Those exercise
entirely different code paths, and a method tuned on one may be untested on the other.

\section{What to compare against}

Three levels of comparison give three different numbers, and conflating them is the most
common overstatement in applied optimization papers.

\begin{enumerate}
\item \textbf{Same family, different linear algebra.} The fair algorithmic comparison.
      A dual active-set method on $(J,R)$ factors against a dual active-set method on
      recursive $LDL\T$ updates \cite{arnstrom2022daqp} isolates the linear algebra and
      nothing else.
\item \textbf{Different algorithm, native API.} A Krylov active-set method against an
      interior-point solver \cite{clarabel2024} or an operator-splitting solver
      \cite{osqp}, each called through its own interface. This measures algorithms, and
      it is the number a paper's abstract should carry.
\item \textbf{Different algorithm, through a modelling layer.} The same comparison with
      the competitor wrapped in a general-purpose modelling front end. This measures
      algorithm \emph{plus} problem-construction overhead, and it is an \emph{end-to-end}
      figure, not an algorithmic one.
\end{enumerate}

The gap between levels 2 and 3 is large and systematic. On the portfolio study below,
the algorithmic gain over an interior-point solver is roughly $5.7\times$ without
shrinkage and $22\times$ under heavy shrinkage; the corresponding end-to-end figures
against the same solver behind a modelling layer are $156\times$ and $661\times$. Both
sets of numbers are true. Only the first pair is about the algorithm, and a paper that
reports only the second pair is not lying but is not informing either.

\emph{Report the fair comparison in the abstract and the end-to-end one in the text,
labelled as such.}

\section{Case study: the long-only portfolio, end to end}

Everything in this book meets here.

\paragraph{The problem.} Minimise $w\T\hat\Sig w - \rho\hat\mu\T w$ subject to
$Bw = c$ and $w \ge 0$, with $\hat\Sig$ estimated from a $T \times n$ matrix of
demeaned returns. The data: $494$ S\&P~500 constituents over five years, $T = 1213$,
so $n/T \approx 0.41$.

\paragraph{Step 1: do not form the covariance.} The variance term is
$\norm{Xw}^2/T$, so the gradient is $w \mapsto X\T(Xw)$: two products, $O(nT)$, no
$n \times n$ storage (Chapter~\ref{ch:operator}, Gram backend). The covariance matrix is
not a primitive of the problem; it is a computational detour.

\paragraph{Step 2: eliminate the balance multipliers.} The $p$ multipliers of $Bw = c$
come out analytically through a $p\times p$ Schur complement, leaving $p+1$ SPD solves
(Chapter~\ref{ch:krylov}, Section~3). For the budget constraint alone this is
$w = v_1/(\onevec\T v_1)$ with $v_1 = \hat\Sig^{-1}\onevec$: a \emph{single} solve, and a
formula visible in Merton \cite{merton1972} sixty years ago but not, until recently,
exploited computationally.

\paragraph{Step 3: enforce non-negativity by an active-set loop.} The primal--dual loop
of Chapter~\ref{ch:guessing} wraps the solver unchanged, with the finite-termination
guarantee of Theorem~\ref{thm:termination} and the inexactness transfer of
Lemma~\ref{lem:inexact}.

\paragraph{Step 4: condition it.} Ledoit--Wolf shrinkage at the intensity practitioners
already use, or RMT cleaning (Chapter~\ref{ch:conditioning}).

\paragraph{Step 5: for a frontier, warm start.} Chapter~\ref{ch:warmstart}.

\subsection{What is measured}

\begin{center}
\footnotesize
\begin{tabular}{@{}lll@{}}
\toprule
Quantity & Without shrinkage & Heavy shrinkage ($\alpha = 0.5$) \\
\midrule
$\kappa(\hat\Sig)$                    & $\approx 109{,}000$ & $\approx 147$ \\
CG iterations                         & $684$   & $82$ \\
Outer active-set steps                & $18$    & $6$ \\
Speedup vs.\ interior point (native)  & $5.7\times$ & $22\times$ \\
Speedup vs.\ interior point (modelling layer) & $156\times$ & $661\times$ \\
\bottomrule
\end{tabular}
\end{center}

Four readings, in decreasing order of how flattering they are.

\emph{The conditioning effect is real and compound.} Shrinkage cuts CG iterations by a
factor of eight, as Proposition~\ref{prop:cg} predicts from the $\kappa$ collapse. It
also cuts the \emph{outer} count from $18$ to $6$, which is not a $\kappa$ effect at
all: a lower condition number moves the unconstrained portfolio toward equal weights, so
fewer assets receive negative weights in the first primal step. Two independent
mechanisms, one parameter.

\emph{Scaling favours the method.} At $n = 3000$ on synthetic factor-model data with $T$
fixed, the matrix-free method takes about $0.081$ seconds and its runtime grows
substantially \emph{sub}quadratically, while interior-point methods scale as $O(n^3)$
per iteration. The gap widens with $n$.

\emph{A well-engineered general solver is still slower, and this is the informative
comparison.} An operator-splitting solver called through its \emph{direct} API---no
modelling overhead---remains $6$ to $10\times$ slower. That is the honest statement of
what exploiting structure buys: not three orders of magnitude, but roughly one, and
reliably.

\emph{The three-order-of-magnitude figures are mostly about something else.} A $50$-point
frontier for $n = 500$ completes about $1{,}455\times$ faster than a naive loop over a
modelling layer. That number is real and it is what a user experiences, but most of it
is problem-construction overhead paid $50$ times, not algorithmic superiority. Say so.

\subsection{The statistical objective, checked separately}

A faster answer to the wrong problem is worthless, so the conditioning choice has to be
validated on its own terms. A rolling-window out-of-sample study on the same universe
finds realised out-of-sample variance minimised at $\alpha \approx 0.3$---inside the
heavy-shrinkage range that gives the conditioning benefit---with turnover falling
monotonically in $\alpha$. And RMT-cleaned estimates carry no measurable statistical
penalty relative to standard shrinkage on these data.

So the statistical and numerical objectives are reconciled at the intensities
practitioners already use, and the apparent tension noted in
Chapter~\ref{ch:conditioning} is an artefact of the Frobenius surrogate.

\emph{Notice what this required: a separate experiment, with a separate metric, that
could have come out the other way.} A computational paper that only measures time has
not established that its method should be used.

\section{Reporting}

A checklist for a benchmark section, derived from the above.

\begin{enumerate}
\item State the regime. Report where the overhead floor ends.
\item Report the \emph{fair} comparison---same family, or native API---as the headline,
      and label end-to-end figures as end-to-end.
\item Report iteration counts and wall time separately, and never mix warm and cold
      iteration counts in one column (Chapter~\ref{ch:warmstart}).
\item Say what the test family's active-set size is, since methods differ enormously
      between small and large active sets.
\item Design the family so that competing complexity hypotheses predict different
      curves, and say which hypotheses you were testing.
\item Verify correctness independently of timing: every returned point should pass the
      certificate of Chapter~\ref{ch:certificate} against the original data, and agree
      with a reference solver.
\item Where a method wins by an approximation---a cleaned estimator, a regularised
      operator---validate that approximation on its own metric.
\item Report what did not work, and on what family the hazard you did not hit would have
      fired.
\end{enumerate}

\section{Closing}

The book began with the claim that the analysis of the convex quadratic program is easy
and the combinatorics hard, and that every method is a strategy for searching among
active sets. Four strategies later, the strategies have turned out to share more than
they differ: one certificate makes all of them safe; one operator interface serves all
of them; one linear-algebra kernel---a free block, bordered by active normals, updated
rather than rebuilt---sits inside all of them; and one parameter, allowed to vary, turns
the whole subject into a single curve that four literatures traced independently.

If there is one thing to carry away, it is the certificate principle of
Chapter~\ref{ch:certificate}, because it generalises far past quadratic programming: a
method that may fail, but that can recognise its own success, needs a certificate and
not a convergence theory. Ask of any algorithm you meet whether it can check its own
answer. If it can, you may be as inventive as you like about how it finds one.

\begin{notes}
The regime observation and the measurement that the fast path beats the exact walk by
more than the gap between two language implementations are from the Goldfarb--Idnani note \cite{schmelzer2026quadprog}. The
portfolio measurements are from \cite{schmelzer2026matrixfree}; the frontier-sweep and
RMT figures from \cite{schmelzer2026rmt}; the out-of-sample study from
\cite{schmelzer2026matrixfree}, with \cite{demiguel2009} for the wider evidence on
shrinkage and out-of-sample performance. The solvers benchmarked against are Clarabel
\cite{clarabel2024} and OSQP \cite{osqp}, with CVXPY \cite{diamond2016} the modelling
layer.
\end{notes}

\begin{exercises}

\begin{exercise}
Given timings $t(n)$ at $n = 50, 100, 200, 400, 800, 1600$, describe how to estimate the
local exponent $d\log t/d\log n$ and how to locate the end of the overhead floor.
\end{exercise}

\begin{exercise}
Explain why a single power-law fit across both regimes yields an exponent that describes
neither, and construct synthetic timings for which the fitted exponent is $1.0$ although
neither regime has that exponent.
\end{exercise}

\begin{exercise}
For each of the three comparison levels of Section~3, name a question it answers and a
question it cannot.
\end{exercise}

\begin{exercise}
The outer active-set count fell from $18$ to $6$ under shrinkage. Give the mechanism, and
explain why it is \emph{not} a consequence of Proposition~\ref{prop:cg}.
\end{exercise}

\begin{exercise}
Suppose a method is $1{,}455\times$ faster than a baseline end to end and $22\times$
faster algorithmically. What is the remaining factor measuring, and how would you
measure it directly?
\end{exercise}

\begin{exercise}
Design a test family on which the block-pivoting rank hazard of
Chapter~\ref{ch:guessing} fires, and one on which it never does. What does the existence
of both say about published benchmark suites?
\end{exercise}

\begin{exercise}[Implementation]
Assemble the full pipeline of Section~4 on real or simulated returns: Gram operator,
Schur elimination, active-set loop, shrinkage, warm-started frontier sweep. Verify every
returned portfolio against the certificate of Chapter~\ref{ch:certificate} and against a
general-purpose QP solver. Then reproduce the table of Section~4.1 on your own data and
comment on every figure that differs.
\end{exercise}

\begin{exercise}[Implementation]
Run a rolling-window out-of-sample study over shrinkage intensities
$\alpha \in \{0, 0.017, 0.1, 0.3, 0.5, 0.7\}$, recording realised out-of-sample variance
and turnover at each. Does your data reproduce the reconciliation of Section~4.2? If not,
what differs?
\end{exercise}

\end{exercises}


\appendix
\chapter{Notation and the Rosetta Stone}
\label{app:notation}

This book fixes one notation across six research papers and four literatures that did
not share one. This appendix records the choices and translates them.

\section{Notation used in this book}

\begin{center}
\footnotesize
\begin{tabular}{@{}lll@{}}
\toprule
Symbol & Meaning & First used \\
\midrule
$G$, $H$ & the quadratic form (Hessian); $G$ in Parts I--II, $H$ in Part III
         & Ch.~\ref{ch:oneproblem}, Ch.~\ref{ch:parametric} \\
$a$, $d$ & the linear term; $-a\T x$ in Parts I--II, $-\lambda d\T x$ in Part III
         & Ch.~\ref{ch:oneproblem}, Ch.~\ref{ch:parametric} \\
$C$, $b$ & constraint normals (columns) and right-hand sides, $C\T x \ge b$
         & Ch.~\ref{ch:oneproblem} \\
$m_{\mathrm{eq}}$ & number of leading constraints treated as equalities
         & Ch.~\ref{ch:oneproblem} \\
$A$ & in Parts I--II, the working-set normals $C_{\cdot\Active}$;
      in Part III, the equality system & Ch.~\ref{ch:oneproblem} \\
$\Active$ & working set (algorithm state) or active set (property of a point)
         & Def.~\ref{def:activeset} \\
$\Free$, $\Bound$ & free and blocked (bound) index sets
         & Ch.~\ref{ch:geometry} \\
$\lambda$, $u$, $\nu$, $\zeta$ & multipliers; $\lambda$ full-length, $u$ working-set,
      $\nu,\zeta$ in Part III & Ch.~\ref{ch:certificate} \\
$x_u = G^{-1}a$ & the unconstrained minimiser & Prop.~\ref{prop:sub} \\
$A\T G^{-1}A$ & the dual Hessian of a working set & Def.~\ref{def:dualhessian} \\
$J$, $R$ & the carried factors, $J\T GJ = I$, $J\T A = [R;0]$
         & Ch.~\ref{ch:walking} \\
$z$, $r$ & the primal and dual step directions & \eqref{eq:dzr} \\
$s$ & slacks (Part II) or the reduced gradient $Ax - b$ (bound-constrained case)
         & Ch.~\ref{ch:walking}, Ch.~\ref{ch:geometry} \\
$\kappa$ & spectral condition number $\lambda_{\max}/\lambda_{\min}$
         & Ch.~\ref{ch:freeblock} \\
$\Sig$, $\mu$, $w$ & covariance, expected returns, portfolio weights
         & Ch.~\ref{ch:oneproblem} \\
$X$, $y$, $\beta$ & design, response, regression coefficients
         & Ch.~\ref{ch:homotopy} \\
\bottomrule
\end{tabular}
\end{center}

Two conventions inherited from the reference implementations and used throughout
Parts~I and~II: the linear term is \emph{subtracted}, so the unconstrained minimiser is
$G^{-1}a$ rather than $-G^{-1}a$; and constraints are held \emph{column-wise}, so the
system is $C\T x \ge b$ rather than $Cx \le b$.

One genuine collision the reader should watch for. In Parts~I and~II, $A$ is the matrix
of working-set normals. In Part~III, following the parametric literature, $A$ is the
equality-constraint system and the working set is carried by the partition
$(\Free,\Bound,\mathcal{S})$. The chapters say which is in force; the collision is
retained rather than resolved because both usages are standard in their own literature
and a reader going on to the primary sources will meet both.

\section{The Rosetta stone}

Each row is one component of the common scheme of Chapter~\ref{ch:parametric}; the
columns give the name it goes by in the four literatures that developed it
independently. The leftmost column is the neutral optimization term of which all four
are special cases.
\index{Rosetta stone}

\begin{center}
\footnotesize
\begin{tabular}{@{}p{0.15\textwidth}p{0.19\textwidth}p{0.19\textwidth}p{0.17\textwidth}p{0.19\textwidth}@{}}
\toprule
Common scheme & Portfolio selection (CLA) & Regression (LASSO / LARS) &
Kernel methods (SVM path) & Control (explicit MPC) \\
\midrule
Homotopy parameter $\lambda$ & risk--return trade-off & penalty $\lambda$ &
cost $C$ & state $x$ (a \emph{vector}) \\[2pt]
Quadratic form $H$ (as operator) & covariance $\Sig$ & Gram matrix $X\T X$ &
kernel Gram $\kappa(X,X)$ & state--input cost \\[2pt]
Active set & assets held off their bounds (the free set) &
the support (active coefficients) & margin support vectors &
constraints active on a critical region \\[2pt]
Face mechanism & imposed by box bounds & induced by the $\ell_1$ subgradient &
induced by the margin & imposed by state/input constraints \\[2pt]
Breakpoint (event) & turning point / corner portfolio & knot / LARS event &
a point crossing the margin & critical-region boundary \\[2pt]
Event selection (ratio test) & next asset to enter or leave the free set &
next variable to enter or leave & next point to cross the margin &
next region boundary reached \\[2pt]
Solution path & efficient frontier & regularization path & regularization path &
explicit (piecewise-affine) control law \\
\bottomrule
\end{tabular}
\end{center}

The control column describes the \emph{multi-parameter} generalisation of
Remark~\ref{rem:vector}: the interval of breakpoints becomes a polyhedral partition into
critical regions and the homotopy becomes region traversal. The remaining entries are
its scalar slice.

\section{The four events, by dialect}

Chapter~\ref{ch:parametric} named four event families and Table~\ref{tab:events} read
two dialects. For reference, the full sorting:

\begin{center}
\footnotesize
\begin{tabular}{@{}llll@{}}
\toprule
Method & Events populated & Active set & Face mechanism \\
\midrule
Forward LARS               & $\mathrm{D}_1$
                           & monotone      & induced \\
LASSO homotopy             & $\mathrm{P}_1$, $\mathrm{D}_1$
                           & non-monotone  & induced \\
Critical Line Algorithm    & $\mathrm{P}_1$, $\mathrm{D}_1$
                           & non-monotone  & imposed (box) \\
CLA with coupling rows     & all four
                           & non-monotone  & imposed \\
Constrained LASSO          & all four
                           & non-monotone  & mixed \\
SVM path                   & $\mathrm{P}_1$, $\mathrm{D}_1$
                           & non-monotone  & induced (margin) \\
Explicit MPC               & all four, over a vector parameter
                           & non-monotone  & imposed \\
\bottomrule
\end{tabular}
\end{center}

The Critical Line Algorithm and the LASSO homotopy populate the same two families. The
difference between them is that the CLA's feasibility intervals are \emph{fixed} (a box)
while the LASSO's travel with $\lambda$ (the interval $[-\lambda,\lambda]$). That single
difference is all that separates an imposed active set from an induced one at the level
of the algorithm, and Theorem~\ref{thm:identity} is what makes the two interconvertible.

\chapter{A Short History}
\label{app:history}

The methods in this book were invented four times, in four fields, over fifty years,
largely without cross-reference. This appendix tells that story, because knowing it
changes how one reads the literature: a graduate student who meets least angle
regression as a statistical innovation and the Critical Line Algorithm as a piece of
finance folklore will not think to look for the third and fourth copies.

The account is a sketch, not a survey. Each of the four literatures has a history deep
enough to fill an article of its own.

\section{The 1950s: optimization and finance}

The active-set treatment of the parametric quadratic program is a creation of the
1950s, and the Critical Line Algorithm is its earliest fully worked instance.

The decade was the founding one for operations research, and its founding algorithm was
Dantzig's simplex method \cite{dantzig1951}, whose vertex-to-vertex walk along the edges
of a polyhedron the homotopy of Chapter~\ref{ch:parametric} still rides---the ratio test
and the Bland tie-break of that chapter are literally the simplex method's. The same
decade supplied the optimality conditions on which everything in this book rests:
Kuhn and Tucker \cite{kuhntucker1951}, anticipated by Karush \cite{karush1939},
characterised the constrained optimum by the stationarity and complementary-slackness
relations whose parametric solution is the bordered system of
Lemma~\ref{lem:affine}.

Markowitz worked alongside Dantzig at the RAND Corporation and took the method as his
starting point. Having posed mean--variance selection \cite{markowitz1952}, he
recognised that the efficient set is obtained by solving a quadratic program
\emph{parametrically} in the risk--return trade-off, and \cite{markowitz1956} gave the
Critical Line Algorithm, developed in full in the 1959 monograph \cite{markowitz1959}.

The idea did not appear in a vacuum. Gass and Saaty \cite{gasssaaty1955} introduced
\emph{parametric} linear programming, the direct linear ancestor: follow the optimal
vertices as a cost coefficient varies, make the objective quadratic, and that trajectory
becomes the critical line. Frank and Wolfe \cite{frankwolfe1956}, Beale
\cite{beale1959}, and Wolfe \cite{wolfe1959} gave general quadratic-programming
algorithms over the same years, several appearing in the same journal as Markowitz's---
the Critical Line Algorithm contemporary with, and independent of, the general theory
that would later be seen to contain it.

Two threads then ran from this beginning for half a century, both inside finance. On the
modelling side, Sharpe \cite{sharpe1963}, again at Markowitz's suggestion, introduced the
single-index model, writing the covariance as a diagonal matrix plus a rank-one term---
precisely the structure Chapter~\ref{ch:operator} exploits, and the ancestor of every
modern multi-factor risk model. On the algorithmic side the CLA was re-derived, refined,
and scaled across decades, almost always from within portfolio theory
\cite{perold1984,hirschberger2010,niedermayer2010,markowitz2000,stein2008,bailey2013,markowitz2021}.
The lineage is unbroken; what it never asked was whether the algorithm belonged to a
family larger than portfolio selection.

\begin{remark}[Why the method never travelled]
One can only speculate, but the honest reading is that it was easy to do without. A
single efficient point is a quadratic program a generic solver dispatches in
milliseconds, and one approximates the frontier by re-solving as the return tilt is
varied. The parametric machinery that returns the whole frontier exactly bought little
that a loop over a black-box solver did not.

The surviving implementations bear this out: the canonical reference codes support
neither a factor risk model in place of a dense covariance nor inequality constraints
beyond simple bounds---precisely the generality Chapter~\ref{ch:operator} draws out.

The neglect cuts close to home. Markowitz's 1952 paper is among the most-cited articles
in all of economics, while the 1956 paper that actually gave the algorithm is
comparatively forgotten, with no clean citation record of its own, references scattering
across the journal article, its 1955 RAND memorandum, and later reprints. If finance
largely overlooked its own computational contribution, it can hardly fault statistics for
later arriving at the same homotopy without reference to it.
\end{remark}

\section{The 1990s and 2000s: statistics}

Statistics arrived at the same device around forty years later, by a different door, and
at first without reference to any of this.

The door was variable selection and shrinkage, not optimization. The question was not how
to trace a frontier but how the LASSO of Tibshirani \cite{tibshirani1996} chooses, as its
penalty relaxes, which coefficients to admit---and the discovery was that it does so
along a piecewise-linear path. Osborne, Presnell and Turlach \cite{osborne2000} gave the
homotopy explicitly; Efron, Hastie, Johnstone and Tibshirani \cite{efron2004} gave least
angle regression, read the path's corners as the events of a single forward procedure,
and made the piecewise-linear path famous well beyond the LASSO itself. Hastie et al.\
\cite{hastie2004svm} carried the construction to the support vector machine; Tibshirani
and Taylor \cite{tibshirani2011} extended it to the generalized lasso, with the fused
lasso and trend filtering among its special cases, by a dual path of the same kind.

What had been an optimization fact about parametric quadratic programs was being rebuilt,
theorem by theorem, as a native part of regularised estimation.

The statistical development reached even the governing principle. Rosset and Zhu
\cite{rosset2007} characterised, within statistics, exactly when a regularised solution
path is piecewise linear---a quadratic loss against a polyhedral penalty or
constraint---which is the criterion the parametric-programming literature had carried
since the 1950s, now restated in the field's own vocabulary. Tibshirani
\cite{tibshirani2013} settled when the LASSO solution, and so the path, is unique.

None of this leaned on the optimization lineage, and there was little reason it should
have: the statistical versions answered statistical questions---about selection, degrees
of freedom, and prediction---that the portfolio and control literatures had never posed.
That the underlying object was the same went unremarked because no one had occasion to
remark it.

The bridge back was late and partial. Brodie et al.\ \cite{brodie2009} observed that a
gross-exposure-constrained Markowitz problem \emph{is} a LASSO---the special case of
Theorem~\ref{thm:identity} without general constraints---and Gaines, Kim and Zhou
\cite{gaines2018} later traced the constrained LASSO path directly, under arbitrary
linear equalities and inequalities, by an active-set homotopy on the same KKT system.
Even so, the observation travelled as a curiosity rather than as the statement that two
large algorithmic literatures describe one structure.

\section{The 2000s: control, and a partial crossing}

The same device was found once more, in control, and this time the crossing was partly
deliberate.

Through the 2000s the constrained linear-quadratic regulator was being solved
\emph{explicitly}, as a function of the state rather than afresh at each sampling
instant. Bemporad et al.\ \cite{bemporad2002} showed that the optimal control law is
piecewise affine over a polyhedral partition of the state space, and Tøndel, Johansen and
Bemporad \cite{tondel2003} gave the multi-parametric quadratic-programming algorithm
computing it by traversing those regions. This is the common scheme carrying a
\emph{vector} parameter (Remark~\ref{rem:vector}).

Roll \cite{roll2008} then made the link to the scalar case plain, tracing
single-parameter piecewise-linear solution paths in the control setting and explicitly
generalising the LARS/LASSO construction. Control thus arrived at both the scalar
homotopy and its multi-parameter form, and arrived---in Roll's case---already aware of
where statistics had been: the lone instance in this story of a field reaching the device
knowing the others had been there.

\section{The single-point solvers}

The parametric story is the striking one, but the single-point methods of Part~II have
their own chronology and it is worth recording.

The dual active-set method is Goldfarb and Idnani
\cite{goldfarb1982,goldfarb1983}. Its restriction to positive definite Hessians was
lifted by Boland \cite{boland1997}; Forsgren, Gill and Wong \cite{forsgren2016} proved
finite termination for paired primal and dual active-set methods more generally; and
Arnström, Bemporad and Axehill \cite{arnstrom2022daqp} rebuilt the same dual walk on
recursive $LDL\T$ updates for embedded model predictive control---the point at which the
control and optimization threads of this appendix rejoin. Arnström and Axehill
\cite{arnstrom2022certification} additionally compute exact worst-case iteration counts
for parametric QP families.

The complementarity thread is older and separate. Non-negative least squares is Lawson
and Hanson \cite{lawson1974}; principal pivoting and the least-index rule are Murty
\cite{murty1988}; block principal pivoting with the anti-cycling guard of
Theorem~\ref{thm:termination} is Júdice and Pires \cite{judice1994}, with
\cite{cottle1992} the standard reference for linear complementarity. The primal--dual
active set strategy read as semismooth Newton is Hintermüller, Ito and Kunisch
\cite{hintermuller2002}. The Krylov side is Hestenes and Stiefel \cite{hestenes1952},
and the feasible-point family that Remark~\ref{rem:feasible} contrasts runs from
Bertsekas \cite{bertsekas1982} through Moré and Toraldo \cite{moretoraldo1991} to
Dostál \cite{dostal1997,dostalschoberl2005}.

\section{One structure, four disconnected literatures}

Stand back and the shape is the same in each account: a parametric convex quadratic, an
active set, a piecewise-linear path traced by a ratio test.

Optimization has, all along, had the general language for this---parametric programming,
continuation and homotopy methods, active-set sensitivity analysis---within which every
instance is a special case. In that precise sense none of the device is new. What is
striking is not that the structure existed but that the literatures using it stayed
largely disconnected. The portfolio, regression, kernel, and control versions were
developed, taught, and implemented in parallel, each with its own vocabulary
(Appendix~\ref{app:notation}), its own canonical references, and its own software, and
the general framework that subsumes them was seldom the frame in which any one of them
was presented. A practitioner of one could, and routinely did, work a career without
meeting the others.

It is that disconnection, rather than any depth hidden in the mathematics, that kept the
shared structure from view. And it is a reason to be suspicious, in one's own field, of
any method presented as native to it.

\begin{quote}
That an idea is reached more than once, independently, is less an oversight to correct
than a measure of how good it is.
\end{quote}

\chapter{Software and Reproducibility}
\label{app:software}

The methods in this book are implemented, open source, and testable, and the exercises
marked \emph{Implementation} assume the reader will write code. This appendix says what
exists, what each package is for, and how to check an implementation against something
other than itself.

\section{Reference implementations}

\paragraph{Dual active-set QP.} The method of Chapter~\ref{ch:walking}, together with
the guessing fast path of Chapter~\ref{ch:guessing} and the factor reuse of
Chapter~\ref{ch:warmstart}, is implemented in the \texttt{cvx-quadprog} package
\cite{cvxquadprog}. Its signature carries the two conventions this book inherited: the
linear term is subtracted and constraints are held column-wise.

\paragraph{Non-negative conjugate gradients.} The loop of Chapter~\ref{ch:guessing} with
the matrix-free Krylov inner solver of Chapter~\ref{ch:krylov}, including the
equality-augmented variant, the regularising split, and the Nyström preconditioners of
Chapter~\ref{ch:conditioning}, is released as \texttt{nncg}
\cite{schmelzer2026nncg}. Its continuous-integration suite is the paper's synthetic
study, which is a good model to copy: \emph{the test suite is the experiment}.

\paragraph{The Critical Line Algorithm and the path tracer.} The homotopy of
Chapter~\ref{ch:parametric}, the operator contract of Chapter~\ref{ch:operator} with its
four backends, and the LASSO client of Chapter~\ref{ch:homotopy} are implemented in
\texttt{cvxcla} \cite{cvxcla}. It factors into a problem-agnostic path tracer of which
the frontier tracer and the LASSO solver are two clients---which is the software form of
Theorem~\ref{thm:identity}.

\paragraph{General-purpose comparators.} Clarabel \cite{clarabel2024} (interior point)
and OSQP \cite{osqp} (operator splitting) are the solvers Chapter~\ref{ch:measurement}
benchmarks against, and both have native APIs worth calling directly. CVXPY
\cite{diamond2016} is the modelling layer whose overhead that chapter warns about
attributing to an algorithm.

\section{How to validate an implementation}

Four levels, in increasing order of how much they tell you. Every exercise in this book
marked \emph{Implementation} can be checked at one of them.

\paragraph{1. The certificate.} Every returned point must satisfy the KKT
test~\eqref{eq:certificate} \emph{against the original data}, not against the
construction that produced it. This is cheap, it is self-contained, and it catches most
errors. Make it an assertion in your test suite, not a diagnostic you run when
suspicious.

\paragraph{2. An independent solver.} Solve the same instance with a general-purpose QP
solver and compare minimisers. Agreement to solver precision---typically $10^{-8}$ or
better on well-conditioned data---is the expectation. For a parametric method, compare
at each of a grid of parameter values: the traced segment evaluated at a grid point
should reproduce the independently solved QP there.

\paragraph{3. A closed form the code does not use.} Chapter~\ref{ch:walking} gives the
model: the iteration computes $z$ and $r$ from carried factors, and
Proposition~\ref{prop:invariance} gives closed forms in $(G,A,n_*)$ alone. Computing
both and comparing to machine precision tests the factorisation machinery against the
mathematics rather than against a second implementation of itself. Similarly,
Proposition~\ref{prop:increment} lets a running objective be checked against a fresh
evaluation at every step.

\paragraph{4. A different field's reference.} The most convincing check available for
Part~III, and one that only the identity makes possible: trace a LASSO path with a
frontier tracer and compare against an established LARS implementation, breakpoint by
breakpoint. Agreement at the $10^{-15}$ level on a standardised design, including on
high-dimensional problems, is achievable and is a genuinely independent test---different
authors, different field, different code lineage.

\section{Two hazards worth building tests for}

\paragraph{The hazard that never fires on your test family.} Chapter~\ref{ch:measurement}
made the general point; here is the concrete test. Duplicate a column of the constraint
matrix and re-run your guessing fast path. If your solver has no rank test and no
certificate fallback, it will now return plausible wrong answers, and it will do so on
the sort of model any programmatic assembly can produce. Add the test.

\paragraph{The degeneracy that only appears at scale.} The feasible-set projection of
Chapter~\ref{ch:numerics}, Section~6, addresses a failure that appears only after
hundreds of steps on a near-rank-deficient operator. A test suite that runs only small
well-conditioned instances will never see it. Build a controlled rank-deficiency sweep
and assert both branches: that the trace completes and matches a reference while the free
block stays full rank, and that the solver \emph{declines} once it does not.

\section{On reproducing numbers}

Every empirical figure quoted in this book comes from the companion papers, and each is
reproducible from their experiment code---the S\&P~500 and FTSE~100 studies, the
synthetic scaling studies, the constraint-family sweeps behind the $93\%$ figure of
Chapter~\ref{ch:guessing}, and the out-of-sample study of Chapter~\ref{ch:measurement}.

Two habits from those papers are worth copying regardless of subject. First, the figures
quoted in the text are generated by the same run that produces the tables, so prose and
table cannot drift apart. Second, empirical claims made anywhere in a paper are measured
in its own experiments rather than quoted from elsewhere. Both are cheap to adopt and
they eliminate an entire category of error that peer review does not catch.


\backmatter
{\footnotesize
\bibliographystyle{plain}
\bibliography{bib/refs}
}
\printindex

\end{document}
